% =====================================================================
%  R. Nielsen, OM „DEN GODE VILLIE“ SOM MAGT I VIDENSKABEN.
%  Kjøbenhavn. Forlagt af den Gyldendalske Boghandel (F. Hegel.)
%  Trykt hos J. H. Schultz. 1867.   141 pp. + [1] p. Rettelser.
%
%  Diplomatic transcription from the Royal Danish Library scan
%  (`bibliotek/Nielsen, Rasmus/1867-gode-villie.pdf`,
%  https://www.kb.dk/e-mat/dod/110308002961-color.pdf, 153 PDF pages,
%  A4 page boxes, one 300 ppi colour JPEG per page,
%  sha256 37c4b9234fa21f264b71df85cf833e42a1437c37bba4fe35f9e45b5bfef04589).
%  KB copy, shelfmark label "DA 1.-2 S 3 8".
%
%  The exact title, from the title page (PDF 6): the words Den gode Villie
%  stand in Danish low-high quotation marks, „Den gode Villie“, which is
%  how KB's catalogue renders its »Om "Den gode Villie" som Magt i
%  Videnskaben«. The catalog id is `gode-vilje`; the book spells "Villie".
%
%  This header is the edition's apparatus. The transcription harness
%  (pagemap.py, ocr.sh, spacing.py, check.py, splice.py, joints.py) is not
%  tracked in git and is unrunnable without the scan, so everything an
%  outside reader needs in order to check this text against the scan is
%  recorded here.
%
% ---------------------------------------------------------------------
%  1. PAGE MAP  (verified by eye 21 Sep 2026; see pagemap.py)
% ---------------------------------------------------------------------
%
%      printed 1--141   =   PDF page - 7      (PDF 8--148)
%
%  Uniform. The head of every one of PDF 9--148 was rendered and its folio
%  read on the image: 2--141 in unbroken sequence, no leaf missing or
%  doubled, no misprinted folio. p. 1 (PDF 8) carries no head folio (opening
%  page), only the signature "1" at the foot. The embedded text layer fails
%  to read about twenty folios (e.g. PDF 15, 39, 45, 71); the image reads
%  them all as expected.
%
%  Unnumbered leaves (PDF pages):
%      PDF   1      KB digitisation notice -- not part of the book
%      PDF   2      front board (marbled paper)
%      PDF   3      inside board: KB shelfmark label, barcode
%      PDF   4--5   blank endleaves
%      PDF   6      TITLE PAGE (with a KB ownership stamp, not transcribed)
%      PDF   7      verso of title, blank
%      PDF 149      "Rettelser." (errata), = p. [142]
%      PDF 150--152 blank
%      PDF 153      back board
%  There is no Forord, no Indhold, and no separate epigraph leaf.
%
% ---------------------------------------------------------------------
%  2. STRUCTURE  (from the images, not the OCR)
% ---------------------------------------------------------------------
%
%      pp.  1--8    untitled opening. It begins at once, with no head, in a
%                   quotation: Martensen's quotation of Aristotle, Metaph.
%                   XII 10, quoting Iliad 2.204 in Greek type:
%                   „Οὐκ ἀγαθὸν πολυκοιρανίη, εἷς κοίρανος ἔστω.
%                   This is body text of p. 1, not an epigraph.
%      p.   9       I.   Naturlærens Skjæbne i Monarchiet.
%      p.  40       II.  Historiens Vilkaar i Monarchiet.
%      p.  75       III. Philosophiens Stilling i Monarchiet.   (to p. 141)
%
%  The text itself calls these "Afdeling" (første / tredie Afdeling). The
%  heads are centred: a roman numeral on one line, the title below in
%  larger Fraktur. Set as \section* with the numeral, \\, and the title.
%  There are NO § heads: every § in the book is a citation (Martensen's
%  Dogmatik § 132, Grundideernes Logik § 13 ...). Inside the parts the
%  argument is divided by run-in, letterspaced numbered propositions --
%  "1) at en Troesvidenskab (Theologie) er utænkelig." (p. 34), "2) ..."
%  (p. 35), "3) ..." (p. 35); "1) Det Phlyariske i sit negative Moment."
%  (p. 91), 2) p. 92, 3) p. 93; "1) Theismen i Grundideernes Logik er kun
%  en omskreven Pantheisme." (p. 118), 2) p. 131 -- which stay run-in.
%
% ---------------------------------------------------------------------
%  3. TYPOGRAPHY AND TRANSCRIPTION CONVENTIONS
% ---------------------------------------------------------------------
%
%  Body in Fraktur. aa (not å); ø is the slashed Fraktur ø (the embedded
%  layer misreads it as "ö" -- "religiös" -- but the print has ø: checked
%  at p. 25 "religiøse"). German quotations are in Fraktur (p. 46 note).
%  Emphasis is by LETTERSPACING -> \emph{}. Latin set in antiqua -> \textit{}
%  (p. 46 note "Omnia et singula ...", p. 91 "credo, ut intelligam",
%  "Metaphys. XII, 10" on p. 1 has the numeral in antiqua). Spaced antiqua
%  -> \emph{\textit{}}. Fat-face (halbfett) Fraktur: none in the book. Greek in Greek type (p. 1; p. 46 note Θέσις),
%  typed directly (textalpha). Quotation marks „...“, low-high; the closing
%  sort is a straight double prime ″ in this fount, transcribed “.
%  Footnotes marked *), **) per printed page (p. 46 and p. 71 carry two):
%  \thefootnote gives *), **), ***), reset by \setcounter{footnote}{0}%
%  after every page marker. Fraktur double hyphen "=" at line end is a
%  word-break hyphen, resolved. Fraktur I/J: one sort; house convention I
%  for I-words (Idee, Incarnation), J for J-words (Jeg). Signature marks at
%  page feet ("1", "3*" ...) are not transcribed.
%  Plain o for ø (as in Logik og Psychologie, same printer, same year) does
%  NOT occur: every ø checked on the image, in all ten batches, has its
%  slash. Printed line breaks are kept in the source except in pp. 54--66,
%  where the lines are reflowed; the words are the same either way.
%
% ---------------------------------------------------------------------
%  4. ERRATA  (the book's own "Rettelser", PDF 149; NOT applied)
% ---------------------------------------------------------------------
%
%  Transcribed as printed at the end of this file. The text is left as
%  printed at each site and a % comment there points to the Rettelser:
%      p. 17 l. 4 fra oven          legitim;   læs: legitim“;
%      p. 43 l. 17 fra neden        saadan     læs: saadant
%      p. 46 Anm. l. 7 fra neden    Sanctus    læs: Sancti
%            (the KB copy carries a hand correction in ink at this word)
%      p. 61 l. 10 fra oven         hedder det om  læs: hedder det, som
%                                   ovenfor anført, om
%      p. 61 Anm. Lin. 1 fra neden  S. 3       læs: S. 7.
%
% ---------------------------------------------------------------------
%  5. WORD-ACCURACY
% ---------------------------------------------------------------------
%
%  Word-accuracy: every batch diffed against the scan's embedded layer
%  before splicing (ocrdiff.py). 10 batches, 909 candidates, 0 transcription
%  errors corrected, 0 unresolved (every candidate was the layer's own
%  Fraktur misreading: sk->fl, ae->cr, ø->o, Martensen->Mortensen; one
%  error, "vaier" for "veier" on p. 1-13, was caught by its batch agent
%  before the diff ran). Two further checks, 21 Sep 2026:
%   (a) a whole-book word alignment against the same layer (87% of words
%       matched after normalising the layer's systematic confusions); all
%       79 short divergences were examined and the doubtful ones read at
%       the image: 0 transcription errors, all layer displacement or noise;
%   (b) an independent collation at the image of six randomly drawn pages
%       (pp. 2, 3, 61, 63, 106, 141; ~1,600 words, every word, mark and
%       emphasis extent): 0 discrepancies.
%  See TRANSCRIPTION-PLAYBOOK.md §3 step 4 and §7.
%
% ---------------------------------------------------------------------
%  6. PRINTER'S DEFECTS LOG
% ---------------------------------------------------------------------
%
%  All transcribed as printed and logged by a % comment at the site.
%      p.   6      οὐσίας: no accent visible over the ι at 300 ppi (doubtful)
%      p.   7      „videnskabelig“ with a broken final g
%      pp. 16, 19  ἀσθηνῆ / ἀσθηνη καὶ πτωχ(ὰ/α) στοιχεια: η for ε, accents
%                  wanting; p. 19 κτισις, φυσις unaccented
%      p.  25 n.   MDCCVLIV for MDCCXLIV
%      p.  47      „ before „Historien om Syndefaldet“ never closed
%      p.  54      stray full stop in „med. den“
%      p.  56      stray mark over the final e of „omtvistede“
%      p.  63      „Menen“ for „Mening“
%      p.  64      χριστᾷ for χριστῷ; a surplus closing “
%      p.  71 n.   „N. H. Clausen“, initials reversed (p. 66: H. N.)
%      p.  72      stray full stop in „Theo-|log.“
%      pp. 87--88  a quotation resumed without a reopening „
%      p.  94      full stop (or tailless comma) after „folkelige“
%      p. 121      „hverkeu“, turned n
%      p. 131      outer quotation never closed
%      p. 136      a connective wanting („… Nødvendighedsbestemmelser
%                  Frihedsbestemmelser …“)
%  QUOTE BALANCE: „ 519, “ 520, balance -1, and it must stay -1. It is the
%  sum of: +1 p. 17 (the missing “ the Rettelser supply), +1 p. 47, -1 p. 64,
%  -2 pp. 87--88, +1 p. 131, -1 the Rettelser entry for p. 17.
%  KB-copy hand marks (ink corrections at pp. 43, 46, 61, 92; pencil in the
%  margins of pp. 42, 58, 106) are noted at the sites, never transcribed.
%
% =====================================================================
\documentclass[11pt]{book}
\usepackage[utf8]{inputenc}
\usepackage[T1]{fontenc}
\usepackage{libertinus}
\usepackage{libertinust1math}
\usepackage{textalpha}
\usepackage{graphicx}
\DeclareUnicodeCharacter{0254}{\reflectbox{c}}
\usepackage{amsmath}
\usepackage[danish]{babel}
\usepackage{fancyhdr}
\setlength{\headheight}{28pt}
\addtolength{\topmargin}{-13.5pt}
\pagestyle{fancy}
\fancyhf{}
\fancyhead[LE,RO]{\thepage}
\fancyhead[RE]{\textit{Om „Den gode Villie“ som Magt i Videnskaben}}
\fancyhead[LO]{\textit{\leftmark}}
\renewcommand{\thefootnote}{\ifcase\value{footnote}\or *)\or **)\or ***)\else
  \arabic{footnote})\fi}
\usepackage{setspace}
\setstretch{1.3}
\usepackage{marginnote}
\newcommand{\opage}[1]{\marginnote{\footnotesize\textit{[#1]}}}
\newcommand{\apage}[1]{\marginnote{\footnotesize\textit{[#1]}}}

\begin{document}

\frontmatter
\begin{titlepage}
\centering
\vspace*{3cm}
{\LARGE Om}\\[1.5em]
{\Huge „Den gode Villie“}\\[1.5em]
{\LARGE som Magt i Videnskaben.}\\[2em]
{\small Af}\\[1em]
{\large R. Nielsen.}\\[5em]
Kjøbenhavn.\\[0.5em]
Forlagt af den Gyldendalske Boghandel (F. Hegel.)\\[0.3em]
{\footnotesize Trykt hos J. H. Schultz.}\\[0.5em]
1867.
\end{titlepage}

\mainmatter
\markboth{Om „Den gode Villie“}{Om „Den gode Villie“}

% Heads, set by the batch agent at the page where they are printed:
%   p.  9   \section*{I.\\ Naturlærens Skjæbne i Monarchiet.}
%   p. 40   \section*{II.\\ Historiens Vilkaar i Monarchiet.}
%   p. 75   \section*{III.\\ Philosophiens Stilling i Monarchiet.}
% each followed by \markboth{<title without stop>}{<same>}.

% --- p. 1 ---
\opage{1}\setcounter{footnote}{0}%
„Οὐκ ἀγαθὸν πολυκοιρανίη, εἷς κοίρανος ἔστω.
„\emph{Mange til Herskab er aldrig til Gavn, kun Een
være Hersker}“. Ved at give dette Homeriske Ord en
videnskabelig Anvendelse har Aristoteles villet sige, at mange
Principers Herredømme er ikke til Gavn, men til Fordær\-%
velse, ikke blot for Stater, men ogsaa for videnskabelige
Systemer, hvor de medføre Anarchie og indre Opløsning,
men at eet Princip bør være det Herskende og Altstyrende,
ligesom en Hær kun bør have een øverste Feltherre, da
ellers Alt kommer i Uorden. Han tilføier: „det Værende
vil ikke regjeres slet.“ (Metaphys. XII, 10).“ Med disse
Ord har Biskop Martensen indledet, hvad han kalder „Slut\-%
ningsbetragtninger“ til sit Leilighedsskrift: \emph{Om Tro og
Viden}; det er i Sandhed mærkelige Ord, Ord, der vel
fortjene, at vi alvorlig lægge os dem paa Sinde.

Hvad vilde Theologiens Antagonister, den negative
Kritiks Ordførere, hvad vilde en D. Strauß, en Bruno
Bauer, en L. Feuerbach vel tænke og sige, naar de læste
disse indledende Ord i et Skrift, hvis Forfatter de endnu
ikke kjendte? De vilde vel sige: „See, det er da en frisk
og resolut Begyndelse; saaledes bør en fornuftig Mand gribe
Tingen fra Grunden af: ingen Tvetydigheder, ingen Omsvøb!
Mange Principers Herredømme er ikke til Gavn, men til For\-%
dærvelse for videnskabelige Systemer; een være Hersker! Det
er Fornuften selv, den kongelige Fornuft, den selvskrevne Ene\-%
hersker, for hvis ubestridelige Ret vi saalænge have stridt og
% --- p. 2 ---
\opage{2}\setcounter{footnote}{0}%
arbeidet. Først naar Fornuftprincipet aabenlyst er anerkjendt
for eneherskende Princip i Aandens Rige, kan man faae
Bugt med det halvvidenskabelige Uvæsen, der kalder sig
Theologie, og omsider vente Befrielse for al den uklare
Gemytlighed, den skimlede, halvreflecterede Umiddelbarhed,
hvormed speculative Seminarister og tankeløse Katechismus\-%
præster kun altfor længe have opvartet os. Mange Prin\-%
cipers Herredømme er ikke til Gavn, kun Een være Hersker:
freidigt, ja med Ungdomsfriskhed har Forfatteren taget
Ordet for Kritik og Videnskab, for Fornuftens Autonomie,
den videnskabelige Selvbevidstheds absolute Eenhed!“ Ja
saa træffende, saa vækkende ere Biskop Martensens Ord, at
Fornuft-Autonomiens allerivrigste Forsvarere ligefrem kunne
tage dem til Indtægt.

Men, det forstaaer sig, naar Autonomiens Forkæmpere
gik videre i Texten og til deres store Overraskelse erfarede,
hvad her egentlig sigtes til: vilde de jo vistnok slaae om i en
anden Tone. „Saa det er ikke Fornuftprincipet, men Troes\-%
principet, altsaa det gode gamle Mirakelprincip, der med nogle
moderne Omsvøb tilbyder sig at holde Styr paa Viden\-%
skaberne. Een være Hersker: Troen i Theologien, Theo\-%
logien i Troen! Et videnskabeligt Troesmonarchie, et theo\-%
logisk absolut Universalmonarchie: hvilket Uhyre af en Tanke,
hvilket Tanke-Uhyre! Men primer Theologien, gaaer den
i Barndom? Den nøies altsaa ikke med at være som
hidtil en forgjemt, forklemt, forhutlet Videnskab for sig selv,
en Videnskab, som, naar den vrede Kritik opfordrede den
til at drage Klingen, bad saa bevægeligt: Lad mig være, til
slig en Ære er jeg for ydmyg, --- nei, Fortidens Minder
spøge i dens gamle Hjerne; den nærer ærgjerrige Planer,
den vil tilveirs, den vil som Videnskaben med det over\-%
naturlige Princip, altsaa dog som Mirakelvidenskab --- \textit{risum
teneatis} --- være, hvad det Absolutes Philosophie trods sin
dybsindige Speculation ikke kunde være, den centrale Viden\-%
% --- p. 3 ---
\opage{3}\setcounter{footnote}{0}%
skabernes Videnskab. Hvad er her da skeet? Ere Conjunc\-%
turerne overordentlige, Constellationerne lykkevarslende; traadte
Solen maaskee ind i Løvens Tegn, just som Mars var i
Conjunction med Venus? Eller er det Klogskabshensyn,
fiin Beregning, der nu seer Qvotienten af de gunstige og
de mulige Betingelser nærme sig Eenheden? Det Over\-%
ordentlige maa forudsættes, dersom det skal lykkes med Uni\-%
versalmonarchiet; skulle de fribaarne Videnskaber virkelig
underkaste sig en Mirakelvidenskab, da maa der skee et
videnskabeligt Mirakel.“

Vi have tænkt os, hvad Theologiens Plage-Aander,
„Fritænkerne“, „Vantroens Børn, der have forspiist sig i
Kundskabens Frugt“, vilde sige, naar de betragtede Sagen
fra Fornuftens Standpunkt. Men dersom de nu ikke ind\-%
skrænkede sig til udelukkende at see paa den Fornuften tilføiede
Uret, men tillige gave sig til at betragte Theologiens Vil\-%
kaar, hvordan mon de da vilde blive tilmode? De vilde
blive ilde tilmode; thi det vilde ærgre dem, om Theologien,
den af dem saa ofte haanede og bespottede Theologie, virkelig
skulde tage et saadant Opsving, sætte sig paa Thronen, gribe
Troens Scepter og sige homerisk: Een være Hersker! Ja,
det vilde ærgre dem, og Ærgrelsen vilde forøges, naar de i
Phantasien ikke kunde Andet end udmale sig al den Ære
og Glæde, en saadan Opkomst maatte forskaffe Theologerne.
„Hvilken Sang, hvilken Klang! Ringer det ikke for vore
Øren, som hørte vi Seiershymner fra den længe beleirede
Fæstning. See Dogmatikens Banner veier høit fra alle
Taarne; Hermeneutiken illuminerer; Exegetiken stikker sin
Fane ud; selv Indledningen til det Nye Testamente stikker
en Klud ud --- den har kun en Klud. Ja, der er Røre i
Theologernes Leir: det gamle, nær havde vi sagt hovedløse,
Facultet opløfter Hovedet, skygger for Øinene, skimter det
nye Lys, lytter til det glade Budskab, og det glade Budskab er
Universalmonarchiets Evangelium.“

% --- p. 4 ---
\opage{4}\setcounter{footnote}{0}%
Slige Stemninger og Udtalelser fra et Standpunkt,
der ikke blot vil have Theologien, men ogsaa Religionen op\-%
hævet, antydes ingenlunde for at kaste Skygge paa det af
Biskop Martensen opstillede Princip; tvertimod, de antydes
netop for at sætte Principet i et klarere Lys, forsaavidt de
gjøre det kjendeligt, at her i Theologiens Navn er udtalt en
stor og resolut Tanke. Vi ere ganske rigtigt komne til et
sidste Dilemma: enten maa Theologien bort, eller den maa
indtage sin oprindelige Plads som høieste Videnskab. I
Charakteer af høieste Videnskab maa Theologien da bestemmes
som herskende Videnskab, som den Videnskab, der --- natur\-%
ligviis ikke vilkaarlig, men med en velbegrundet Ret, ikke
umiddelbart, men middelbart, ikke despotisk, men organisk ---
hævder Primatet over samtlige Videnskaber, efterdi det Lys,
hvori den seer Tilværelsen, er det høieste og reneste Lys, er
Aabenbaringens hellige Lys, det Sandhedens Lys, hvori alle
endelige Erkjendelser maae finde deres sidste og høieste
Forklaring. Tanken er, som man seer, i sig selv baade stor
og skjøn.

I den absolut guddommelige Idee er den levende Sand\-%
hed Eet med den hellige Kjærlighed, og denne Kjærlighed
ikke blot Sandheden i sig selv, men Kilden til al Sandhed.
Der kan i den guddommelige Viden umulig være to absolut
ueensartede Sandheder, en upersonlig og en personlig. Jo
mere den menneskelige Viden antager Lighed med den gud\-%
dommelige, jo mere vil Stykværket hæves, vil det Uperson\-%
lige underordne sig det Personlige, og Skilsmissen mellem
de ueensartede Sandheder forsvinde. Men den menneskelige
Viden kan kun blive en sand Afglands af den guddommelige
paa den Betingelse, at den i høieste Forstand udgaaer
fra Troen, er og vil være en Viden i Troen, af Troen.

Den høieste, reneste Viden i Troen er en speculativ
Viden. „Speculation“, siger Schelling, „er Alt, er i Gud
at skue, at betragte det, som er. Videnskaben selv har kun
% --- p. 5 ---
\opage{5}\setcounter{footnote}{0}%
forsaavidt Værd, som den er speculativ, det er Contemplation
af Gud, som han er. Men den Tid vil komme, da Viden\-%
skaberne mere og mere ville ophøre, og den umiddelbare
Erkjendelse indtræde. Kun i den høieste Videnskab lukker
det dødelige Øie sig, idet ikke mere Mennesket seer, men den
evige Seen selv er seende i Mennesket.“ Det høie Maal,
hvortil herved sigtes, kan umuligt naaes i abstract Specu\-%
lation; det kan kun --- ved Tilnærmelse --- naaes i den
Speculation, hvis Udgangspunkt er den levende Gud, altsaa
kun i den fyldige, indholdsrige Speculation, der begynder
og ender med Troen.

Saaledes opfattet faaer det af Biskop Martensen op\-%
stillede Princip da ogsaa en gjennemgribende Betydning for
Udviklingen af Videnskabernes Encyclopædie. „Encyclopæ\-%
dien,“ siger Hegel, „er egentlig en Fremstilling af Philoso\-%
phiens almindelige Indhold; thi hvad der i Videnskaberne
er grundet paa Fornuft afhænger af Philosophien, hvad der
derimod er i dem af det, der beroer paa vilkaarlige og udvortes
Bestemmelser, eller, som det hedder, er positivt og statutarisk,
ligger udenfor den.“ Hermed er det udtalt, at samtlige Vi\-%
denskaber udgjøre eet eneste Videnskabernes System. Den
systematiske Videnskabernes Eenhed betinges af, at der med
Ideens Ret gjøres Forskjel imellem lavere og høiere, under- og
overordnede Videnskaber, og beroer i sidste Instans paa,
at der gives en absolut Videnskabernes Videnskab. Havde
den abstract logiske Viden været istand til at magte det Ab\-%
solute, da vilde det ganske vist have været Philosophien, der
var bleven Centralvidenskab. Men da det, som ovenfor an\-%
tydet, kun er den troende, følgelig ikke den autonomiske, men
den theonomiske Speculation, den Speculation, der begynder og
ender med det ethiske Gudsbegreb, som kan give den menne\-%
skelige Viden sit høieste Indhold: saa bliver det ikke Philo\-%
sophien, men Theologien, vi maae erkjende for den egentlige
Centralvidenskab.

% --- p. 6 ---
\opage{6}\setcounter{footnote}{0}%
En Ideelære, der sigter paa Christendommen, „maa“,
siger Biskop Martensen, (Leilighedsskr. S. 122) „først og frem\-%
mest tage Sigte paa det Godes Idee, som den, der er ἐπέ\-%
κεινα τῆς οὐσίας, og sørge for at der kommer en grundig % οὐσίας: no acute is visible over the iota on the scan (p. 6 l. 4); absence not taken as a printer's defect
Subordination ind i Ideeriget.“ Hermed stemmer, hvad
Biskoppen udtaler i en tidligere Sammenhæng (S. 102):
„Naar Religions- eller Aabenbarings-\emph{Philosophien} --- thi
en Religionsphilosophie, der ikke er Aabenbaringsphilosophie,
er kun en Halvhed --- er den \emph{universelle} Videnskab, for\-%
saavidt som den erkjender Aabenbaringen i dens Forhold til
Mythologien, til Naturens og Historiens hele Mangfoldig\-%
hed: saa er Theologien den \emph{centrale} Videnskab, idet den
udelukkende fordyber sig i det Ene, i Guds Rige som
saadant.“

Naar vi med disse Oplysninger ret betænke, hvilken
Vægt Biskop Martensen lægger paa sit homerisk-aristoteliske
Ord: Een være Hersker! da kunne vi ikke noget Øieblik
tvivle om, at han bogstavelig og for ramme Alvor vil fordre
en, om end middelbar, Subordination af alle relative Viden\-%
skaber under sin centrale Videnskab, at han med andre Ord,
vil oprette et nyt theologisk Universalmonarchie.

Men et nyt theologisk Universalmonarchie, hvilken uhyre
Tanke! Mon en Tanke af denne Høide kan, uden at bukke
sig saa dybt, at den bliver ganske kroget, gaae ind ad Virke\-%
lighedens Dør? Ogsaa jeg har mødt den høie Tanke; den
er mange Gange traadt mig imøde, naar jeg betragtede
Dogmatiken, ikke fra Skyggesiden, men fra Lyssiden; thi,
ihvorvel jeg, ifølge min Opgave, som oftest har maattet
gjennemvandre dens Skyggepartier, har jeg derfor ikke glemt,
hvor jeg endnu kan finde de solbeskinnede Steder. Men naar
jeg saa, efter at have beskuet Idealet i al dets Herlighed,
alvorlig spurgte mig selv: kan et saadant Ideal realiseres?
har jeg bestandig faaet et benegtende Svar.

Det er jo vistnok saa, at det Godes Idee, opfattet som
% --- p. 7 ---
\opage{7}\setcounter{footnote}{0}%
Kjærlighedens, den hellige Villies Idee er i guddommelig Viden
Eet med den evige Sandhedsidee; men heraf følger jo ingenlunde,
at Sandhedsideen i en menneskelig Viden er Eet med det Godes
Idee i menneskelig existentiel Villen. Det er jo vistnok
saa, at der i Viisdommen selv ikke kan tænkes nogen Kløft
imellem Troen og Videnskaben: men hvorfor er Gud ingen
Videnskabsmand? Fordi Gud ikke er nogen Troende. At
tænke sig den guddommelige Viden som en videnskabelig, være % videnskabelig: the final g is a broken sort, only a fragment inked
sig philosophisk eller theologisk Viden, er ligesaa gudsbespotte\-%
ligt, som at tale om en troende Gud. Hvorfor skal et
Menneske opfatte det Absolute i Troen? Fordi den reneste
og høieste Udvikling af den menneskelige Viden er, ikke gud\-%
dommelig, men videnskabelig. Hvorfor skal et Menneske op\-%
fatte det Absolute i Viden? Fordi Troen er en menneskelig,
ikke en guddommelig, Erkjendelse. Hvorfor viser det Abso\-%
lute, der opfattes i menneskelig Viden, sig saa ueensartet med det
Absolute, der opfattes i Troen? Fordi de relative Opfattelses\-%
former gribe det Absolute, ikke i dets centrale Dybde, saa\-%
ledes som det af os \textit{in abstracto} maa tænkes at være med
uendelig Concretion i det guddommelige Selv, men i to
polære Extremer, af hvilke det ene er betegnet ved den
almene, i rationel Nødvendighed grundende Tænkning, det
andet ved det punktuelle, i Samvittigheden spirende Villies\-%
liv. Her er i det Absolute Noget, jeg absolut skal troe,
fordi der i det Absolute er Noget, jeg absolut skal vide.

Hvad dernæst Forholdet imellem Videnskaberne indbyrdes
angaaer, er det virkelig, encyclopædisk seet, en egen Sag
med den strenge Subordination. Hver Videnskab er paa sit
Omraade autonomisk, er inden sine Landemærker en Sou\-%
verain. Vistnok kan den ene Souverain subordinere sig
den anden, men bestandig med et Forbehold, der sikkrer
Selvstændigheden. Skulde Fagvidenskaberne samle sig i Sub\-%
ordination under en høieste, eneherskende Centralvidenskab, da
vilde de vist begynde deres Hilsen med de bekjendte Ord:
% --- p. 8 ---
\opage{8}\setcounter{footnote}{0}%
Vi, der hver for sig ere ligesaa gode som Du og tilsammen bedre
end Du, erkjende o. s. v. Jeg forstaaer jo vel, at Theologien
ei vil lægge an paa at beherske de øvrige Videnskaber umid\-%
delbart, at den f. Ex. ikke vil tvinge enten Astronomien eller
Geologien til, om jeg saa maa sige, at snakke Bibelen efter
Munden; jeg forstaaer, at den kun vil drage de reale,
underordnede Videnskaber ind under Lyset af den hellige
Kjærligheds, den aabenbarede Viisdoms Idee. Men jeg
forstaaer ogsaa, at dersom en saadan theologisk Belysning
virkelig skulde have videnskabelig Betydning, da maatte de
subordinerede Videnskaber kunne bringes til at indsee, at der
ved Hjælp af overnaturlige Forudsætninger opnaaes en over\-%
ordnet Viden, en Viden, der er istand til at anbringe
ideale Correctiver og videnskabeligt forklare, hvad der paa de
reale Forskningers, de exacte Analysers naturlige Stand\-%
punkt viser sig urimeligt, selvmodsigende, ubegribeligt. Men
en saadan Indrømmelse kunne og ville de forskende Viden\-%
skaber vist ikke gjøre.

Da nu Biskop Martensen i sit homerisk-aristoteliske
Ord: „Een være Hersker!“ har givet os Tilsagn om at
ville drage Omsorg for, „at der kommer en grundig Sub\-%
ordination ind i Ideeriget“, og paa denne Maade oprette
et theologisk Universalmonarchie: saa have vi, med Forbehold
af vore Betænkeligheder og vor egen Synsmaade, i det
Følgende at see til, paa hvilke Betingelser han da vil være
istand til at udføre sin Plan, hvorledes han egentlig vil
bære sig ad med at oprette Monarchiet. Vi skulle med
stadigt Henblik paa Forholdet mellem Viden og Villie, theo\-%
retisk og praktisk Erkjendelse, aabne Forhandlinger med de til
Subordination bestemte Videnskaber, saaledes, at vi særlig
henlede Opmærksomheden paa: \emph{Naturlærens Skjæbne},
\emph{Historiens Vilkaar} og \emph{Philosophiens Stilling i
Monarchiet}.

% --- p. 9 ---
\opage{9}\setcounter{footnote}{0}%
\section*{I.\\ Naturlærens Skjæbne i Monarchiet.}
\addcontentsline{toc}{section}{I. Naturlærens Skjæbne i Monarchiet}
\markboth{Naturlærens Skjæbne i Monarchiet}{Naturlærens Skjæbne i Monarchiet}

Der er noget Haardt i det Ord Skjæbne. Idet vi
bruge et saa haardt Ord, maae vi strax afværge en Misfor\-%
staaelse, som om Meningen var, at Theologien, naar den nu
igjen kommer i Besiddelse af Enevoldsmagten, udtrykkelig
havde til Hensigt at behandle Naturlæren haardt. Langtfra!
Hvor skulde en Videnskab, hvis Princip er Kjærligheden
selv og „den gode Villie“, kunne behandle Nogen af sine
Undergivne haardt, at sige, naar den Undergivne frivil\-%
ligt subordinerer sig: Theologien vil ikke være en blind
Skjæbne, den vil tvertimod være et naadigt Forsyn. Men,
det forstaaer sig, de Undergivne maae paaskjønne Naaden;
thi ellers kan det naadige Forsyn let forvandle sig til
Skjæbne.

Den første store Grundsætning, der i Monarchiet maa
indprentes Naturlæren, er den, at efterdi det Praktiske er
theoretisk, og det Theoretiske praktisk, saa maa Theorien
væsentlig rette sig efter Villien: til en god Theorie svarer en
god Villie, til en ond Theorie en ond Villie. Uden at tilegne
sig denne Grundsætning vil Naturlæren ikke engang være
istand til at gjøre Forskjel mellem Sandhed og Løgn. Thi,
siger Biskop Martensen (Leilighedsskr. S. 122): „Modsæt\-%
ningen mellem Sandhed og Løgn er paa eengang practisk og
theoretisk, idet den som Bekræftelse eller Benægtelse af det
Gode som den fuldkomne \emph{Virkelighed} grunder i den gode
eller onde eller dog syndige Villie. Den gode Villie har
den til det Gode svarende Theorie. Men den onde Villie,
der er udtraadt af Kjærlighedens Afhængighed, maa optænke
sig selv en Theorie, tantalisk bestræbe sig for at sætte en ind\-%
bildt Virkelighed istedetfor den objective“. Kan Naturlæren
nu ikke gaae ind paa at gjøre sine Theorier enten middel\-%
bart eller umiddelbart afhængige af Theologiens „gode Villie“,
% --- p. 10 ---
\opage{10}\setcounter{footnote}{0}%
hvad følger saa heraf? Saa følger heraf, at om end Na\-%
turforskninger kunne tillades, naar de opfattes som under\-%
ordnede Forsøg paa at finde sig tilrette med det sandseligt
Virkelige, det blot Logisk-Physiske, saa blive dog Naturviden\-%
skaberne selv at ansee for uægte, i deres inderste Væsen usande
Videnskaber.

Den anden store Grundsætning, der i Monarchiet maa
indprentes Naturlæren, er den, at Culturbevidstheden med
de fra Kirkelæren emanciperede Realvidenskaber hviler paa et
falsk, forbryderisk Grundlag. Det nytter ikke at ville i
Praxis erkjende Friheden, i Theorien Nødvendigheden, i
Praxis æde af Livets Træ, i Theorien af Kundskabens Træ,
i Praxis hævde „den ethiske Villie“, i Theorien den objectiv
gyldige Lov. Det er ikke Virkelighedens, ikke Kjendsgjer\-%
ningernes, ikke Lovenes Skyld, at Naturlæren er udtraadt af
Kjærlighedens Afhængighed; nei, det er Villiens Skyld, ene
og alene Villiens Skyld. Thi, siger Biskop Martensen
(Leilighedsskr. S. 123--124): „Villien vil have det hele, det
altomfattende Primat i Bevidstheden, og dersom den ikke
vilde dette, da var den ikke Villie. Som ethisk Villie vil
den udvælge eller forkaste det Theoretiske, der fremstiller sig
for Bevidstheden, forbeholde sig det afgjørende Ja og Nei,
Placet og Veto, efterdi den selv skal være ansvarlig og af\-%
lægge Regnskab, ogsaa for hele den indre Huusholdning“.
Fremdeles (S. 125--26): „At æde af Kundskabens Træ beroer
paa en Villie og er en Handling, hvis vigtigste Conseqvenser
ere Personlighedsconseqvenser. Og Prometheus, der har lært
Menneskene Videnskab og hver en Kunst og fremkaldt Cul\-%
turbevidsthed paa Jorden, maa, lænket til Klippen, bøde for
den \emph{Handling}, ved hvilken han brød igjennem til sin
Kundskab, fordi han forskaffede sig den ad en Vei, der ikke
var legitim, og fordi han kun gjorde Menneskene kloge, men
ikke gode --- heri et Forbillede paa Nutidens Culturbevidst\-%
hed med dens mange baade aabenbare og hemmelige Lidelser“.
% --- p. 11 ---
\opage{11}\setcounter{footnote}{0}%
Dersom Naturlæren ikke kan faae denne Grundsætning i sit
Hoved, hvad følger saa heraf? Saa følger naturligviis heraf,
at ingen sand Christen kan være Naturforsker, og ingen Na\-%
turforsker en sand Christen.

Men ere Naturvidenskaberne da ikke virkelige og sande
Videnskaber? Man kunde fristes til at antage et saadant
Spørgsmaal for en utidig Spøg, en literær Uanstændighed,
hvis man ikke udtrykkelig tog Hensyn paa ovenstaaende Ud\-%
talelser. Thi dersom vi ellers tør kalde en Lære, der giver
os en systematisk sammenhængende Viden om noget Virkeligt,
en virkelig Videnskab, da er Naturlæren med alle dens Af\-%
delinger ei alene en virkelig, men man kunde fristes til at
sige, den virkeligste af alle virkelige Videnskaber. Kan nu
nogen Videnskab tænkes at være virkelig uden at være sand?
Man skulde ikke troe det. Skille vi Virkeligheden fra Sand\-%
heden, saa bliver Virkeligheden jo usand, og Sandheden
uvirkelig. Er nu en uvirkelig Sandhed usand, og en usand
Virkelighed uvirkelig, saa følger nødvendig heraf, at Sandhed
og Virkelighed maae paa alle virkelige Erkjendelsers Trin
være forenede; og saa følger igjen heraf, at Naturlæren, saa\-%
sandt den er en virkelig Videnskab, maa være en sand Viden\-%
skab. Skal nu Naturlæren fremdeles bevare sin Charakteer
som virkelig og sand Videnskab, da kan det kun skee derved,
at den træder i aabenbar Modsigelse med den af Biskop
Martensen udtalte Fordring: Een være Hersker, nemlig den
gode Villie.

Paa Naturlærens Standpunkt er det saa langt fra, at
vore videnskabelige Sandhedserkjendelser tør rette sig efter
vor gode Villie, at den redelige Forsker tvertimod føler sig
ethisk forpligtet til at lade Villien bøie sig for de strenge
Love, respectere Kjendsgjerninger og sund Logik, og uden
Omsvøb holde sig til det Logisk-Physiske. „Det er ikke sub\-%
jective Følelser, ikke fromme Ønsker, men Indsigt i Tin\-%
genes Væsen, der afleder Lynstraalen og sætter Dampvognen
% --- p. 12 ---
\opage{12}\setcounter{footnote}{0}%
i Bevægelse. Vistnok er Videnskab umulig uden Subjec\-%
tivitet; men Subjectiviteten er tjenende. Her sandses, her
føles, her tænkes; men det i Sandsning og Tænkning blot
Subjective maa overalt udrenses som det Tilfældige, det
Forstyrrende, det Usande: hver Forudsætning er objectiv“.
(Phil. Propæd. S. 74--75).

Hvordan vil det nu under disse „objective Forudsæt\-%
ninger“ i Virkeligheden tage sig ud, naar Biskop Martensen
med sin homerisk-aristotelisk-dogmatiske Grundtanke optræder
midt i Naturforskernes Kreds, og med den Fordring: Een
være Hersker! gjør den gode Villies Ja og Nei, dens Placet
og Veto, til Høiesteret for de forskende Videnskaber? Mon
de Naturkyndige ved en saadan Overraskelse ikke ville tabe
Mund og Mæle, naar den alvorlige Mand, efter at
have rettet en saadan Opfordring til Forsamlingen, tilføier
en Forklaring som denne: „Først naar man fra alle sub\-%
ordinerede Ideer er stegen op til det Gode som til det kon\-%
gelige Princip, kan der blive Tale om den levende, positive
Sandhed som dets Aabenbarelse i Modsætning til en blot
negativ, abstract-logisk Sandhed, selv om denne bestemmes
som Eenhed af Viden og Magt. Og først naar det Gode
er bestemt som den absolute Kjærlighed, naar Modsætningen
mellem Godt og Ondt er erkjendt som den første Grund\-%
modsætning, kan der blive Tale om den levende Modsætning
mellem Sandhed og Løgn, der er noget heelt andet end den
blot logiske Modsætning mellem Rigtigt og Urigtigt, Væsen
og Skin“? (Leilighedsskr. S. 121).

Hvad Naturlæren i de enkelte Grene faaer ud af sine
enkelte Forskninger, hvorledes Astronomen f. Ex. bærer sig
ad med Perturbationsproblemerne, hvorledes Geognosten be\-%
stemmer sine Mineraler og Bjergmasser, derom vil Theolo\-%
gien ikke spørge. Theologien er, hvad Kjendsgjerninger og
deres Analyse angaaer, ikke nøieregnende, ikke smaalig; den
plager ikke sine Undergivne med Forhandlinger om det Spe\-%
% --- p. 13 ---
\opage{13}\setcounter{footnote}{0}%
cielle; den seer paa det Totale. Det er altsaa ikke for de
astronomiske Enkeltheder, men for den astronomiske Grund\-%
anskuelse, ikke for de geognostiske Specialiteter, men for den
hele geognostiske Totalanskuelse, at de paagjældende Viden\-%
skaber skulle staae Theologien til Regnskab.

Men heri stikker just Knuden. I de forskende, iagt\-%
tagende og beregnende Fagvidenskaber maa Opfattelsen af det
Hele nødvendig bestemmes ved Opfattelsen af det Enkelte, og
omvendt. Hvad vil det nu egentlig sige, at en subordineret
Fagvidenskab skal gjøre Theologiens „gode Villie“ Regnskab
for en Totalanskuelse, den ikke skal retfærdiggjøre ved specielle
Undersøgelser? Det tager sig underligt ud! Theologien
sætter sig paa Thronen og siger til den subordinerede Viden\-%
skab: Kom og gjør Regnskab for din Totalanskuelse! Den
Subordinerede nærmer sig og siger med et dybt Buk: „Deres
Majestæt, her er mit Regnskab, mine Analyser, mine Be\-%
regninger!“ Theologien svarer: „Dine Regnskaber og Speciali\-%
teter spørger jeg jo ikke om; men Du skal indestaae mig for en
Totalanskuelse, der er en Frugt, ikke af Kundskabens Træ,
men af Livets Træ“. Den Subordinerede undskylder sig:
„Jeg kan ikke faae Andet ind i min Totalanskuelse, end hvad
der med rationel Nødvendighed fremgaaer af mine Special\-%
undersøgelser; hav Medlidenhed med mig! Min Villie er
ikke ond; men jeg kan ikke Andet.“ --- Den naadige Theo\-%
logie begynder at blive lidt unaadig: „Du fornærmer mig
med din Paatrængenhed. Du veed meget godt, at jeg aldrig
vil nedlade mig til at forstaae dine Regnskaber“.

De, som i dette Skyggerids mene at opdage enten For\-%
dreielser eller Overdrivelser, ville maaskee behage at sige os,
hvori disse Fordreielser, disse Overdrivelser skulle bestaae.
Naturforskerne maae i al Fald kunne vente sig en Forklaring af,
hvorledes Biskop Martensen egentlig har tænkt sig Natur\-%
lærens Afhængighed af „det kongelige Princip“. Sagen er
vigtig; Øieblikket spændende. Vil Biskoppen virkelig gjøre
% --- p. 14 ---
\opage{14}\setcounter{footnote}{0}%
Alvor af at sætte den kirkelig strenge Dom, han har udtalt
over Prometheus og Nutidens Culturbevidsthed, igjennem med
kirkelig Strenghed, da maae Naturvidenskaberne tilsammen
berede sig paa ubetinget Underkastelse.

Naturlæren roser sig af sine store Opdagelser; men disse
store Opdagelser ere jo Frugter af Kundskabens Træ, for\-%
budne Frugter. Hine saa straalende Lys paa Videnskabernes
Himmel, disse glimrende Stjerner, der omgive Navne som
Kopernikus og Tycho Brahe, Galilei og Rømer, Kepler og
Newton, Volta, Ørsted, Laplace med Udødeligheds-Krandse,
hvad ere de vel? Fulgurationer af den prometheiske Gnist,
Flammeblus af hiin Gudernes Ild, Titanen bragte fra
Himlen til Jorden. Hvilken Dom, hvilken streng Dom maa
Biskop Martensen i Conseqvens af sit kongelige Princip ---
paa Monarchiets Vegne --- nu ikke see sig nødsaget til at
fælde over disse syndige Frembringelser, disse Tegn paa for\-%
fængelig Storhed, disse skuffende Luftsyn! Hine Flamme\-%
blus ere jo Blus af kosmisk, ikke af hyperkosmisk, af dæ\-%
monisk, ikke af kanonisk Oprindelse: den prometheiske Gnist
er et Ran fra Guderne; Frugten af Kundskabens Træ er en
forbuden Frugt, hvis Nydelse har bragt Forbandelse over al
Naturen.

Naturlæren er i Safter og Marv, i Stamme og Grene
autonomisk og driver det aldrig til at blive theonomisk. Na\-%
turlæren kan ikke, vil ikke sige \textit{credo, ut intelligam}; den kan
ikke, har aldrig kunnet og vil aldrig kunne udvikle sig af
nogen kirkelig Bekjendelse; Naturlæren staaer og maa blive
staaende som „en verdslig, af Troen reent uafhængig“, Troes\-%
videnskaben modsigende Videnskab, „brammende og prangende
i fuldeste Flor ved Siden af Troen“; Naturlæren, det nytter
ikke at ville skjule det, bliver ifølge det kongelige Princip at
henføre til „den onde Villie, der er udtraadt af Kjærlighedens
Afhængighed“, og desaarsag „maa optænke sig selv en Theorie,
tantalisk bestræbe sig for at sætte en indbildt Virkelighed
% --- p. 15 ---
\opage{15}\setcounter{footnote}{0}%
istedetfor den virkelige“. Vistnok er der noget uhyggeligt
Skærende i en saadan Fordømmelsesdom; men Ingen, som
er istand til at samle to Tanker, skal kunne negte, at Dom\-%
men er en uimodsigelig Conseqvens af „den gode Villie“, det
„kongelige“ Videnskabernes „Princip“, hvormed Biskop Mar\-%
tensen nu vil gjennemføre sit homerisk-aristotelisk-dogmatiske
Universalmonarchie.

Hvad skulle Naturvidenskaberne nu gjøre: underkaste sig
Dommen, eller excipere Forum? „De skulle“, vil man maa\-%
skee sige, „omvende sig; de skulle oprigtigt forsage „den onde
Villie, der er udtraadt af Kjærlighedens Afhængighed“; de
skulle betænke, at naar der siges om Djævelen som Løgnens
Fader, at han ikke blev bestaaende i Sandheden, da har dette
til sin Forudsætning, at han ikke blev bestaaende i Kjærlig\-%
heden og Lydigheden“; de skulle betænke, hvad der „kan sees
af Historien om Menneskets Syndefald og de tvende absolut
ueensartede Træer“, at Modsætningen mellem Godt og Ondt
er den første Grundmodsætning, og at der først kan blive
Tale om „den levende Modsætning mellem Sandhed og
Løgn“, naar hiin er erkjendt som den første; de skulle be\-%
tænke, at saasandt de fremture i deres eensidige Vedhængen
ved det Logisk-Physiske, uden at bøie sig ydmygt under „den
Tænkning, som i Skarphed, Sikkerhed o. s. v. veier op imod
den logisk-physiske, hvis Fremstilling af Miraklet f. Ex. veier
op imod de physicalske Indvendinger“: da vil det gaae dem
ilde; da vil den høiere Videnskab, den Videnskab, der ikke
har ugudelige Kundskaber fra noget Kundskabens Træ, efter\-%
som den, hvad Naturlæren angaaer, slet ingen Kundskaber
har, men alene Gudfrygtighedens \textit{credo} fra Livets Træ,
komme straffende over dem; da vil det kongelige Videnskabernes
Princip, den gode Villie, for at forebygge „Anarchie og
indre Opløsning“, see sig nødsaget til at støde dem ud af
Videnskabernes Samfund. Udstødte af Videnskabernes Sam\-%
fund som af en høiere ethisk Lysregion, nemlig Kjærligheds\-%
% --- p. 16 ---
\opage{16}\setcounter{footnote}{0}%
regionen, kunne de saa med deres under en lavere Viden
hørende ἀσθηνῆ καὶ πτωχὰ στοιχεια % sic: printed ἀσθηνῆ (for ἀσθενῆ) and στοιχεια without circumflex (Gal. 4,9 στοιχεῖα); checked at 600 dpi
(kraftløse og fattige
Børnelærdomme) faae Lov til at sidde i Inductionernes,
Forskningernes og Roderiernes, de sublunariske Forskningers,
de geognostiske Roderiers mørke Hul, i uvidenskabelig Graad
og Tænders Gnidsel.

Det er Naturlæren umuligt at hæve sig over det Logisk-%
Physiske; det er Naturlæren umuligt at lade Theologernes
„gode Villie“ raade for Virkeligheden og bestemme Sand\-%
heden. „Maatte vi“ --- siger den ellers saa fredelige H. C.
Ørsted i dyb Uvillie over det theologisk-biskoppelige Subjec\-%
tivitetsvæsen --- „ikke skamme os i vort Inderste, naar vi
grebe os selv i at ville have en anden Sandhed end den
virkelige?“ Naturlæren kan ikke --- det er dens Constitution
imod --- sige \textit{credo ut intelligam}; den kan ikke bøie Nakken
under en troende Enehersker: den kan fra sin Side ikke give
efter. Men Monarchiets „gode Villie“ kan fra sin Side,
saasandt den virkelig er et Princip, heller ikke give efter.
„Den gode Villie“ maa som \emph{Princip} være conseqvent; den
maa som „\emph{det kongelige Princip}“ gjøre Fordring paa
Overherredømmet; den maa som det \emph{kongelige Viden\-%
skabernes Princip} gjøre Fordring paa \emph{videnskabeligt}
Overherredømme. Magt staaer her imod Magt, „Tegn imod
Tegn“: det trækker op til en antik Tragedie.

Det antikt Tragiske kjendes just derpaa, at to ideale og
forsaavidt lige berettigede Magter støde sammen i uforsonlig
Strid, en Strid, der ender med, at den svagere bukker under
for den stærkere. Bør det nu antages, at det kongelige
Princip i absolut Forstand maa være det stærkere: da seer
det sort ud for Naturvidenskaberne. Et Dommens Uveir
staaer i Horizonten; det er den gamle Tragediernes Tra\-%
gedie, Tragedien med „den bundne Prometheus“, der be\-%
gynder forfra. Endnu engang maa Prometheus, „der har
lært Menneskene Videnskab og hver en Kunst og fremkaldt
% --- p. 17 ---
\opage{17}\setcounter{footnote}{0}%
Culturbevidsthed paa Jorden, lænket til Klippen, bøde for
den Handling, ved hvilken han brød igjennem til sin Kund\-%
skab, fordi han forskaffede sig den ad en Vei, der ikke var
legitim; % Rettelser: "Side 17 Linie 4 fra oven legitim; læs: legitim“;" -- not applied; the closing quote of „der har lært ... (p. 16) is missing in print, so „ exceeds “ by one here. The KB copy shows a thin hand-drawn “ after "legitim" (manuscript, not transcribed).
endnu engang: o haarde Skjæbne! Og vi, som med
det nittende Aarhundredes „Culturbevidsthed“ allerede drømte
om at have seet den befriede Prometheus, seet Lænkerne falde
af Titanens mægtige Lemmer, seet det Blik, han fæstede paa
den døende Glente, seet Herakles med Buen omstraalet af et
saa guddommeligt Lys, at Natten selv, Brutaliteternes Nat,
maatte flye og skjule sig i Klippernes Huler. Drømmen
er ude; ak, det var kun en Drøm, en indbildt Virkelighed,
et Kogleri af den onde Villie. Thi, siger Biskop Martensen:
„den onde Villie, der er udtraadt af Kjærlighedens Af\-%
hængighed, maa optænke sig selv en Theorie, tantalisk be\-%
stræbe sig for at sætte en indbildt Virkelighed istedetfor den
objective“.

I den høie, antike Tragedie er Alt symbolsk. Pro\-%
metheus og Io, Kratos og Bia, Hermes, Okeanos, alle i
Tragedien optrædende Figurer ere symbolske. Prometheus,
Kunsternes Opfinder, Naturvidenskabernes Fader, er Cultur\-%
bevidstheden, den autonomiske Bevidsthed, som ved at udelukke
Troen er bleven „for absolut i Viden“; medens Io, Zeuses
ulykkelige Elskerinde, en, om vi saa maae sige, bissende Sub\-%
jectivitet, der er stukken af Paradoxets Brems, er den Be\-%
vidsthed, som af misforstaaet Theonomie har udelukket Viden
og er bleven „for absolut i Troen“. Kratos og Bia, der
udsendes af Zeus for at binde Prometheus, ere vel efter
Ordene \emph{Vold} og \emph{Magt}, men i Symbolet er det \emph{Kjær\-%
ligheden og den gode Villie}, der udsendes af det kon\-%
gelige Princip for at forebygge „Anarchie og indre Opløs\-%
ning“ i Monarchiet.

Kun een Forandring vil foregaae med den gamle Tra\-%
gedie, en Forandring, der da forsaavidt ogsaa er symbolsk,
som den jo er et Vidnesbyrd om, at „det System af Love
% --- p. 18 ---
\opage{18}\setcounter{footnote}{0}%
og Kræfter, hvilket vi kalde Natur, befinder sig i en teleolo\-%
gisk Udvikling, en fortsat Skabelse“; Forandringen er charak\-%
teristisk: Prometheus vil nemlig denne Gang ikke blive lænket
til en Klippe, og det af den gode Grund, at det kongelige
Princip i den hele naturlige Virkelighed ikke eier en eneste
Klippe; derimod vil han blive bunden til et Træ, et af de
tvende i Paradiis „absolut ueensartede Træer“, det Træ
nemlig, hvis Frugter vare Mennesket forbudne.

Den nye Straffelov vil imidlertid blive ligesaa haard
som den gamle. Thi hvorvel Titanen bliver bunden til
Kundskabens Træ, altsaa det Træ, hvis Frugter smage ham
bedst, har „den gode Villie“, Villien, der alene har „den til
det Gode svarende Theorie“, naturligviis sørget for at faae
ham bagbunden saaledes, at han i sin lange Straffetid ikke
skal blive fed af „den forbudne Æden“, men nøies med
„Propædeutikens magre Kost“, og saaledes nok faae „Kjær\-%
ligheden“ at føle.

Vistnok staaer Naturlæren efter Monarchiets Rangfor\-%
ordning paa et temmelig underordnet Trin blandt Viden\-%
skaberne; thi, siger Biskop Martensen (Leilighedsskr. S. 109):
„Naturvidenskaben har som Erfaringsvidenskab ingen andre
Grunde at støtte sig til end Induction og Analogie. Af en
Række Phænomener, der under de samme Betingelser ved\-%
varende gjentage sig som de samme, udleder man hvad man
kalder en Lov; men det indsees let, at ubetinget Almeengyl\-%
dighed og Nødvendighed ad denne Vei umuligt kan frem\-%
komme“. Men for Naturforskerne stiller Sagen sig ander\-%
ledes. „I min Bog, Aanden i Naturen, har jeg“ --- siger
H. C. Ørsted --- „med fuld Overbeviisning hyldet den gamle
Lære om Naturlovenes Uforanderlighed“. Denne „gamle
Lære“ er for Forskeren saa afgjørende, at Forskningens
videnskabelige Charakteer staaer og falder med dens Gyldighed.

Vi tvivle naturligviis ikke om, at Biskop Martensen
ved Hjælp af det kongelige Princip samt ved \textit{credo ut in\-%
% --- p. 19 ---
\opage{19}\setcounter{footnote}{0}%
telligam} vil kunne føre et strengt Beviis for, at Natur\-%
lovene ikke ere uforanderlige. „Overhovedet“, siger han (S.
109), „veed Naturvidenskaben, som blot Erfaringsvidenskab,
slet Intet om det i \emph{absolut} Forstand Mulige eller Umulige.
Og i Henseende til Virkeligheden og Kjendsgjerningerne be\-%
væger den sig kun i en Midte, indenfor en Kreds af Be\-%
tingelser, der nu engang ere factisk givne, den veed ikke selv
hvorledes, medens Begyndelse og Ende ere indhyllede i
uigjennemtrængeligt Mørke. Endnu bliver den os Svaret
skyldig paa det gamle Spørgsmaal: hvad der var først,
Ægget eller Hønen? Oldenkjærnen eller Egen?“ Vi ville
nu gaae ind paa denne Synsmaade; vi ville endog ind\-%
rømme, at Theologien er den eneste Videnskab, der fra sit
overnaturlige Standpunkt ikke bliver Svaret skyldig, naar der
alvorlig spørges, hvorledes det i Egen og Hønen egentlig for\-%
holder sig med Skabelse og Natur. Ja, det skal end ikke
betvivles, at jo Dogmatiken ved at bestemme det rette For\-%
hold mellem Creatianismen og Traducianismen vil komme
til det uomtvistelige Resultat, at vi oprindelig have det crea\-%
tiansk Naturlige i Egen, det traduciansk Creaturlige i Olden\-%
kjernen, Creaturet (κτισις) i Hønen, Naturen (φυσις) i
Ægget. % Greek unaccented as printed (checked at 600 dpi)
Men alt dette forudsat, saa bliver dog det for
Naturlærens Skjæbne i Monarchiet skjæbnesvangre Spørgs\-%
maal: hvorledes vil Theologien bære sig ad med at overbe\-%
vise Naturforskerne om, at Naturvidenskabens Lærdomme høre
til ἀσθηνη καὶ πτωχα στοιχεια % sic, as printed: ἀσθηνη without circumflex, πτωχα without grave (contrast p. 16 ἀσθηνῆ, πτωχὰ); checked at 600 dpi
(kraftløse og fattige Børne\-%
lærdomme)? endnu tilbage.

Naturforskerne ere stolte af deres Videnskab. Disse med
den dogmatiske Historie om Menneskets Syndefald og „de tvende
absolut ueensartede Træer“ kun lidet fortrolige Mænd faae
det vist aldrig i deres Hoved, at Spørgsmaalet om Natur\-%
lovenes Uforanderlighed skal afgjøres i Theologien. Man
kalde det en Grille, en Fixidee, en Indskydelse af Djævelen,
en Virkning af den onde Villie, der er udtraadt af Kjærlig\-%
% --- p. 20 ---
\opage{20}\setcounter{footnote}{0}%
hedens Afhængighed, eller hvad man vil; men vist er det, at
enhver af Forskerne er sikker paa sin Videnskabs Realitet.
En Astronom f. Ex. holder sig i den Grad forvisset om, at
hans Opfattelse af Stjernehimlen er videnskabelig, at hvis
det kongelige Princip skulde finde sig foranlediget til endnu
engang at opkaste Tvivl om Grundtanken i det kopernikanske
System, vil han være dristig nok til at erklære en saadan
Tvivl for uvidenskabelig. Ingen er villigere end Geologen
til at indrømme, at der med Hensyn paa Jordens Udvik\-%
lingshistorie endnu staaer Meget, uendelig Meget, tilbage at
undersøge; men at Geologien ikke skulde være en Videnskab,
en virkelig og sand Videnskab, det faaer man ham aldrig til
at indrømme; ja, Geologen er i sin videnskabelige Iver endog
istand til at vove sig saa yderlig, at hvis Nogen spørger ham
om den Egn, hvori Paradiis virkelig har ligget, vil han hen\-%
vise den Spørgende til Ammestuen. Hvad her er sagt om
Astronomen og Geologen, gjælder paa tilsvarende Maade om
Botanikeren. Man stille engang en Botaniker den „høiere
videnskabelige“ Opgave, at han i Forening med Geologen skal,
for at sikkre den „\emph{objective} Virkelighed“, undersøge Beskaffen\-%
heden af den Jordbund, hvori „de tvende \emph{absolut ueens\-%
artede} Træer“ maae antages at have voxet, om de begge
have voxet f. Ex. i Engbund eller i Agerland, i leret eller
sandet Jord, eller om det Jordsmon, hvoraf Livets Træ
fremskød, maaskee har været \emph{absolut ueensartet} med det,
hvori Kundskabens Træ slog Rødder: da vil han gjøre store
Øine. Vil man erklære, at hans Videnskab er en „kraftløs
og fattig Børnelærdom“, siden han ikke engang kan bestemme
os de tvende \emph{absolut ueensartede} Træers Plads i Systemet,
men bliver os Svaret skyldig paa det Spørgsmaal, om de
begge høre f. Ex. til Gramineerne, til Rosaceerne eller maa\-%
skee til Nellikefamilien: da er den ellers saa godmodige Bo\-%
taniker endog istand til at blive uartig imod „den gode
Villie“ og forløbe sig i den Grad, at han underdanigst beder
% --- p. 21 ---
\opage{21}\setcounter{footnote}{0}%
„det kongelige Princip“ med den hele dogmatiske Træskole reise ad
Helvede til. At et saadant Udbrud er høist upassende, maae
vi indrømme; men Sligt kan jo skee i Fortvivlelse. Natur\-%
forskerne blive saa fortvivlede, naar en Dogmatiker fra Theo\-%
nomiens overnaturlige Høide med Overlegenhed seer ned paa
deres Videnskabs naturlige Autonomie.

Det vil blive det „kongelige Princip“ en haard Dyst, inden
det faaer „Naturalismen“ grundig udryddet af Naturlæren
og Forskerne overbeviste om, at der virkelig kan gives en
Videnskab, hvis Grunderkjendelser og første Forudsætninger
befinde sig i aabenbar Strid, med hvad de maae erklære for
vitterlige Kjendsgjerninger, en Videnskab, hvis „Fremstilling
af Miraklet veier op imod de physicalske Indvendinger“, en
Videnskab, der paa Grund af „den gode Villie“ gjør det
objectivt uimodsigeligt, at Skabelsen ender med et Paradiis,
at der forud for Syndefaldet ikke har været Spor af For\-%
krænkelighed, Smerte og Død paa Jorden. Hvad disse
Empirikere, navnlig Geologerne, dog kunne hitte paa for at
give deres „kraftløse og fattige Børnelærdom“ et Slags Skin
af virkelig og sand Videnskab! Istedenfor at begynde, hvad
de kalde Jordklodens Historie, med den rette historiske, nemlig
den overhistoriske Begyndelse, med Paradiset, „de tvende absolut
ueensartede Træer“ og „den forbudne Æden“, begynde de,
tænk engang, med Graavakken! altsaa med „at optænke sig
selv en Theorie, tantalisk bestræbe sig for at sætte en indbildt
Virkelighed istedetfor den objective“.

Ved en Række uhyggelige, for Dogmatikeren lidet opbyg\-%
gelige Billeder søge Nutidens Geologer at anskueliggjøre en
ingenlunde objectiv, men tantalisk indbildt Virkelighed af en
umaalelig Urtid, idet de --- naturligviis kun hen i Luften til
deres egen „og de Uvidendes Fornøielse“ --- paa naturalistisk
Maade stræbe at objectivere sig Dyrelivet i de forskjellige
Jordperioder ligefra Graavakken og Kulkalkstenen indtil Kridt- og
Molassetiden, bestandig under den tilsnegne Forudsætning,
% --- p. 22 ---
\opage{22}\setcounter{footnote}{0}%
at endogsaa Kridt- og Molassetiden gaaer forud for Synde\-%
faldet. For at gjøre os et Begreb om den Naturalisme,
der kan aabenbare sig i slige tantaliske Objectiveringer, be\-%
høve vi blot at kaste et Blik paa et af disse Billeder, det
Billede f. Ex., der viser os „Jurahavets og dets Kysters Be\-%
boere i fuld Virksomhed.“ Hvor tantalisk indbildt er ikke
denne Virksomhed! „I Luften“, siger Geologen\footnote{Geologiske Billeder af Bernhard Cotta. Efter Originalens tredie
Oplag ved C. Fogh. Kjøbhvn 1859. S. 243 flgd.}, „seer
man to store Flyveøgler af Slægten Pterodaktylus, maaskee
to iversyge Rivaler, der i Kampens Hede have ladet Kamp\-%
prisen af Øie, thi den stakkels Pterodaktylushun falder netop
som Bytte for en langhalset forsulten Plesiosaurus, der har
faaet fat i dens ene Vinge, og snart vil gjøre Ende paa
den trods dens ængstelige Skrig. For Millioner Aar siden
ere dens Angstskrig døde hen i Luften,“ o. s. v.

At den hele Opfattelse af Juraperioden i et saadant
Billede maa beroe paa et uhyre Sandsebedrag, kan jo Dog\-%
matiken i Conseqvens af det kongelige Princip udtrykkelig
bevise. Det kan nemlig ved Frugter, der hentes, ikke fra
Kundskabens, men fra Livets Træ, videnskabeligt godtgjøres,
at „Naturalismen“ her gjør sig skyldig i et uhyre Hysteron-%
Proteron; thi det kan bevises, at en saadan Dyrekamp, der
kun kan foregaae i en af Herrens Forbandelse ramt Natur,
umulig har kunnet finde Sted før Syndefaldet. Her staaer
Slaget; her er just Naturalismens Achilleshæl, et af de
svage Punkter, hvorfra Troesvidenskaben seierrig vil kunne
angribe Vantroens Børnelærdom.

Beviset er ved \textit{credo, ut intelligam} uigjendriveligt og
kan fattes af enhver, der er istand til at følge dets Tanke\-%
gang. Før Syndefaldet var Alnaturen i sin Uskyldighed
theonomisk; først efter Syndefaldet --- det kan sees af „den
christelige Dogmatik“ --- blev den formedelst Slangen og det
% --- p. 23 ---
\opage{23}\setcounter{footnote}{0}%
kosmiske Princip autonomisk. Men nu betragte man hvilke
Partier af hiint Jurabillede man vil, i alle Scener, fra alle
Sider vise sig Kjendsgjerninger, der vidne om, at Auto\-%
nomien allerede har fortrængt Theonomien, at det Værende
bliver regjeret slet. See der længere fremme i Billedet
ligger en mægtig Krokodil! Allerede Nærværelsen af en
Krokodil varsler om, hvordan det Værende her bliver regjeret.
Krokodilen ligger, som man seer, under Bananer, Cycadeer
og Bregnetræer; den vil da, kan jeg tænke, heller ikke nøies
med „Propædeutikens magre Kost“, men lurer øiensynlig
paa en velsmagende Dessert. At en saadan Dessert, hvor
velsmagende den for Resten kan være, absolut maa falde, ikke
\emph{før} Syndefaldet, men \emph{efter} Syndefaldet, efterdi det,
nærmere beseet, er en syndig, nemlig en objectiv-natur\-%
% line-end hyphen after "natur" taken as a word break (objectiv-natursyndig); it may be a compound hyphen (natur-syndig)
syndig Dessert, kan exegetisk udledes af Contexten. Af
det umiddelbart Følgende sees nemlig, at en velpantsret
Havskildpadde er --- naturligviis autonomisk, ikke theonomisk
--- kommen Krokodilen i Forkjøbet og har grebet et prægtigt
Belemnitdyr, som den er i Begreb med at slæbe paa Land.
Allevegne Spor af Forkrænkelighed, Spor af Synd, af
Forbandelse! Længere henne sluger en Ichtyosaurus --- og % a stray raised dot stands before "Ichtyosaurus" in print (spacing material); not transcribed
det skulde være skeet før Syndefaldet? --- en heel Blæk\-%
sprutte, medens en anden Blæksprutte --- naturligviis efter
Forbandelsen --- tilligemed en korthalet Krebs søger Frelse i
iilsom Flugt. Er, sige vi, alt dette Virkelighed; er det, hvad vore
Naturalister jo paastaae, virkelig engang skeet, og har det,
som allerede viist, umulig kunnet skee \emph{før} Syndefaldet: saa
maa det nødvendigviis være skeet \emph{efter} Syndefaldet. \textit{Quod
erat demonstrandum}: her er Beviset.

Fra Theonomiens overnaturlige Standpunkt har Theo\-%
logien jo vistnok Ret; men hvordan vil det kongelige Princip
bære sig ad med at faae Naturforskerne løftede op paa
Theonomiens Standpunkt? Det er, i Alvor, et stort og
vigtigt Skridt, Biskop Martensen ved denne Leilighed har
% --- p. 24 ---
\opage{24}\setcounter{footnote}{0}%
besluttet sig til, idet han ved at udtale Dommen og bryde
Staven over Naturlærens prometheiske Forudsætninger har
udfordret „Culturbevidstheden“ til Kamp paa Liv og Død.
Hvad saa end Udfaldet bliver, skal det indrømmes ham, at han
ved sin Optræden paa Kamppladsen har, om end ikke umid\-%
delbart, saa dog middelbart, leveret et Bidrag, der blotter
Problemets inderste Kjerne.

Ogsaa vi have fra vor Side længe været opmærksomme
paa den under Asken ulmende Glød og med stigende Spæn\-%
ding ventet paa det Øieblik, da den vilde slaae ud i Lue.
Vi have desaarsag paa bedste Maade stræbt at gjøre os Rede,
ikke blot for Forholdet mellem Tro og Viden i speculativ
Almindelighed, men, at vi ikke skulde forfeile den concrete
Virkelighed, udtrykkelig for Forholdet mellem Religion og
Naturlære. Vi have, det forstaaer sig, ikke kunnet svinge os
saa høit som Biskop Martensen, og ere med vore Under\-%
søgelser komne til mere spagfærdige Resultater.

Naar vi da spurgte os selv: hvorledes er det muligt, at
to saa ueensartede Principer som Tro og Viden kunne for\-%
enes i een Bevidsthed, da tænkte vi herved ligesaa meget paa
den naturvidenskabelige som paa den philosophiske Bevidsthed.
Kan Naturforskeren ikke forene Tro og Viden i een Bevidst\-%
hed, da kan Philosophen det heller ikke; kan Naturforskeren det,
saa kan Philosophen det ogsaa. At Theologer og Præster
let dresseres til at forene Tro og Viden, kunne vi vel for\-%
staae; men med Naturforskeren er det en egen Sag. Vi
have under Problemets Udvikling holdt Øiet fæstet paa de
naturvidenskabelige Vanskeligheder og det af gode Grunde.
For det Første finde vi hos Naturforskeren, hvad man just
ikke altid finder hos Theologer og Præster, Kjendskab til og
desaarsag Respect for en virkelig Natur med virkelige Kjends\-%
gjerninger og virkelige Love. For det Andet kunne vi med
temmelig Sikkerhed forudsætte om Naturforskeren, hvad man
ikke saadan uden videre tør forudsætte om Theologer og
% --- p. 25 ---
\opage{25}\setcounter{footnote}{0}%
Præster, at dersom han virkelig giver sig af med at tænke
over religiøse Problemer, da maa hans Tænkning have en
egen personlig Oprindelighed; den beroer da ikke paa officiel
Dressur, men udspringer af uimodstaaelig indre Trang.

Ved Tilbageblik paa Naturvidenskabernes Fremskridt og
Culturens Gang i de to sidste Aarhundreder meente vi, hvad
de religiøse Naturforskeres Opfattelse af Tro og Viden an\-%
gaaer, at have opdaget en umiskjendelig Stræben efter at
lade Adskillelsen mellem en religiøs og videnskabelig, praktisk
og theoretisk Erkjendelse træde i Kraft. Vi meente at finde
en saadan Adskillelse, vel ikke theoretisk udviklet, men dog
factisk anerkjendt f. Ex. af Newton. Newtons Mechanik
bærer ligesaa lidt Spor af Troens overnaturlige „gode
Villie“ som d'Alemberts eller Laplaces; og dog var New\-%
ton jo, som bekjendt, en religiøs Mand, en Mand, der
skrev Anmærkninger til Propheten Daniel, forklarede Apoca\-%
lypsen og troede paa Mirakler og Spaadomme.\footnote{For Newton har det prophetiske Ord umiddelbar guddommelig
Autoritet. \textit{Auctoritas Imperatorum, Regum, Principum humana
est: Conciliorum, Episcoporum auctoritas humana. Prophetarum
auctoritas divina est \dots\ Isaaci Newtoni Opuscula Lausannæ et
Genevæ MDCCVLIV. Observationes ad Danielis prophetæ
vaticinia p. 291.} % sic MDCCVLIV (for MDCCXLIV), as printed
Hvad Udtydningen af Spaadommene angaaer,
advarer Newton vel imod at gaae formeget i det Enkelte, men
tilføier dog. Anfr. Skr. \textit{p. 292: Rerum futurarum prædictiones
Ecclesiæ conditionem per omnia secula respiciunt.} Hvor umid\-%
delbart han har opfattet de apocalyptiske Visioner, sees af følgende
Udtalelse: \textit{Post paratas lucernas, vidit Joannes apereri januam
Templi ac voce, quasi buccinæ, evocabatur ad eam magni atrii
portam, quæ Orientem spectat, ut illic visa videret. Opuscul.
XXV. Observat. in Apocalyp. p. 451.} Det er af disse og
lignende Udtalelser klart, at Newton i det Hele maa ansee Er\-%
kjendelserne af det Naturlige og Erkjendelserne af det Overnatur\-%
% [printed p. 26: the note continues at the foot]
lige for ligesaa ueensartede som det Naturlige og det Overnatur\-%
lige. Hører Erkjendelsen af det Naturlige nu ind under Viden\-%
skaben, og kan der ikke gives absolut ueensartede Videnskaber: saa
følger jo heraf, at Erkjendelsen af det Overnaturlige maa være en
Erkjendelse til praktisk Brug, en praktisk Erkjendelse.} Havde
% --- p. 26 ---
\opage{26}\setcounter{footnote}{0}%
imidlertid nogen Physiker, medens Newton endnu var ved
fuld Kraft, vovet at angribe \textit{Principia philosophiæ natu\-%
ralis mathematica} ved Argumenter, hentede fra hvad Gud
i sit Ord antages at have sagt: da vilde den gudfrygtige
Newton vist ikke have taget i Betænkning at afvise et saa
taabeligt Forsøg med Ord, der let kunde tydes til samme
Underviisning som Slangens Ord: „Det er ikke sandt, hvad
Gud har sagt (nemlig fra Videnskabens Standpunkt)“.

Vilde i vore Dage en Bibeltheolog imødegaae Optikens
Forklaring af Regnbuens Dannelse ved at anføre Guds
Ord til Noah (1 Moseb. 9, 13): „Min Bue vil jeg sætte
paa Skyen, og den skal være et Pagtens Tegn imellem mig
og imellem Jorden,“ og tilføie, at Regnbuen udtrykkelig
skyldes et Mirakel, efter hvad Gud selv har sagt: vilde da
ikke selv den gudfrygtigste af alle videnskabelige Optikere give
et Svar, der kunde udlægges til den Underviisning: „Det
er ikke sandt, hvad Gud har sagt (nemlig fra Videnskabens
Standpunkt)“.

Lad en Dogmatiker, der med det kongelige Princip saa
fortræffeligt veed at forsvare Underets Mulighed, engang for\-%
søge at forklare gudfrygtige Zoologer Muligheden af Miraklet
med Bileams Aseninde (4de Moseb. 22, 21 flgd.)! „Og
Bileam stod op om Morgenen og sadlede sin Aseninde og
gik med Moabs Fyrster. \dots\ Men Jehovas Engel stod i
en Huulvei mellem Viingaardene. \dots\ Og da Aseninden saae
Jehovas Engel, da lagde den sig ned under Bileam; saa
blev Bileam vred og slog Aseninden med Kjeppen. Da aab\-%
nede Jehova Asenindens Mund, og hun sagde til Bileam:

% --- p. 27 ---
\opage{27}\setcounter{footnote}{0}%
hvad har jeg gjort Dig, at Du har slaaet mig nu tre Gange.
\dots{} er jeg ikke Din Aseninde, som Du har redet paa fra den
Tid Du var til, indtil denne Dag“ o. s. v. Hvorledes vil
Dogmatikeren med det kongelige Princip nu bære sig ad
med at fremstille dette Mirakel saaledes, at Fremstillingen
„veier op mod de physicalske Indvendinger“? Beraaber han
sig paa, hvad Engelen har sagt, saa svare Zoologerne som
med een Mund: „det er ikke sandt, hvad Engelen har sagt
(nemlig fra Videnskabens Standpunkt)“. Beraaber han sig
paa, hvad Aseninden har sagt, saa svare de i Chor: „det er
ikke sandt, hvad Aseninden har sagt (nemlig fra Videnskabens
Standpunkt)“. Beraaber han sig paa, hvad Gud i sit Ord
har sagt, saa blive de ved: „det er ikke sandt, hvad Gud har
sagt (nemlig fra Videnskabens Standpunkt)“. % a small raised point after "nemlig" (an inked space, as on p. 31) is not transcribed

Hvad Naturforskeren indvender mod alle særskilte Mirakler
i den hele Bibelhistorie, indvender han ogsaa imod Grund\-%
miraklerne: Skabelse, Incarnation og Inspiration. „Det
er ikke sandt, at Gud i Begyndelsen har skabt Himlen og
Jorden (nemlig fra Videnskabens Standpunkt); det er ikke
sandt, at Guds eenbaarne Søn er undfangen af den Hellig
Aand og født af Jomfru Maria (nemlig fra Videnskabens
Standpunkt); det er ikke sandt, at den Hellig Aand paa
Pintsefesten er kommen i Stormveir og Ildsluer og har i et
eneste Nu lært Apostlene de mange fremmede Tungemaal
(nemlig fra Videnskabens Standpunkt)“.

For ikke at forhaste sig med nogen endelig, uigjenkaldelig
Dom over disse Udtalelser, gjør „den gode Villie“ udentvivl
bedst i at gjentage for sig selv, hvad det er for et Stand\-%
punkt, hvorfra her tales, at det bestandig er fra \emph{Viden\-%
skabens}, særlig \emph{Naturvidenskabens} Standpunkt. Paa det
naturvidenskabelige Standpunkt er det nemlig vedtaget, hvad
der ikke altid er Skik og Brug blandt Theologer og Præster,
at holde strengt over Adskillelsen mellem \emph{hvad man veed},
og \emph{hvad man ikke veed}. Det ligger i Sagens Natur,
% --- p. 28 ---
\opage{28}\setcounter{footnote}{0}%
at den, der veed Noget tilgavns, vil mærke sig Grændsen
mellem Viden og Ikke-Viden. Ved at besinde sig paa den
menneskelige Videns Grændser vil en Naturforsker let komme
til følgende, i den philosophiske Propædeutik (S. 67) udtalte
Resultat: „I Kraft af den empiriske Forsknings indre oprinde\-%
lige Sammenhæng med det Aprioriske gives der en væsentlig
Viden; men ifølge Tænkningens og Sandsningens Ueens\-%
artethed er Realerkjendelsen underlagt endelige Betingelser, og
al vor Viden begrændset. Videns Grændse er en dobbelt:
en \emph{relativ}, forsaavidt her altid er Noget, vi endnu ikke
vide, men som vi dog med Tiden kunne faae at vide; en
\emph{absolut}, forsaavidt en endelig Adskillelse af Bevidsthedens
og Virkelighedens Indhold er begrundet i vor Organisation:
her er altsaa Noget, vi i denne Tilværelse aldrig faae at
vide.“

Af Forholdet mellem det Relative og Absolute i vor
Videns Begrændsning vil Forholdet mellem Tro og Viden
kunne opklares. Havde den menneskelige Viden kun relative
Grændser, da maatte der naturligviis altid kunne gives „\textit{plus
ultra}“, da maatte alle vore Erkjendelser uden Forskjel om\-%
sider blive at begrunde i Videnskaben. Hvad Videnskaben
paa disse Vilkaar ikke kunde begrunde, hvad den fandt selv\-%
modsigende og forkasteligt, blev da ikke blot usandt fra Vi\-%
denskabens Standpunkt, men det blev med det Samme at
ansee for usandt i sig selv: i denne Forstand vilde det ganske
rigtig være en Umulighed, at Noget kunde være sandt i
Troen, som ikke var sandt i Viden. Har den menneskelige
Viden derimod foruden de relative Grændser en absolut % a stray point stands before "relative" (a speck or raised quad), not transcribed
Grændse, saa at her altsaa er Noget, vi i denne Tilværelse
aldrig faae at vide: da er her et Mysterium, en Tilværel\-%
sens Gaade, hvis Opløsning i Videnskaben maa blive usand.
Til dette Mysterium maae henføres det personlige Livs ab\-%
solut ideale Forudsætninger: at en hellig almægtig Villie
er Altilværelsens Ophav; at Menneskets Væsen, ethisk seet,
% --- p. 29 ---
\opage{29}\setcounter{footnote}{0}%
befinder sig i en oprindelig Strid med sit Ideal; at ethvert
Individ, der slaaer Øiet op i Aandens Rige, maa erkjende
sig \emph{skyldig}; at Hjertet først kan finde Fred, naar Skylden
er indseet, tilstaaet og sonet, og Vished om et evigt Liv
i Personlighedens og Sandhedens Rige erhvervet. Jo mere
Samvittigheden skærper disse ideale Forudsætninger, desto
dybere er den religiøse Trang; jo dybere, stærkere og inder\-%
ligere den religiøse Trang er, desto mere absolut personlig
er Afgjørelsen, og jo mere absolut personlig Afgjørelsen er,
desto tydeligere indsees det, at det personlige Livs absolut
ideale Forudsætninger maae haves i Tro, ikke i Viden.

Men hvorfor nu ikke \emph{baade} i Tro \emph{og} i Viden, \emph{først}
i Tro \emph{og saa} i Viden, først \textit{credo} og saa \textit{intelligam}?
Ja hvorfor ikke? Det er i Reglen lettere at gjøre Mod\-%
sætningen mellem Tro og Viden indlysende for Naturforskere
og religiøse Individualiteter end for Theologer og speculerende
Præster. Den religiøse Individualitet \emph{føler}, hvad Tro er;
Naturforskeren \emph{veed}, hvad Viden er; men de, hvis Tro
har en Bismag af Viden, vide ikke, hvad \emph{Tro} er, og de,
hvis Viden indtil det Utrolige smager af Tro, vide ikke,
hvad \emph{Viden} er.

For Naturforskeren er Troen absolut ueensartet med
Viden; heri har han baade en Fristelse til at opgive og en
Mulighed for at bevare Troen. Hvorledes han bestaaer i
denne Fristelse, hvorledes han personlig benytter denne Mu\-%
lighed, derfor skal han i sin Samvittighed staae Ethiken til
Regnskab. Vi kunne i denne Henseende underskrive Biskop
Martensens Ord, at Villien ethisk vil „udvælge eller forkaste
det Theoretiske, der fremstiller sig for Bevidstheden, forbe\-%
holde sig det afgjørende Ja og Nei, Placet og Veto, efterdi
den selv skal være ansvarlig og aflægge Regnskab, ogsaa for
hele den indre Huusholdning“. Men lader os nu særlig
betragte først Fristelsen og dernæst Muligheden. Fristelsen
er, at gjøre Tro og Viden noget eensartede og lade Erkjen\-%
% --- p. 30 ---
\opage{30}\setcounter{footnote}{0}%
delsen af Videns absolute Grændse træde i Baggrunden.
At Troen gjøres eensartet med Viden, kjendes derpaa, at
der af Videns Forudsætninger udledes Troesconseqvenser
og af Troens Forudsætninger Vidensconseqvenser. Saa ofte
Naturforskeren hildes i denne Fristelse, mærker han strax,
hvor urimelig Troen er, hvorlidt den lader sig forene med
Viden.

Vi vælge som Exempel Forholdet imellem Natur og
Skabelse. Man vil af Troens Forudsætninger forgjæves
anstrenge sig med at uddrage Conseqvenser til Bedste for
Videnskaben; thi Skabelsesmiraklet ophæver Videnskabens
Naturbegreb. Man vil af Videns Forudsætninger forgjæves
anstrenge sig med at udvikle Conseqvenser til Bedste for
Troen; thi Videns Naturbegreb gjør Skabelsesmiraklet til
en Phantasie. „At anvende Skabelsestanken til en For\-%
staaelse af det Skabte vil sige: at see den hele Verden under
det dobbelte Synspunkt: \emph{Intet og dog virkelig Noget}.
Det Begreb, hvori denne Modsigelse løses, er Begrebet Al\-%
magt; det Skabte er i sig et Intet, men paa Almagtens
Bud er det blevet til Noget.“ Hermed er Troens Forud\-%
sætning sat, men det videnskabelige Naturbegreb er hævet;
thi hvorledes skal man her tænke sig Overgangen fra Ska\-%
belse til Naturfrembringelse? „Et Forhold, der er begyndt
med et Mirakel, maa ogsaa fortsættes ved et Mirakel: ikke
blot i nogle Øieblikke, ikke blot om vi saa maa sige, en 8
à 14 Dage efter Skabelsen, men i Aarhundreder, i Aar\-%
tusinder, i Evigheders Evighed, maae de skabte Ting vedblive
i sig at være, hvad de oprindelig ere og have været, et Intet,
et reent Intet for den Almægtige.“ Forsaavidt Vandet er
skabt, har det ikke blot en mirakuløs Oprindelse, men en
mirakuløs Tilværen; paa Almagtens Bud er det et virkeligt
Noget, og dog for Almagten selv, den almægtige Villie, et
Intet: \emph{der er ingen Natur i det skabte Vand}. For\-%
saavidt Vandet har en naturlig Tilværelse, forsaavidt der
% --- p. 31 ---
\opage{31}\setcounter{footnote}{0}%
med andre Ord er Natur i Vandet, fremkommer det ved
chemisk Forening af Ilt og Brint, har altsaa ikke blot en
naturlig Tilværelse, men en naturlig Oprindelse; \emph{men den
naturlige Oprindelse er jo ingen Skabelse}. % small inked points after "Ord", "Brint," and "naturlig" (spacing material) are not transcribed

Som Tilværelsen er, saaledes er Oprindelsen: det
videnskabelige Naturbegreb ophæver Skabelsen; som Oprin\-%
delsen er, saaledes er Tilværelsen: Skabelsen ophæver Natur\-%
begrebet. „Er Adam skabt, saa er Adam ingen Natur, saa
er der ingen Natur i Adam. Er der ingen Natur i Adam
og Eva, saa kunne Adams og Evas Børn heller ikke \emph{fødes};
de maae alle \emph{skabes}.“ (Om Theologiens Naturbegreb med
særligt Hensyn til Malebranche). Er der en virkelig Men\-%
neskenatur, altsaa \emph{Natur}, i Adams og Evas Efterkommere,
saa har der ogsaa været en virkelig Menneskenatur, altsaa
\emph{Natur}, i Adam og Eva, saa have Adam og Eva været
virkelige Naturindivider, havt en virkelig naturlig Tilværelse.
Men dersom Adam og Eva have havt en naturlig Til\-%
værelse, da have de ogsaa havt en ligesaa naturlig Oprin\-%
delse som Nogen af deres Efterkommere. Af det Mirakuløse
kan det Naturlige ikke afledes; af det Naturlige kan det
Mirakuløse ikke afledes: det Mirakuløse er absolut ueens\-%
artet med det Naturlige.

Fra dette Synspunkt vil en troende Individualitet let
fristes til at bryde Staven over Naturlæren for at bevare
Troen, medens Naturforskeren omvendt føler sig fristet til
at opgive Troen for at hævde Videnskaben. Fristelsen er, at
Tro og Viden saa let gjøres lidt eensartede, uagtet de ere
absolut ueensartede. Betones det i Ueensartetheden Absolute
tilstrækkeligt skarpt, saa høres en skrigende Modsigelse: \emph{hvad
der er Sandhed i Troen er Usandhed i Viden\-%
skaben; hvad der er Sandhed i Videnskaben er
Usandhed i Troen}. Men i Modsigelsen selv ligger Mu\-%
ligheden til dens Løsning.

Det er vor Videns absolute Grændse, det er Tilværel\-%
% --- p. 32 ---
\opage{32}\setcounter{footnote}{0}%
sens Gaade, hvorpaa saavel den Videndes som den Troendes
Opmærksomhed her maa fæste sig. Har al menneskelig Viden
en absolut Grændse, er her virkelig et Aandens Mysterium,
en Tilværelsens Gaade, da er her ogsaa Noget, der maa
\emph{troes}, fordi det ikke kan \emph{vides}. Naar nu det i Gaaden
Skjulte, det, der ene og alene er bestemt til at være Gjen\-%
stand for Tro, ikke desmindre drages ind under Viden og
Videnskab, hvad saa? Saa overskrider Viden sin absolute
Grændse. Og naar Viden har overskredet sin absolute
Grændse, hvad saa? Saa giver Videnskaben sig frisk væk
ifærd med at udgrunde, begrunde og begribe det, som ifølge
sin Natur hverken kan eller skal videnskabeligt begribes. Men
naar Videnskaben giver sig af med at begribe det Ubegribe\-%
lige, hvad saa? Saa kommer den nødvendigviis til mod\-%
sigende Begreber. Hvad der kun kan udtrykkes i modsigende
Begreber maa være usandt; Troens Mysterier lade sig kun
udtrykke i modsigende Begreber: altsaa ere de usande. At
Noget kan være sandt i Troen og usandt i Videnskaben,
følger nu ligefrem deraf, at \emph{Videnskaben ved at ville
begribe, hvad der ligger udenfor dens Grændse,
selv bliver usand}.

Men hvorledes kan Noget, der virkelig er sandt i Viden\-%
skaben, være usandt i Troen? Svar: Fordi det ikke blot
er Viden, men ogsaa Troen, der har en absolut Grændse:
at Viden har Troen til Grændse, forudsætter nødvendigt,
at Troen har Viden til Grændse. Men at Troen har
Viden til Grændse, vil sige, at et Troesbegreb umuligt kan
udvikles til et Vidensbegreb. Idet Troen støtter sig til Mi\-%
raklet, forudsætter den en virkelig Natur med virkelige Love.
Men hvorledes forholder Naturen sig til Miraklet? Ja det
er jo netop dette, der er Tilværelsens Gaade. Hvis nu
Troen, forglemmende sin absolute Forudsætning, nemlig Gaa\-%
den, overskrider sin Grændse, og begynder at ville omkalfatre
det videnskabelige Naturbegreb til Bedste for Miraklet: da
% --- p. 33 ---
\opage{33}\setcounter{footnote}{0}%
bliver Naturlærens Naturbegreb ikke blot en Usandhed i
Troen, men \emph{Troen, der har indladt sig paa Noget,
den ikke magter, bliver selv usand}.

At Tro og Viden lade sig forene i een Bevidsthed, ikke
\emph{endskjøndt}, men \emph{fordi} de ere absolut ueensartede Prin\-%
ciper, vil da med andre Ord sige, at de kun lade sig forene
i et \emph{absolut Mysterium}. Det absolute Mysterium er paa
een Gang exclusiv Grændse og negativ Midte. Som exclusiv
Grændse sætter Mysteriet Skilsmisse imellem Tro og Viden
som imellem to Regioner. Hver Region har sit Ubetingede,
sit Midtpunkt: Troens Midtpunkt er Friheden, den almæg\-%
tige, hellige Villie; Videns Midtpunkt er Nødvendigheden,
den rationelle, i uforanderlige Love sig aabenbarende Nød\-%
vendighed. Naar nu den Troende seer mod Troens, den
Vidende mod Videns Midtpunkt, ere Extremerne spændte til
det Yderste: hvad der er sandt i Troen, at nemlig det
Ubetingede er Miraklet, er usandt i Viden; hvad der er
sandt i Viden, at nemlig det Ubetingede maa søges i For\-%
nuftnødvendigheden, er usandt i Troen. Som negativ Midte
forener Mysteriet derimod de adskilte Extremer og viser sig som
Regionernes Gjennemgangsmedium: gjennem Mysteriet gaaer
Veien fra Tro til Viden, fra Viden til Tro. Idet den
Troende og den Vidende begge paa een Gang rette deres
Blik mod Mysteriet, udlignes Striden, og Extremerne enes.
Seet fra Videns Standpunkt er nu Miraklet muligt. Thi
vel er det umuligt, at der nogensinde skulde kunne skee et
vilkaarligt Brud paa de evige med Fornuftnødvendigheden
identiske Naturlove; men med Lovenes Uforanderlighed kan
jo, naar de i Mysteriet skjulte Mellembestemmelser underfor\-%
staaes, den Aabenbarelse af Villiens, Viisdommens og Kjær\-%
lighedens Magt, hvori den Troende seer Miraklet, meget
vel tænkes at bestaae. Seet fra Troens Standpunkt kan
Fornuftnødvendigheden med de uforanderlige Naturlove for\-%
enes med Forudsætningen af Miraklet. Thi vel er det
% --- p. 34 ---
\opage{34}\setcounter{footnote}{0}%
umuligt, at Nødvendigheden skulde kunne hemme Friheden,
og Naturens Love indskrænke Guds ubetingede Villie; men
naar de i Mysteriet skjulte Mellembestemmelser underfor\-%
staaes, lader det sig jo tænke, at Naturlovene, langt fra at
brydes, tvertimod i allerhøieste Forstand opfyldes, hver Gang
en guddommelig Villiesact fuldbyrdes.

Af denne Forsoning følger:

1) \emph{at en Troesvidenskab (Theologie) er utæn\-%
kelig.} % the numeral "1)" is not letterspaced

Skulde en Troesvidenskab være tænkelig, da maatte
Troesbegreber lade sig forvandle til Vidensbegreber, og om\-%
vendt; men det er, ifølge disse Begrebers absolut ueensartede
Charakteermærke, umuligt. Vidensbegrebet er mærket af
Nødvendigheden, den rationelle Nødvendighed, hvori Friheds\-%
bestemmelsen --- paa Grund af Mysteriet --- er et forsvin\-%
dende Moment; Troesbegrebet derimod er mærket af Fri\-%
heden, den guddommelige overgribende Villiesact, hvori Nød\-%
vendighedsbestemmelsen --- paa Grund af Mysteriet --- uop\-%
hørligt forsvinder. Begreberne: Skabelse, Paradiis, Synde\-%
fald o. desl. ere Troesbegreber, ikke Vidensbegreber; hvorimod
Begreberne: Graavakke, Juraperiode, Kridtperiode ere Vi\-%
densbegreber, ikke Troesbegreber. Den Art Virkelighed, der
maa tillægges Paradiset og Syndefaldet er ligesaalidt at søge i
Jura og Kridt, som den Virkelighed, der maa tillægges Jura
og Kridt, er at søge i Paradiis og Syndefald. Graavakke,
Jura og Kridt høre til en anden Region end Skabelse, Pa\-%
radiis og Syndefald. Det vil fra dette Synspunkt kunne
skjønnes, hvormegen Vægt, der kan lægges paa Biskop
Martensens Ord, naar han (S. 133) siger om min Philo\-%
sophie: „Hos ham“ (Prof. N.) „staaer en saakaldt moderne
Bevidsthed,“ --- om Graavakke --- „staaer en verdslig af
Troen reent uafhængig og Troen“ --- udi Vedkommendes
egen Indbildning --- „modsigende Viden“ --- om Dyrelivet
% --- p. 35 ---
\opage{35}\setcounter{footnote}{0}%
i Juraperioden --- „brammende og prangende i fuldeste Flor
ved Siden af Troen.“

2) \emph{at der imellem Troen og Viden ikke kan
tænkes noget Concordat.}

Et Concordat, d. v. s. en endelig vilkaarlig Overeens\-%
komst, kan kun tænkes mellem relativt eensartede Magter,
der paa udvortes og endelig Maade indskrænke hinanden:
Indskrænkningen forudsætter en udenfra kommende Begrænds\-%
ning. Men Tro og Viden kunne --- paa Grund af My\-%
steriet --- umulig indskrænke hinanden. Vel er der i Troen
en Grændse for Viden, forsaavidt det, der kan troes, ikke
kan vides; men denne Grændse er Eet med Videns egen
iboende Grændse. Med sin iboende Grændse, sin indre
Begrændsning, er Viden ligeoverfor Troen aldeles uind\-%
skrænket, saa uindskrænket, som om der slet ingen Tro var;
ligeledes er Troen med sin iboende Vidensbegrændsning
uindskrænket ligeoverfor Viden, saa uindskrænket, som om der
slet ingen Viden var. Saalidt som Viden nogensinde er
istand til at angribe Troen, saalidt vil Troen nogensinde
være istand til at angribe Viden; saalidt som Viden er
istand til at rette sig efter Troen, ligesaa lidt er Troen
istand til at rette sig efter Viden. Kunne de nu, hvad det
absolute Mysterium absolut forhindrer, hverken stride mod
hinanden eller understøtte hinanden, hverken gavne eller skade
hinanden, saa er jo en endelig Overeenskomst imellem dem
utænkelig.\footnote{Om Biskop Martensens Opdagelse af Concordatet og hans Phan\-%
taseren over de mange Concordater vil der blive forhandlet i tredie
Afdeling.}

3) \emph{at Adskillelsen af Tro og Viden ikke kan
halvere den menneskelige Bevidsthed.}

Bevidstheden halveres ved relativt eensartede, men ikke
ved absolut ueensartede Factorer. Naar en Bevidsthed hal\-%
% --- p. 36 ---
\opage{36}\setcounter{footnote}{0}%
veres i Tænken, faaer den en halv Tanke istedenfor en heel:
det, der mangler, er Tanke; halveres den i Villie, faaer den
en halv Villie istedenfor en heel; det, der mangler, er Villie.
Derimod kan man ikke i strengere Forstand sige, at en Be\-%
vidsthed halveres af Tænken og Villen saaledes, at Tan\-%
kens Halvdeel bliver Villie, Villiens Halvdeel Tanke; thi
Tanke kan kun i høist uegentlig Mening forvandles til
Villie, og Villie til Tanke. Staaer det da fast, at Troen
paa ingen Maade kan forvandles til Viden eller Viden til
Tro, saa staaer det ogsaa fast, at Troen ligesaa umulig kan
halveres ved Viden som Viden ved Tro. Kan en Bevidst\-%
hed ikke halveres af Venskab og Mathematik, af Kjærlighed
og Chemie, saa kan den endnu mindre halveres af Ethik og
Astronomie, af Religion og Geologie. Der gives mathematiske
Venner, men der gives ingen venskabelig Mathematik; der
gives forelskede Chemikere, men ingen chemisk Forelskelse, ingen
forelsket Chemie. Der kan gives ethiske og uethiske Astronomer,
men Astronomien er hverken ethisk eller uethisk. Videnskaben
er med andre Ord hverken ethisk eller uethisk, hverken christe\-%
lig eller uchristelig, hverken troende eller vantro: Kjærlig\-%
hed er praktisk, men Videnskab er theoretisk. Det Busse\-%
mandsagtige i Biskop Martensens Pathos (S. 79): „Du
skal kun praktisk elske og erkjende mig, men theoretisk skal
Du“ --- hvad der er umuligt --- „fornægte mig! Praktisk
skal Du holde Dig til den Tro, som er overantvortet min
Menighed; men theoretisk skal Du“ --- i Kraft af en theo\-%
logisk Indbildning --- „holde Dig til Atheister og Antichrister,
skal vandre Fornægteres Vei og sidde i Spotteres Sæde,“ ---
falder altsaa bort af sig selv.

Anderledes, ganske anderledes, bliver Naturlæren og med
den Naturforskeren at behandle i Universalmonarchiet. Her
nytter det ikke, at Naturforskeren vil have sin Religion i en
praktisk, sin Videnskab i en theoretisk, af Religionen uafhængig,
Erkjendelse; thi Biskop Martensen fordrer, at den theoretiske
% --- p. 37 ---
\opage{37}\setcounter{footnote}{0}%
Erkjendelse skal gaae i Eet med den praktiske, den praktiske
med den theoretiske. I Monarchiet er det ikke nok, at Astro\-%
nomen selv er troende, nei, han skal være troende i Conse\-%
qvens af en troende Astronomie; ikke nok, at Geologen per\-%
sonlig vil være Christen, nei, han skal være Christen i det
Praktiske paa Grund af en christelig Geologie i det Theo\-%
retiske. Thi, siger Biskop Martensen (S. 133): „Det nytter
slet ikke, at man, naar man kommer i Forlegenhed, forsikkrer,
at det Praktiske, at Livet alligevel har Primatet fremfor
Videnskaben. Thi om saadanne Forsikkringer maa der siges,
at de have Livets Skin, men fornægte dets Kraft, saalænge
som en hedensk, opblæst Viden istedetfor at kastes af, skal
conserveres, udvikles og opelskes, pleies, fredes og næres ved
Siden af Troen i een og samme Bevidsthed“. I Monar\-%
chiet er det ikke nok, at en samvittighedsfuld Naturforsker vil
staae Religion og Ethik til Regnskab for sit personlige For\-%
hold til Videnskabens upersonlige Sandheder, at han ligesaa
lidt vil indføre Vantro som Tro, ligesaa lidt Hedenskab
som Christendom i sin videnskabelige Viden, at den Villie
altsaa, hvormed han ærligt og redeligt hengiver sig til Viden\-%
skabens Dyrkelse, er ethisk bestemt ved den Villie, hvormed
han personlig nærer og udvikler Troen i Livet. Nei, der
skal mere til i Monarchiet. Den Villie, hvormed en Na\-%
turforsker dyrker sin Videnskab, skal være en saa troende
Villie, at den kan gjøre Videnskaben selv troende. Kan hans
Villie ikke bringe det dertil, at den gjør hans Videnskab
troende: da faaer han jo to Villier, en i Videnskaben og en
anden i Troen; han vil da „ikke Eet men To, og har et
dobbelt Formaal“. Men vee den Naturforsker, der i Mon\-%
archiet „ikke vil Eet, men To, og har et dobbelt Formaal“!
Det vil gaae ham med sin Naturlære ligesaa ulykkeligt, som
det er gaaet os med „vor nyeste Philosophie“; og det er jo
nu desværre en bekjendt Sag, at det er gaaet os meget
ulykkeligt. Thi, siger Biskop Martensen (S. 139): „Som
% --- p. 38 ---
\opage{38}\setcounter{footnote}{0}%
vidende vil den“ (vor nyeste Philosophie) „Menneskets Eman\-% "vidende" is not letterspaced, though its counterpart "troende" below is
cipation fra Christendommen, ja fra den levende Gud. Den
har da sit Hjem i Vantroens Leir, hvorfra den forkynder, at
ligeoverfor den moderne Videnskab og den paa Naturviden\-%
skaberne byggede Culturbevidsthed er Christendommen bleven
en afmægtig og uholdbar Sag. Som \emph{troende} vil den
Lydigheds- og Autoritetsforholdet og hævde Troens fuld\-%
komne Uafhængighed af Videnskaben, dog kun som et absolut
Absurdum, --- en Religiøsitet, der maa betegnes som Obscu\-%
rantisme, hvilken dog deri er forskjellig fra anden Obscu\-%
rantisme, at den fra Først til Sidst er \emph{theoretiseret}.
Hvorledes de to Villier gaae sammen til een, er ikke be\-%
gribeliggjort og vil ifølge Sagens Natur aldrig blive det“.

At vi ved denne Leilighed presse os med vor Philoso\-%
phie saa nær ind paa Naturlæren, skeer ingenlunde for
Philosophiens Skyld, thi Philosophien maa forsvare sig selv;
nei det skeer kun for vor egen Skyld. Det er jo naturligt,
at vi, der i vor Stilling til Monarchiet ere saa ulykkelige,
maae nære en inderlig Sympathie for andre, høistærede Med\-%
ulykkelige. Det er af Sympathie og Bekymring, at vi til Slut\-%
ning opkaste det Spørgsmaal: Hvad skal der blive af de
ulykkelige Naturforskere? Vi tale her ikke om alle Natur\-%
forskere, men kun om de ulykkelige. Hvilke ere vel de ulyk\-%
kelige Naturforskere? Ulykkelige kalde vi de Naturforskere,
der ikke blot ville dyrke deres Videnskab, men ogsaa deres
Gud. Hvorfor ere de ulykkelige? Fordi de ligerviis som
„vor nyeste Philosophie“ ikke ville Eet, men To, og altsaa
have to Villier. Dette kan „det kongelige Princip“ paa ingen
Maade tillade. At Christus i „Dogmatiken“ faaer Lov til
at have to Villier, er en anden Sag; det kan Principet for\-%
svare, thi det er begribeligt. „Modsætningen mellem det
naturlige og det aandelige, det kosmiske og det hellige Prin\-%
cip“ maa jo „nødvendigt findes i den anden Adam som en
Dobbelthed i hans Villie; da den anden Adam ikke kan
% --- p. 39 ---
\opage{39}\setcounter{footnote}{0}%
tænkes monotheletisk (hvilket vilde være at tænke ham mono\-%
physitisk), men maa tænkes dyotheletisk: saa rører sig jo ogsaa
det Princip i den anden Adam, der gjør den førstes Synde\-%
fald \emph{muligt}“. Anderledes med en ulykkelig Naturforsker.
Vistnok rører sig ogsaa i den naturforskende Adam „det
Princip, der gjorde den første Adams Synd mulig“; men
desuagtet maa Dogmatiken tænke sig enhver Naturforsker,
endog den ulykkelige, ikke dyotheletisk, men monotheletisk,
efterdi den absolut maa tænke ham, ikke dyophysitisk, men
monophysitisk. Maa Dogmatiken nu, ifølge en høiere viden\-%
skabelig Nødvendighed, tænke Naturforskeren monotheletisk, saa
maa Naturforskeren ogsaa være monotheletisk. Men, spørge
vi, hvorledes kan en Naturforsker med to Villier være mo\-%
notheletisk? Saasnart en Naturforsker gjør Mine til at
ville have to Villier og opføre sig dyotheletisk, maa han øie\-%
blikkelig Eet af To: enten ud af Menigheden eller ud af
Videnskaben. Een være Hersker, den gode Villie! Men to gode
Villier kunne umulig rummes i een Naturforsker; er Villien
til at anerkjende Naturlovene god, saa er Villien til at an\-%
erkjende Christi Lov ikke god; er Villien til at anerkjende
Christi Lov god, saa er Villien til at studere den „forbandede“
Natur, ikke god. Endnu engang: enten ud af Menigheden
eller ud af Videnskaben! Her er ingen Barmhjertighed. En
Naturforsker med to Villier kan Biskop Martensen ikke taale
i sit Monarchie, altsaa heller ikke i sin Menighed; thi Me\-%
nighed og Monarchie og Monarchie og Menighed ville ikke
To, men Eet. Skulde Biskop Martensen taale en dyotheletisk
Naturforsker i Menigheden, da maatte det være begribelig\-%
gjort, hvorledes to ueensartede Villier, en theoretisk og en
praktisk, gik sammen til een; men, siger Biskop Martensen,
„hvorledes de to Villier gaae sammen til een, er ikke be\-%
gribeliggjort og vil ifølge Sagens Natur aldrig blive det“.

% --- p. 40 ---
\opage{40}\setcounter{footnote}{0}%
\section*{II.\\ Historiens Vilkaar i Monarchiet.}
\addcontentsline{toc}{section}{II. Historiens Vilkaar i Monarchiet}
\markboth{Historiens Vilkaar i Monarchiet}{Historiens Vilkaar i Monarchiet}

Den dybe Kloft imellem Naturforskerens og Dogmati\-%
kerens Gjenstand, Methode og Formaal har til alle Tider
været saa iøinefaldende, at vi allerede ved Begyndelsen af
første Afdeling nok anede, hvordan det vilde gaae Naturlæren,
naar det theologiske Universalmonarchie blev oprettet. Med
Historien er det nu en anden Sag, hvis vi ellers tør fæste
Lid til „den christelige Dogmatiks“ meget anerkjendende, for
Historien særdeles hæderlige Vidnesbyrd. „Thi“, siger Bi\-%
skop Martensen (Dogmat. § 12), „skjøndt det vistnok er
Skaberaanden, der igjennem Naturen taler til den skabte
Aand, saa taler Naturen med sit stumme Sprog dog kun
paa en indirecte og forblommet Maade om Skaberens evige
Magt og Guddommelighed; den directe, utvetydige Aaben\-%
baring kan kun findes i Aandens egen Verden, i Ordets,
i Samvittighedens, i Frihedens Verden, eller i \emph{Historien}.
Aabenbaring og Historie ere derfor uadskillelige, men dersom
der ikke gaves nogen anden Historie end \emph{Verdenshistorien},
da vilde Guds Aabenbaring dog endnu savne sit adæqvate
Medium.... Vel lyder Guds hellige Røst igjennem de
verdenshistoriske Røster, vel er Guds Gjerning i de verdens\-%
historiske Begivenheder; men i den verdenshistoriske Tummel
sammenblandes Guds Røst og Menneskenes Røster for vort
Øre, og den hellige Forsynsvillie, hvilken vi øine gjennem
de menneskelige Skjæbner, skjuler sig atter for vort Blik under
Verdensløbets rastløse Strøm. Skal der i Sandhed være
Tale om en hellig Gudsaabenbaring, da maa der være en
Historie i Historien, da maa der indenfor Verdenshistorien
være en \emph{hellig Historie}, i hvilken Gud aabenbarer sig
som Gud, en Historie, i hvilken \emph{det hellige Verdens\-%
formaal} aabenbarer sig som saadant, hvor Guds Ord saa\-%
ledes indfatter sig i det menneskelige Ord, at det sidste bliver
% --- p. 41 ---
\opage{41}\setcounter{footnote}{0}%
det rene Organ for det første, og hvor Guds Gjerning saa\-%
ledes indfatter sig i Menneskets Gjerning, at den sidste bliver
fuldkommen gjennemsigtig for den første“.

Dette er nu baade træffende og smukt; tiltalende især,
saalænge man seer de historiske Kjendsgjerninger paa et Par
Miles speculative Afstand\footnote{I min Licentiatdisputats
\textit{De speculativa historiæ sacræ tractandæ methodo}, oversat af
B. C. Bøggild 1842, vil man f. Ex. § 10 og § 11 finde lignende
Grundtanker udtalte.}. Men naar Kjendsgjerningerne
i deres Enkeltheder og Mangfoldighed rykkes Betragteren
nærmere; naar Forholdet imellem den historiske (\textit{fides histo\-%
rica}) og den religiøse Tro (\textit{fides religiosa}) nøiagtig skal be\-%
stemmes; naar Kritikens Ret i dens Væsen og Omfang
skal erkjendes: da komme Vanskelighederne, og Vanskelig\-%
hederne ere mange. Uden videnskabelig Kritik er Hi\-%
storien ingen Videnskab; men uden en kritisk Maalestok
gives der ingen videnskabelig Kritik. Hvad er det nu, der
sikkrer Kritiken en Maalestok? Det Logisk-Physiske! Som
Videnskab er da Historien i dette Punkt ligesaa autonomisk
som Naturlæren. Historien kan, naar den sees fra denne
Side, ligesaa lidt opgive det Logisk-Physiske som Naturlæren;
den kan ligesaa lidt som Naturlæren lade sig regjere af „den
gode Villie“.

Naar Monarchiet forlanger af Historikeren, at han i
sidste Instans skal lade Alt komme an paa „den gode Villie“,
da er Meningen ingenlunde den, at han skal ligefrem gjøre
Brud paa det Logisk-Physiske, handle vilkaarligt med Kjends\-%
gjerninger og lade Øiemedet i jesuitisk Forstand hellige Midlet.
Nei, Vanskeligheden er af subtilere Art. Den gode Villie
kan nemlig ikke adskilles fra det Godes Idee; men det Godes
Idee maa ifølge det kongelige Princip absolut være Eet med
det Sandes Idee: det Gode er sandt, det Sande er godt.
% --- p. 42 ---
\opage{42}\setcounter{footnote}{0}%
Thi, siger Biskop Martensen: „Først naar man fra alle sub\-%
ordinerede Ideer er stegen op til det Gode, som til det kon\-%
gelige Princip ... kan der være Tale om den levende Mod\-%
sætning mellem Sandhed og Løgn, der er noget heelt Andet
end den blot logiske Modsætning mellem Rigtigt og Urigtigt,
Væsen og Skin“. Her er Knudepunktet. Det er Forholdet
mellem Idee og Virkelighed, Idee og Kjendsgjerning, hvor\-%
paa Eftertrykket falder. Bestemmes Ideen ved Kjendsgjer\-%
ningen, saa hersker Virkeligheden; men i Virkeligheden naaer
man ikke ud over Modsætningen mellem Rigtigt og Urigtigt,
Væsen og Skin. Bestemmes Kjendsgjerningen ved Ideen,
saa sigtes der directe paa Modsætningen, den levende Mod\-%
sætning mellem Sandhed og Løgn; men hvordan gaaer det
saa med den blot logiske Modsætning mellem Rigtigt og
Urigtigt, Væsen og Skin? Jo mere Vægt her lægges paa
Virkeligheden, desto mere trænges Ideen tilbage; jo mere
Vægt her lægges paa Ideen, desto mere trænges Virkelig\-%
heden tilbage. Men heraf følger igjen, at den kritiske Maale\-%
stok maa forandre sig, eftersom den anvendes paa \emph{Verdens\-%
historien}, paa \emph{Kirkehistorien} eller paa \emph{den hellige
Historie}. I Verdenshistorien har Virkeligheden med dens
Kjendsgjerninger afgjørende Overvægt over Ideen: først
Kjendsgjerningerne, siden kunne vi tale om Ideen. I den
hellige Historie har Ideen som det Helliges, det absolute
Godes Idee ubetinget Overvægt over de ydre endelige Kjends\-%
gjerninger: først Ideen, saa kunne vi tale om det historisk
Virkelige. I Kirkehistorien vakler Kritiken imellem Ideens
og Kjendsgjerningens Forret.

Hvilke Vanskeligheder, der kunne optaarne sig for den
historiske Kritik, naar Monarchiets Princip skal fyldestgjøres,
er nu let at indsee. Vi forudskikke et Exempel, hentet fra
Kirkehistorien, nemlig Kritiken over de Pseudoisidorske De\-%
cretaler. Kom det her alene an paa at vurdere de historiske
Kjendsgjerninger, da kunde Kritiken i denne Sag bringe dem
% [KB copy: pencil manuscript note in the left margin of p. 42, not transcribed]
% --- p. 43 ---
\opage{43}\setcounter{footnote}{0}%
til at tale saa tydeligt, at her end ikke blev Skygge af
Tvivl tilbage. I den Isidorske Samling findes jo Breve, i
hvilke der fra det tredie Aarhundrede citeres en Bibelover\-%
sættelse, der først bliver til efter det fjerde Aarhundrede,
Breve, i hvilke den romerske Biskop Victor fra det andet
Aarhundrede skriver angaaende Paaskefesten til den alexandrinske
Biskop Theophilus i det fjerde Aarhundrede; her udkræves da i
Sandhed en jesuitisk Opløftelse over det Logisk-Physiske, naar
en Historiker desuagtet skal forsøge paa at forsvare Samlin\-%
gens Ægthed. Gjentagne Gange havde allerede katholske
Forskere (en Petrus Comestor i det tolvte, en Marsilius i
det fjortende, en Nicolaus af Cusa i det femtende Aarhun\-%
drede) yttret sig imod Ægtheden, da de magdeburgske Cen\-%
turier omsider leverede et udførligt, fra flere Sider stadfæstet
Beviis for Uægtheden. Ikke desto mindre vovede Jesuiten
Franz Turrianus at træde i Skrankerne for Pseudoisidor.
Jesuiten og hans Client bleve i Begyndelsen af det syttende
Aarhundrede lyste hjem med saadan
% Rettelser: p. 43 l. 17 fra neden "saadan" læs "saadant" -- not applied.
% The KB copy has a "t" added in ink after the word; transcribed as printed.
Eftertryk af Blondell
(\textit{Dav. Blondelli Pseudo-Isidorus et Turrianus vapu\-%
lantes. Genev. 1628}), at selv troende Papister meente, at
nu kunde det være nok. Men nei, endnu i Slutningen af det
attende Aarhundrede kunde en Marchetti (\textit{Saggio critico sopra
la storia di C. Fleury}) ikke overbevise sig om Decretalernes
Uægthed, og der vil vist findes Mange endnu, der ikke kunne
overbevise sig derom. Denne Haardnakkenhed er mere end simpel
Trods; den er et Udtryk for, at Ideen og „den gode Villie“
ere med i Spillet.

Selv en Skabning af Pavekirkens Idee, er Pseudoisidor
en Ideens Bærer. Sættes nu Ideen om Pavens Primat
identisk med det for Kirken høieste Gode og dette høieste
Gode som Historiens „kongelige Princip“, da er Jesuiternes
Udholdenhed i at forsvare Pseudoisidor forklarlig. At
Kjendsgjerningerne gaae deres Bestræbelser imod, kan jo
være dem ubehageligt nok; men derfor tabe de ikke Modet.
% --- p. 44 ---
\opage{44}\setcounter{footnote}{0}%
Gjennem store Vanskeligheder hjælpe de sig ved Distinctioner.
Enhver, der vil kritisere, maa kunne distingvere, og Jesuiter
kunne distingvere. Men blandt alle til historisk Kritik hørende
Distinctioner vil man dog neppe kunne udfinde nogen frugt\-%
barere end den, hvorved Biskop Martensen, tilvisse uden
Bitanke om mulige Misbrug, har skilt „den blot logiske Mod\-%
sætning mellem Rigtigt og Urigtigt fra den levende Mod\-%
sætning mellem Sandhed og Løgn“. Idet den logiske Mod\-%
sætning mellem Rigtigt og Urigtigt nedsættes til et Under\-%
ordnet, bliver den historiske Kritik som rørt af et Slag.
Hvad ændser en af „den gode Villie“ bevæget Historiker vel
nogle Modsigelser imellem enkelte Kjendsgjerninger, naar det
gjælder om at forsvare det Godes Idee? Modsigelser af
denne Art kunne i al Fald kun lede til en Forskjel mellem
det Rigtige og det Urigtige; men at forsvare Ideen er at for\-%
svare Sandheden. Naar vi see, hvorledes Biskop Martensen
(Leilighedsskr. S. 80) kan ivre for at faae Christus aner\-%
kjendt for Konge over alle Videnskaber, og først i dette theo\-%
retiske Kongedømme seer det Godes Idee fyldestgjort, hvor
kan det da undre os, at see Curialisterne ivre for at for\-%
svare et Primat, der ifølge deres „gode Villie“ nødvendig
tilkommer Christi Statholder?

Allerede ved Opfattelsen og Bestemmelsen af en kirke\-%
historisk Idee, viser det sig, hvilken Magt den religiøse Tro
er istand til at udøve over den historiske; men denne Magt
er endnu langt mere afgjørende, naar Talen er om \emph{den
hellige Historie}. I den hellige Historie maa Grundideen
oprindelig bestemmes ved den religiøse, ene og alene ved den
religiøse Tro. Derfor er Dogmatiken, som vi have seet,
saa forvisset om, at den hellige Historie absolut maa være
en Historie, „hvori Guds Ord saaledes indfatter sig i det
menneskelige Ord, at det sidste bliver det rene Organ for
det første, og hvor Guds Gjerning saaledes indfatter sig i
Menneskets Gjerning, at den sidste bliver fuldkommen gjen\-%
% --- p. 45 ---
\opage{45}\setcounter{footnote}{0}%
nemsigtig for den første.“ I Overeensstemmelse hermed
har Biskop Martensen i sine „dogmatiske Oplysninger“
(S. 55) tydelig udtalt sig om Troen som Erkjendelsens første
Betingelse og om Dogmatikerens Forhold til de hellige
Kjendsgjerninger. „Jeg tilegner mig“, siger han, „de Ord,
hvormed Anselmus indleder sit Proslogium, og som i
mine Tanker kunne tjene som Indskrift for enhver Dog\-%
matik: „Ikke vover jeg, o Herre, at udgrunde dine Dyb\-%
heder; thi ingenlunde formaaer min Erkjendelse dette; kun
begjerer jeg i nogen Maade at erkjende din Sandhed, som
mit Hjerte troer og elsker. Thi ikke for at troe søger jeg
at erkjende, men jeg troer for at erkjende. Thi ogsaa dette
troer jeg, at dersom jeg ikke troede, da vilde jeg heller ikke
erkjende.“ --- Idet jeg ikke finder mig foranlediget til at af\-%
vige fra det her anviste og i min Dogmatik fulgte Spor,
vedbliver jeg fremdeles at betragte Dogmatikerens Forhold
til Aabenbaringen paa tilsvarende Maade som Naturforskerens
Forhold til Naturen. Den ægte Naturforsker opgiver ikke
Naturen, fordi han ikke seer sig istand til at forklare ethvert
af dens Phænomener, men opgiver dog heller ikke at trænge
ind i Naturens Indre, saavidt han formaaer det.“

Det fortjener den allerstørste Opmærksomhed, at Biskop
Martensen fremdeles vedbliver at betragte Dogmatikerens For\-%
hold til Aabenbaringen som Naturforskerens Forhold til Na\-%
turen; thi heraf sees, hvor sikker han er paa de hellige
Kjendsgjerninger. Hvorved bestemmes Naturforskerens For\-%
hold til Naturen? Ved uforanderlige Love, ved uafviselige
Kjendsgjerninger. „Den ægte Naturforsker opgiver ikke
Naturen, fordi han ikke seer sig istand til at forklare ethvert
af dens Phænomener.“ Rigtig! Men, hvorfor opgiver Natur\-%
forskeren ikke Naturen? Fordi han ikke kan opgive den,
fordi den i umiddelbar Anskuelighed, i sandselig Nærværelse
afnøder ham Anerkjendelse. Skal nu Aabenbaringen med
sine Kjendsgjerninger og Love paa tilsvarende Maade afnøde
% --- p. 46 ---
\opage{46}\setcounter{footnote}{0}%
Dogmatikeren Anerkjendelse, da maa det jo forudsættes, at
Aabenbaringens Kjendsgjerninger i deres Kreds ere ligesaa
sikkre, ligesaa anskuelige og uafviselige som Naturens, Aaben\-%
baringslovene i deres Sphære ligesaa uforanderlig bestemte
som Naturlovene. Fra denne i sig selv respectable Forud\-%
sætning er den gamle orthodoxe Dogmatik ganske rigtig gaaet
ud; men hvordan er det gaaet med Forudsætningen? Den
historisk-philologiske Kritik har behandlet den haardt, ja meget
haardt.

Skal, hvad den protestantisk theologiske Orthodoxie op\-%
rindelig forudsætter, den Hellige Skrift gjælde for et fuld\-%
komment paalideligt, guddommeligt Aabenbaringsdokument:
da maa den antages at være hævet over al videnskabelig
Kritik. Det var da umuligt at tænke sig, at der i den Hel\-%
lige Skrift kunde findes enten dogmatiske, moralske eller
historiske, chronologiske, topographiske, onomastiske Feil.\footnote{\textit{Omnia
et singula sunt verissima, quæcunque in illa (script.\ sacr.) traduntur:
sive dogmatica illa sint sive moralia, sive Historica, Chronologica,
Topographica, Onomastica, nullaque ignorantia, incogitantia aut
oblivio, nullus memoriæ lapsus Spiritus Sanctus
% Rettelser: p. 46 Anm. l. 7 fra neden "Sanctus" læs "Sancti" -- not
% applied. The KB copy has an ink correction written over the ending of
% the word; the PRINTED "Sanctus" is transcribed.
amanuensibus in consignandis sacris litteris tribui potest aut debet.
Quenst. p. I. pg. 77.} Θέσις.}
Endnu imod Slutningen af det attende Aarhundrede paastode
rettroende Theologer ikke blot, at de hellige Skribenter aldrig
have feilet, men endog, at de ikke have kunnet feile.\footnote{„Ihre
(der heiligen Schrift) Verfasser haben nicht allein nie geirrt, sondern
auch nicht einmal irren können, und ihre Erzählungen sind daher nicht
allein so gewiss, als jede andere historische Wahrheit, sondern auch
untrüglich gewiss.“ Chr. W. F. Walch. Krit. Nachricht. v. d. Qvellen
der Kirchengeschichte. Leipzig 1770. S. 70.} Paa
en saa klippefast Grund kunde den hellige Historie modstaae
Kritikens Angreb. Hellige Varsler, store Vidundere omgive
Christi Komme til Verden, hans Undfangelse, Fødsel og
Barndom; men Ingen kan negte disse Varsler, disse Vid\-%
% --- p. 47 ---
\opage{47}\setcounter{footnote}{0}%
undere, thi det er jo Kjendsgjerninger. Mirakel paa Mirakel
møder os i Jesu offentlige Liv; det er jo Kjendsgjerninger. Hvo
skulde ikke ville troe paa Christi Opstandelse, naar Kjendsgjerningen
er saa historisk, at Ingen kan negte den uden at negte al Historie?
Hvad gjør vel det, at Himmelfarten strider imod Tyngdens
Love, naar Kjendsgjerningen selv er ligesaa vis, som de Kjends\-%
gjerninger, hvoraf man har abstraheret Tyngdens Love? Tør
nu „den ægte Naturforsker“ ikke negte en naturlig Kjends\-%
gjerning, fordi han ikke strax kan forklare den, da tør den
ægte Bibelforsker endnu mindre negte en overnaturlig Kjends\-%
gjerning, fordi han ikke kan begribe den. Fra dette Stand\-%
punkt af er Biskop Martensen conseqvent, naar han „frem\-%
deles vedbliver at betragte Dogmatikerens Forhold til Aaben\-%
baringen paa tilsvarende Maade som Naturforskerens Forhold
til Naturen“.

Men maae vi nu ikke paa Biskop Martensens Vegne
bitterlig beklage, at Standpunktet for henved hundrede Aar
siden er forsvundet, ligesom vi jo alle maae beklage, at det spe\-%
culative Paradiis er forsvundet. Og hvorledes er det forsvundet?
Det er igjen „Historien om Syndefaldet og \emph{de tvende
absolut ueensartede Træer}, af hvilke det enes Frugter
% Printer's defect: the „ opening "Historien om Syndefaldet" is never
% closed; the only closing mark is the one after "forbudne" below.
% Transcribed as printed (quote balance +1).
vare „Theologien forbudne“.\footnote{At Kritikens Frugter virkelig vare
Theologien forbudne, kan blandt Andet sees af et Vers, som Georg Pasor
lægger den fra de puristiske Stridigheder bekjendte Sebastian Pfocken i
Munden:\\
\hspace*{2em}\textit{Diligo philologum, qui eviscerat abdita verba}\\
\hspace*{2em}\textit{Jure novum Christi sed mage fedus amo}\\
\hspace*{2em}\textit{Ergo novo quisquis de federe tenuia decit,}\\
\hspace*{2em}\textit{En pupillam oculi vulnerat ille mei.}\\
\hspace*{2em}--- --- --- --- --- --- --- ---\\
\hspace*{2em}\textit{Falx critica in quosvis per me grassetur. At unis}\\
\hspace*{2em}\textit{A Bibliis fas est abstinuisse manus.}} Kritiken, den historisk viden\-%
skabelige Kritik, er her forbilledligt betegnet ved Kundskabens
% --- p. 48 ---
\opage{48}\setcounter{footnote}{0}%
Træ; Slangen i Træet er den sunde Menneskeforstand, der
altid har viist en syndig Forkjærlighed for det Logisk-Physiske.
Ligerviis som Slangen i Paradiis, hiin første Slange, det
dogmatisk kosmiske Princip, der er ældre end Diævelen, efter\-%
som Diævelen, ifølge vor „christelige Dogmatik“, er en Søn
af det kosmiske Princip, forførte Eva med trædske Ord, der
kunde tydes til den Underviisning: „Det er ikke sandt, hvad
Gud har sagt (nemlig fra Videnskabens Standpunkt)“: saa\-%
ledes kom nu den sunde Menneskeforstand, en naturlig Søn
af Diævelen, skjult i spraglet Ham af logisk-physisk Textur,
krøb kritisk op i Kundskabens Træ, stak Hovedet frem
med „Autonomie“ og begyndte at conversere Theologien
med trædske Ord, der kunde tydes til den Underviisning:
Det er ikke sandt, hvad Bibelen har sagt (nemlig fra Kri\-%
tikens Standpunkt). Og Qvinden --- thi i Symbolet er
Theologien en Qvinde --- saae, at Træet, Kritikens Træ,
var godt at æde af, og at det var lysteligt for Øinene og et
ønskeligt Træ til at faae Forstand af, og hun tog af dets
Frugt, og hun aad, og hun gav ogsaa Facultetet Lidt
med, og det aad. Da oplodes begges Øine, og de saae ---
det Nøgne, de saae i Kritiken den nøgne Virkelighed; og de
tyede til Æsthetikens Lund og syede sig Figenblade, og de
skare Strimler af Philosophiens Ryg og gjorde sig Belter.
Og de (Theologie og Facultet) erkjendte, at de ikke uden
Fare for „en egen Plads udenfor Videnskaben“ kunde frem\-%
deles vedblive at betragte Dogmatikerens Forhold til Aaben\-%
baringen paa tilsvarende Maade som Naturforskerens For\-%
hold til Naturen; de erkjendte, at de ikke fremdeles kunde
vedblive at tale høirøstet om Historiens overnaturlige Kjends\-%
gjerninger, efterdi de overnaturlige Kjendsgjerninger nok ikke
vare historiske, og de historiske ikke overnaturlige.

Den i sit Paradiis saa ungdomsfriske, tillidsfulde,
bibelstolte Skrifttheologie, hvor er den efter Syndefaldet ved
det Logisk-Physiske efterhaanden bleven svækket, bøiet, afbleget,
% --- p. 49 ---
\opage{49}\setcounter{footnote}{0}%
ydmyget! Maa det siges, at Diævelen i Speculationen ---
og det dybeste Speculative er altid det høieste Værende --- er
en dogmatisk Søn af det kosmiske Princip og kun paa denne
Maade kan tænkes som en ægte objectiv videnskabelig Diæ\-%
vel: da maa det ogsaa siges, at den med Kritiken leflende
Theologie er en udogmatisk Datter af den faldne Ortho\-%
doxie. Vel sandt, den Faldne mener, at hun endnu har
Videnskab og Samvittighed. Ja Samvittighed! Theologien
maa vel raisonnere om Samvittighed, naar Samvittigheden
sidder med en halv Speculation og sover i en Paragraph
af Dogmatiken. Lad Samvittigheden vaagne, lad Theologien
engang høre Menighedens Stemme!

„\emph{Hvad har Du}“, spørger Menigheden den ved sit
Syndefald syndige Theologie, „\emph{gjort af mit Juleevan\-%
gelium}?“ Ak, kjære Menighed, bliv ikke vred. Kritiken har
desværre revet Juleevangeliet reent istykker. Men tab ikke
Modet! husk paa, at det er Viden, der skal hjælpe paa Troen;
husk paa, hvad Biskop Martensen siger med Anselmus: „Ikke
vover jeg, o Herre, at udgrunde dine Dybheder“; men jeg troer,
at Viden skal hjælpe paa Troen. Græd ikke! Det Evangelium,
Kritiken har revet istykker, kan Apologetiken jo sætte sammen
igjen; thi den synoptiske Apologetik er dog vel videnskabeligt
i sin gode Ret, naar den omhyggelig sanker Stumperne op
og viser, hvorledes de ganske vist maae have siddet.

Du seer saa mørk paa mig. Hvad kan jeg gjøre for,
at Videnskaben er objectiv? Hvem skulde nu ogsaa have
tænkt, at de historiske Kjendsgjerninger i Juleevangeliet vare
saa forunderligt skjøre, de naturlige Bestanddele ligesaa
skjøre som de overnaturlige. Jeg vil nu slet ikke tale om
det Overnaturlige, om Englene, der sang for Hyrderne, om
Stjernen, der gik foran de Vise --- thi om „den mythiserende
Indklædning“, der træder til „udfyldende og formende“ det
historisk Virkelige, tilstaaer endog Apologetiken, at den
hverken veed ud eller ind --- jeg vil, hvad det Naturlige an\-%
% --- p. 50 ---
\opage{50}\setcounter{footnote}{0}%
gaaer, blot i al Fortrolighed lade dig vide, at det er reent
forkeert med Qvirinius.

„\emph{Hvad har Du}“, spørger Menigheden, den ved sit
Syndefald syndige Theologie, „\emph{gjort af mit Paaskeevan\-%
gelium}?“ Ak, kjære Menighed, Fortællingerne om Christi
Død og Opstandelse frembyde virkelig ikke faa Vanskelig\-%
heder for Videnskaben. Vi ville ikke opholde os ved det
Chronologisk-Synoptiske; thi hvad det Chronologisk-Synop\-%
tiske angaaer, kan Menigheden ikke noksom takke den i Op\-%
findsomhed uudtømmelige Apologetik. Men hvo føler ikke
Vægten, der maa tynge paa Videnskaben, naar han tænker
paa Begivenhederne ved Jesu Død. „Og see, Forhænget i
Templet splittedes i to, fra det Øverste indtil det Nederste;
og Jorden skjalv og Klipperne revnede; og Gravene oplodes;
og mange af de Helliges Legemer opstode; og de gik ud af
Gravene efter hans Opstandelse, og kom ind i den hellige
Stad og aabenbaredes for mange.“

For at stille sig honnet ved dette Mirakel er den synop\-%
tiske Apologetik selv nærved at gjøre et Mirakel, et ægte
apologetisk Mirakel. At Du dog kan faae en lille Fore\-%
stilling om, hvorledes den theologiske Exegetik formaaer at
forsone Tro med Viden, meddeler jeg her følgende fra vore
nyeste Lærebøger hentede Fortolkning: „Herrens Opstandelse
var skeet, da Besidderne af Familiegravene toge de Ødelæg\-%
gelser i Øiesyn, som Jordskjælvet havde anrettet i disse
Klippe-Gravkamre. Gjennem Revnerne i Klipperne og
% "Klippe-Gravkamre": the hyphen stands between wide spaces in a
% loosely justified line.
Steendørene, som ved Jordrystelsen vare løsnede og kastede ud,
vare Levninger af deres Afdøde at see, som de hvilede paa
Gulvene eller paa Steenbænkene. Ved Synet af disse Lev\-%
ninger fik Besidderne af Familiegravene visionære Aaben\-%
barelser“; og saa negte man, at her er en Tænkning „hvis
Fremstilling af Miraklet veier op mod de physicalske Ind\-%
vendinger“.

„Hvor er Troen“, spørger Du, „bleven af?“ Troen er
% --- p. 51 ---
\opage{51}\setcounter{footnote}{0}%
jo objectivt i de visionære Aabenbarelser, Besidderne af Fa\-%
miliegravene fik; subjectivt i den Tro, hvormed Fortolkeren
selv troer paa deres Tro, der fik de visionære Aabenbarelser.
„Og Viden?“ Ja Viden er her da nok af. Her er jo,
naar vi ret tage Ødelæggelserne ved denne Opstandelse i
Øiesyn, meget mere Viden i Fortolkeren, end der er i Evan\-%
gelisten. „Men nu Forsoningen af Tro og Viden?“ Kjære
Menighed! Kan Du ikke see Forsoningen af Tro og Viden
i den hermeneutiske Forsoning af Fortolkningens religiøse og
videnskabelige Momenter: da maa det være, fordi Din Troes\-%
erkjendelse kun er praktisk og subjectiv ikke tillige theoretisk
objectiv. Vel er her ikke givet nogen egentlig Forklaring af
Miraklet --- „Theologien giver sig ikke af med Mirakelfor\-%
klaringer“ --- men en videnskabelig Forklaring er her dog
virkelig givet. Husk paa, kjære Menighed, at vi med vore
Forklaringer vandre i Tro! Ja, vi vandre i Tro, ikke i
Beskuelse; men derfor kunne vi jo gjerne vandre videnskabe\-%
ligt, naar vi kun sørge for, at der under denne vor jordiske
Vandring ikke kommer mere ind i vor Praxis, end der er i
vor Theorie. Vi kunne da sige med Biskop Martensen:
„Ikke vover jeg, o Herre, at udgrunde Dine Dybheder; thi
ingenlunde formaaer min Erkjendelse dette“; men dette vover
jeg at paastaae, at jeg meget godt kan komme til en høiere
videnskabelig Erkjendelse af Dine Mirakler uden at kunne for\-%
klare Dine Mirakler. Jeg vilde jo ellers ikke i \emph{nogen
Maade} kunne erkjende Din Sandhed, som mit Hjerte troer
og elsker. Thi mit Hjertes Tro er ikke blot en subjectiv
og praktisk Tro, min Kjærlighed ikke blot en subjectiv og
praktisk Kjærlighed; min Tro er tillige en objectiv og theo\-%
retisk Tro, min Kjærlighed en objectiv og theoretisk Kjær\-%
lighed. Med Hjertet troe vi, og med Hjertet elske vi:
som Hjertet er, saaledes er Troen og Kjærligheden; som
Troen og Kjærligheden er, saaledes er Hjertet. Husk paa,
kjære Menighed, vi have ikke blot annammet et subjectivt og
% --- p. 52 ---
\opage{52}\setcounter{footnote}{0}%
praktisk Hjerte, vi have i Aanden tillige annammet et theo\-%
retisk-objectivt Hjerte. Ja, jeg erkjender, at mit Hjerte er
theoretisk-objectivt; det vil sige: jeg erkjender det i Troen.
Thi ogsaa dette troer jeg, at dersom jeg ikke troede det,
vilde jeg heller ikke erkjende det.

„Hvad har Du“, spørger Menigheden, „med Dit theo\-%
retiske Hjerte dog gjort af den evangeliske Historie, af For\-%
klarelsen paa Bjerget, af Bespiisningen i Ørken, af Van\-%
dringen paa Havet?“ o. s. v. o. s. v. Vi behøve ikke at
fortsætte disse Forhandlinger for at faae et Indtryk af, at
den med „Culturbevidstheden“ enige og dog uenige, meget
vidende og dog lidet erkjendende Theologie har en ond Sam\-%
vittighed, naar den midt i sin Viden skal gjøre Menigheden
Regnskab for sin Tro. Den, der vil have en virkelig Ind\-%
sigt i det ved Theologien tilveiebragte Misforhold mellem
Tro og Viden, maa ikke lade sig bedaare af dogmatisk specu\-%
lative Talemaader, men gjøre sig bekjendt med den fremad\-%
skridende theologiske Fortolkning af det Nye Testament, navn\-%
lig af de fire Evangelier, og lære den Vei at kjende, der
historisk og conseqvent har ført til en mythisk Behandling af
Jesu Liv.
% After "Liv." a faint hairline stroke, not the printed closing-quote sort
% (no „ opens this passage); not transcribed.

Rigtignok siger Biskop Martensen (Leilighedssk. S. 92),
at det kun er „Schellings aflagte Klæder, med hvilke Strauß
og Andre have smykket sig, for i denne Løvehud at angribe
Christendommen;“ men denne Antagelse hører ifølge den
logiske Modsætning imellem det Rigtige og det Urigtige,
aabenbart til det Halvrigtige, for ikke at sige, det i Grunden
Urigtige. Man behøver blot at kaste et Blik paa den
straußiske „Indledning“ til Jesu Liv for at overbevise sig om
Sagens virkelige Sammenhæng. „Efterat allerede Semler
(Indl. § 8) havde talt om et Slags jødisk Mythologie og
ligefrem kaldt Fortællingerne om Samson og Esther Myther;
efterat Eichhorn paa den (§ 6) anførte Maade var gaaet
foran, blev snart Begrebet Mythus ved Gabler, Schelling
% --- p. 53 ---
\opage{53}\setcounter{footnote}{0}%
og Andre med større Ret opstillet som et for al ældre Historie,
hellig saavel som profan, almindeligt Begreb, efter Heynes
Grundsætning: \textit{a mythis omnis priscorum hominum cum
historia, tum philosophia procedit}, og Bauer vovede end\-%
og at træde op med det gamle og nye Testamentes Mytho\-%
logie (1802). .... Ligesom (§ 9) Eichhorn, paa det gamle
Testamentes Omraade, ved Historien om Syndefaldet blev
ved Sagens Magt trængt over fra sin tidligere naturlige
Forklaringsmaade til den mythiske Opfatning, saaledes Usteri
paa det nytestamentlige Omraade med Hensyn til Fristelses\-%
historien. Han havde i et tidligere Arbeide efter Schleier\-%
macher opfattet den som en af Jesus foredragen, men af
Disciplene misforstaaet Parabel, men indsaae dog snart Van\-%
skelighederne ved denne Opfatning, og idet han havde saavel
den supranaturalistiske som den naturlige Anskuelse af For\-%
tællingerne, efter deres forskjellige Afskygninger, endnu langt
længere bag efter sig, blev der intet Andet tilbage for ham
end at betræde det mythiske Standpunkt, hvilket han da og\-%
saa gjorde i en senere Afhandling og det paa en meget
kraftig Maade“; o. s. v. Theologerne komme, som man
heraf seer, til deres mythiske Forklaringsmaade med en Nød\-%
vendighed, der ikke hænger i Schellings „aflagde Klæder.“ Men
Biskop Martensen, der har kastet et attraaende Blik paa Schellings
\emph{nye} Klæder, kommer hver Gang han seer en mythisk Løvehud
uvilkaarlig til at tænke paa de \emph{gamle} Klæder; thi det er
vist: af Schellings nye Klæder kan der ikke dannes nogen
Løvehud; men der kan nok syes en Faarepels, der „funkler
af Nyhed“.

Hvor Misforholdet imellem Forsvaret og det, man vil
forsvare, er saa colossalt som f. Ex. i den apologetiske For\-%
tolkning af de synoptiske Evangelier; hvor „Concordatet“
imellem Evangelietroen og Theologien er saa uheldigt, for
begge saa piinligt, at Theologien i hveranden Linie mod\-%
siger Troen, og Troen igjen, hvor den seer sit Snit, Theo\-%
% --- p. 54 ---
\opage{54}\setcounter{footnote}{0}%
logien; hvor „Concordatet“, i den Grad afføder Splid, at
de contraherende Parter, trods den Forsigtighed, hvormed
den synoptiske Apologet bestandig afstikker Grændsen, uden
Ophør puffes og kives paa Middelveien, og det uagtet „Middelveien
er bred“; hvor et saadant Phænomen, sige vi, bliver
toneangivende og antager Præg af det Normale: der maa
Grundlaget selv skjule over en uhyre Misforstaaelse.

Vi have i en Række af Aar stadig holdt Øie, ikke blot
med. den speculativt jubilerende% printer's defect: stray full stop after "med" (p. 54 l. 9); transcribed as printed
, altid triumpherende Dogmatisme,
men ogsaa med den meget virksomme, sig klogt modererende, Tankerne halverende Middelveislære. Resultatet af
vore Undersøgelser er, at det ingenlunde er den videnskabelige
Kritik, men den christelige Apologetik, der i Culturbevidsthedens
Navn undergraver Christendommen. Til Stadfæstelse
af denne vor Synsmaade vil det for vort nærværende Øiemed
være tilstrækkeligt at udhæve følgende Hovedpunkter.
Evangelietroen og den videnskabelige Kritik ere i den Grad
uenige angaaende de første Forudsætninger for en Bedømmelse
af den hellige Historie, at Uenigheden kun kan hæves
ved at henføre Tro og Viden til absolut ueensartede Principer.
Den evangeliske Historie indeholder en saadan Forening
af naturlige og overnaturlige Kjendsgjerninger, at
Dommen over ethvert enkelt Evangelium nødvendigviis bliver
afhængig af Dommen over Evangelietroens eget oprindelige
Princip. Ifølge Troesprincipet hører Miraklet, det Overnaturlige,
saa absolut med i Aabenbaringshistorien, har for
den troende Bevidsthed en saa uimodsigelig og ubetinget Gyldighed,
at Aabenbaringstroen væsentlig er en Tro paa det Overnaturlige.
Ifølge Vidensprincipet er Miraklet, det Overnaturlige,
absolut uantageligt: Videnskaben veed \textit{a priori}, at
enhver overnaturlig Begivenhed er mythisk. Det er, naar
vi see nærmere til, ingenlunde Noget, hvorom Evangelisterne
ikke stemme overeens; det er tvertimod det, hvorom de alle
stemme overeens, der fra Grunden af er Kritiken til For\-%
% --- p. 55 ---
\opage{55}\setcounter{footnote}{0}%
argelse, nemlig det Overnaturlige. Hvad der for Alvor
skaber Vanskeligheder er ingenlunde den Omstændighed, at
f. Ex. den ene Evangelist har to Engle ved Opstandelsens
Grav, medens den anden kun har een; nei, det er meget
mere den Omstændighed, at her er en Opstandelse, at her
tales om Engle, at Opstandelse og Engle henføres under det
historisk Virkelige, istedenfor at henføres under Mythologien.

Hvor nu den videnskabelige Kritik faaer Lov til at udvikle
sig frit og conseqvent, som hos Strauß og Renan, der
er strax Resultatet klart. Om der end blandt alle de evangeliske
Beretninger ikke kunde opdages een eneste Uovereensstemmelse,
een eneste historisk Unøiagtighed, saa vilde Viden
dog ligefuldt drive Troen ud; thi Miraklet er, videnskabeligt
seet, en Mythe. Anderledes med Concordatvæsenet i den
christelige Apologetik: her kunne Tro og Viden ligge og
narres med hinanden, saalænge det skal være. Trykket af
„Concordatet“ tør Troen ikke tage Mod til sig og sige det
reent ud, at det Overnaturlige er ligesaa uundværligt for
Troesprincipet som det Naturlige for Vidensprincipet; generet
af „Concordatet“ vil Kritiken fra sin Side heller ikke ud med
Sproget. Idet der paa denne Maade er lagt Dølgsmaal paa
Uenighedens egentlige og væsentlige Grund, saa er Misforstaaelsen
en ulmende Glød --- saa kunne Kjævlerierne om
det Enkelte, de enkelte Beretninger aldrig faae Ende: nu
lader Kritiken, som om den gjerne vilde troe Evangeliet om
Englenes Lovsang, naar der bare ikke havde været den historiske
Bommert med Qvirinius; saa tager Apologetiken til
Gjenmæle og beviser, at Unøiagtigheden meget vel kan undskyldes,
og at man, uden at tage Engle-Aabenbarelserne bogstavelig,
meget vel kan antage „et mirakuløst Grundfaktum“;
nu lader Kritiken, som om den gjerne vilde indrømme det
store Mirakel med „Sønderrivelsen af Tempelforhænget“,
naar den bare kunde „troe, at noget Hensyn skulde have
holdt Apostlene tilbage fra at omtale en saa betydningsfuld
% --- p. 56 ---
\opage{56}\setcounter{footnote}{0}%
Tildragelse“; saa har Apologetiken til Gjengjæld den barnagtige
Tilfredsstillelse at kunne forsvare den omtvistede Tildragelse % a stray mark stands over the final e of "omtvistede" (p. 56 l. 2); not an accent
ved den Bemærkning, at Miraklet med Forhænget
jo ikke var en Tildragelse af „den Betydning, at Tausheden
derom skulde kunne gjælde som Beviis imod det historiske
Faktum“; o. s. v. o. s. v. Skal dette ligesaa videnskabelige
som religiøse Uvæsen standses, da maae Troe og Viden enes
om at forkaste en Apologetik, der er lige utaalelig for begge.
Skal saa Evangelietroen bevares, uden at Kritiken opgiver
sin Ret, da kan det kun skee paa den Betingelse, at Tro og
Viden adskilles som to absolut ueensartede Principer.

Idet Troen faaer Lov at udvikle sine Erkjendelser af
sit eget Princip, bliver Troesprincipet selv kritisk; der danner
sig, hvad Evangeliet angaaer, en Troeskritik i skarp Modsætning
til Videnskritiken. Hvor skarp Modsætningen er, kan
sees deraf, at det, der ifølge Videnskritiken er væsentligt, er
ifølge Troeskritiken uvæsentligt, og omvendt. For Videnskritiken
er det f. Ex. væsentligt at udfinde de historiske Forudsætninger
og Betingelser, under hvilke Forestillingerne om
det Overnaturlige, det Uhistoriske, maae antages at have
dannet sig; til den Ende stræber Videnskritiken med utrættelig
Flid at efterforske Beretningens Oprindelse, dens første
Elementer, dens sagnmæssige Fremvæxt, kort, dens Authentie.
Anderledes med Troeskritiken. Den gaaer principielt ud
fra en ubegribelig Forening af Overnaturligt med Naturligt,
og kan paa ingen Maade tillade noget Forsøg paa at sønderlemme
Evangeliet ved en Udskillen af de ueensartede Elementer;
derimod seer den skarpt efter, om den foreliggende
Beretning har alle til Aabenbaringens Væsen svarende Kjendemærker;
thi den antager ikke Grundmiraklet paa Grund af
den fortalte Begivenhed, men den tilegner sig den fortalte
Begivenhed, forsaavidt denne i alle enkelte individualiserende
Træk anskueliggjør Grundmiraklet. Saavel om Videns- som
om Troeskritiken gjælder det, at enhver af dem efter sit
% --- p. 57 ---
\opage{57}\setcounter{footnote}{0}%
Princip maa arbeide paa at tilveiebringe en Forsoning mellem
Indhold og Form, mellem Sagen og dens Indklædning;
men medens Viden med sin Kritik conseqvent kommer til
Ro i Mythen, kommer Troen med sin Kritik til Ro i
Aabenbaringsordet, det Livets Ord, der lever i Menigheden,
det Lysets Ord, der i Aandens Kraft udgaaer fra Evangeliet
og oplyser Evangeliet.

I Følelsen af, at det nok alligevel ikke hænger rigtig
sammen med den videnskabelige Forsknings religiøse Udbytte,
er „den christelige Apologetik“ da ogsaa ved enhver
passende Anledning betænkt paa en æsthetisk Opfattelse og
Tilegnelse af de hellige Kjendsgjerninger. „For den religiøs-æsthetiske
Betydning af disse Fortællinger“ --- hedder det
om Johannes Døberens og Jesu Fødsels- og Barndoms-Historie% the hyphen falls at the line end (Barndoms=/Historie) and the next word is capitalised: kept as a hyphenated compound (cf. p. 61 "Barndomshistorie")
\footnote{\textit{Dr.} H. N. Clausen. Fortolkning af de synoptiske Evangelier.
Første Deel. Kjøbenhavn 1850. S. 7.}
--- „i deres hele Form, er Vidnesbyrdet givet i
den Magt, hvormed den christelige Kunst i sin skjønneste
Blomstringstid har følt sig fortrinsviis hendraget til denne
Række af Scener, som --- selv en Billedkreds, hvor det
Underfulde og Hemmelighedsfulde træder frem i sin klareste
objective Anskuelighed, i Forening af den festlige Høihed og
Eenfoldighedens rigeste Ynde --- har gjennem billedlige Fremstillinger
og lyriske Bearbeidelser afgivet rige Midler til at
vække og nære den kirkelige Andagts Liv.“

Maae vi nu end fra den humane Dannelses Standpunkt
finde det her Anførte at være baade sundt og træffende sagt,
saa er det dog ligefuldt vist, at den hele Opfattelse øiensynlig
lider af den i Culturbevidstheden grasserende Svaghed, at
foretrække en blid Sammensmeltning af det Religiøse og det
Æsthetiske for en principiel Adskillelse. Er allerede den
æsthetiske Betragtning af den Grund ueensartet med den
% --- p. 58 ---
\opage{58}\setcounter{footnote}{0}%
videnskabelig kritiske, at hiin ikke tillader nogen endelig Strid
mellem Indholdet og den individualiserende Form, gaaer
denne bestandig ud paa at bringe en saadan Strid ind i
Bevidstheden. Men er den æsthetiske Opfattelse relativt
ueensartet med den videnskabelig kritiske, da er den absolut
ueensartet med den religiøse. Den religiøse Bevidsthed lader
Foreningen af de naturlige og overnaturlige Begivenheder
foregaae i en høiere Virkelighed, en Virkelighed, som Troen
levende erkjender, en Virkelighed, hvorom Viden Intet veed.

Før vi gaae over til at omhandle de Fordringer, der
maa gjøres til „den gode Villie“, dersom det skal lykkes den
at anvise Bibelkritiken sin rette Plads i Monarchiet, skulle
vi i faa Træk antyde vor egen Synsmaade. At Forholdet
mellem Tro og Viden som to absolut ueensartede Principer
maa vise sin Indflydelse paa Behandlingen af den hellige
Historie, er en Selvfølge. Vi have under første Afdeling
nærmest opfattet Forholdet med Hensyn paa \emph{Mysteriet},
vi skulle nu opfatte det med særligt Hensyn paa \emph{Bevidstheden}.
Naar Tro og Viden forenes i een Bevidsthed, er
det nødvendigt at anvende Kritik for at forebygge Blandinger.
Kritiken er dobbelt: en Kritik, der sørger for, at de to Erkjendelseskredse
under alle Vexelbestemmelser, Grændsesætninger
og Reflexbelysninger holdes ud fra hinanden, og en Kritik,
der paaseer, at der i enhver af de adskilte Kredse gjøres
Forskjel mellem Væsentligt og Uvæsentligt, Nødvendigt og
Tilfældigt, Gyldigt og Ugyldigt. Under den dobbeltsidige
Anvendelse af Kritiken fremkommer da en \emph{Troeskritik},
der staaer i samme Forhold til \emph{Videnskritiken} som
Troen til Viden: er Troen absolut ueensartet med Viden,
saa er Troeskritiken absolut ueensartet med Videnskritiken.

Til at oplyse Modsætningen mellem Troeskritik og
Videnskritik anføre vi nogle Exempler af Evangelierne, idet
vi særlig henvende Opmærksomheden paa Bestemmelsen af
Forholdet mellem Idee og Virkelighed, mellem det Naturlige
% --- p. 59 ---
\opage{59}\setcounter{footnote}{0}%
og det Overnaturlige, mellem Mythe og Aabenbaring. Det
skal da vise sig: 1) at Troeskritiken ved at nedsætte det
udvortes Historiske til Moment for det indvortes (\textit{fides
religiosa}) bygger paa Ideen, Videnskritiken derimod ved
at gjøre det udvortes Historiske (\textit{fides historica}) til Hovedsagen,
bygger paa Kjendsgjerningen; 2) at Troeskritiken
med den almægtige Villie til Princip og Underet til absolut
Forudsætning bedømmer Vexelvirkningen af Naturligt
og Overnaturligt fra Synspunktet af det Religiøse,
Videnskritiken derimod med stiltiende Negtelse af Underet
arbeider paa at faae de naturlige Bestanddele adskilte fra de
overnaturlige; 3) at Troeskritiken vurderer det historiske
Charakteerbillede efter dets Betydning for en \emph{Aabenbaring},
Videnskritiken derimod efter dets Betydning som \emph{Mythe}.

Ved Evangeliet om Christi Fødsel f. Ex. vil den beskrevne
Modsætning gjøre sig gjældende i alle Henseender.
For Troeskritiken er Spørgsmaalet om, hvad man i historisk
Forstand kalder Ægtheden, et aldeles underordnet Spørgsmaal.
Kritiken staaer paa Menighedens Standpunkt; den
tvivler slet ikke om den historiske Ægthed, har i det Hele
liden eller ingen Interesse af at prøve Kjendsgjerningerne i
deres Enkelthed; derimod har den et skarpt Blik for den religiøse
Ægthed; den lægger an paa at udvikle og tydeliggjøre
den religiøse Modsætning imellem det Kanoniske og det Apokryphiske.
Men dersom den nu fik at vide, hvor unøiagtig
Beretningen er angaaende Qvirinius og Skatteudskrivningen?
Den faaer det aldrig at vide, den kan aldrig faae det at
vide; thi i det Øieblik, den faaer det at vide, er den ikke
Troeskritik, men Videnskritik. Hvad der er Troeskritiken et
Ligegyldigt, et Underordnet og Uvæsentligt, er derimod Videnskritiken
et Afgjørende. Videnskritiken kan ikke og vil ikke
lukke Øinene for vitterlige Uovereensstemmelser; den hverken
kan eller vil i denne Henseende gjøre nogen Undtagelse med
den hellige Historie. Saavidt om Ægtheden.

% --- p. 60 ---
\opage{60}\setcounter{footnote}{0}%
Hvad dernæst Forholdet imellem Beretningens naturlige
og overnaturlige Elementer angaaer, vil det være indlysende,
at Troeskritiken umulig kan lægge an paa at adskille, hvad
den hellige Villie til virkelig og virksom Fremstilling af den
religiøse Idee har sammenføiet. For Troeskritiken have
Englene med deres Lovsang en ligesaa umiddelbar Tilværelse
som Hyrderne med deres Hjord. Forskjellen er ikke den, at
Hyrderne ere virkelige, Englene uvirkelige, eller at Hyrderne
have en egentlig, Englene kun en uegentlig Virkelighed; nei
Forskjellen er tvertimod den, at Englene have en himmelsk og
salig Virkelighed, Hyrderne kun en jordisk og forkrænkelig.
Men dersom der nu opkom Tvivl om det Overnaturliges
Forhold til det Naturlige, f. Ex. Tvivl om, hvorledes Stjernen
paa Himlen kunde bogstavelig vedblive at gaae foran de Vise,
indtil den stod stille over Krybben i Bethlehem? En saadan
Tvivl kan aldrig opkomme; thi i det Øieblik en saadan Tvivl
opkommer, i det Øieblik der lægges an paa at udskille det
Overnaturlige som en mistænkelig Tilsætning, i det Øieblik
har Troeskritiken givet Plads for en Videnskritik. Men,
vil man maaskee indvende, naar ingen Reflexioner anstilles
over Muligheden af det i Phænomenerne Overnaturlige, saa
er her jo ingen Kritik. Vi svare: her er en Kritik, forsaavidt
her med Hensyn paa det Religiøse selv maa gjøres en
skarp Adskillelse mellem Væsentligt og Uvæsentligt. Ja der
kan i en forøvrigt kanonisk Beretning være apokryphiske, det
Religiøse fordunklende Bestanddele; Troeskritiken retter sig
ikke efter en ubestemt Følelse, den retter sig efter et Princip.

Hvad endelig Opfattelsen af en Evangelieberetning som
historisk Charakteerbillede angaaer, vil jo Modsætningen mellem
Troeskritiken og Videnskritiken være iøinefaldende. For
Troeskritiken er det en Opgave af afgjørende Betydning at
gjøre Skilsmisse imellem Phantasiens Ideer og den absolute
Villiesidee; kun ved at gjennemreflectere denne Adskillelse kan
den Troende klare sig Troens Princip og i Conseqvens af
% --- p. 61 ---
\opage{61}\setcounter{footnote}{0}%
dette Princip bestemme Modsætningen imellem Aabenbaring
og Mythe. For Videnskritiken gives der ingen Aabenbaring;
hvad vi kalde Aabenbaring er kun en egen Form af det
Mythiske: at Christendommen er den fuldkomneste Religion,
vil sige, at Christendommen er den fuldkomneste Mythologie.
For Videnskritiken med Mythen til Maal maa Fuldendelsen
af det historiske Charakteerbillede bestaae deri, at det gjør
Phantasie-Ideen anskuelig, at det tilfredsstiller æsthetisk. „For
den religiøs-æsthetiske Betydning af disse Fortællinger“ ---
hedder det om % Rettelser: "hedder det om" (p. 61 l. 10 fra oven), læs: "hedder det, som ovenfor anført, om"; the KB copy has an ink caret here and "som ovenfor anført" in the margin. Printed text kept.
Johannes Døberens og Jesu Fødsels- og
Barndomshistorie\footnote{\textit{Dr.} H. N. Clausen. Fortolkning af de synoptiske Evangelier.
Første Deel. Kjøbenhavn 1850. S. 3.% Rettelser: "S. 3" (p. 61 Anm. Lin. 1 fra neden), læs: "S. 7."; the KB copy has the 3 struck through in ink and 7 written beside it. Printed "3" kept.
}
--- „i deres hele Form, er Vidnesbyrdet
givet i den Magt, hvormed den christelige Kunst i sin skjønneste
Blomstringstid har følt sig fortrinsviis hendraget til
denne Række af Scener, som --- selv en Billedkreds, hvor
det Underfulde og Hemmelighedsfulde træder frem i sin
klareste objective Anskuelse, i Forening af den festlige Høihed
og Eenfoldighedens rige Ynde --- har gjennem billedlige
Fremstillinger og lyriske Bearbeidelser, afgivet rige Midler til
at vække og nære den kirkelige Andagts Liv“. Dette er
tilvisse baade træffende og smukt sagt --- nemlig fra et religiøst-æsthetisk
Synspunkt. Men, spørge vi, er Principet for
det kirkelige Andagtsliv et æsthetisk Princip? Ved blot religiøs-æsthetiske
Fortrin vil Aabenbaringen aldrig kunne adskilles
fra Mythen; æsthetisk-religiøse Fortrin vil en Strauß
jo ogsaa indrømme vore kanoniske Forevangelier. Nei Modsætningen,
den qvalitative Modsætning mellem Videns- og
Troeskritiken viser sig just deri, at Troeskritiken sætter en
Sphæreforskjel imellem æsthetisk Contemplation og „kirkelig
Andagt“. At Andagten kan have det Æsthetiske som Moment,
skal ikke benegtes; men i Principet er det Æsthetiske
% --- p. 62 ---
\opage{62}\setcounter{footnote}{0}%
(Billedskjønheden, Poesien) et for Andagten selv Tilfældigt,
det Religiøs-Ethiske det ene Fornødne, det absolut Væsentlige.

Troeskritiken er som Kritik ikke opbyggelig, den bevæger
sig i Reflexioner og Begrebsanalyser ligesaavel som Videnskritiken;
men da hele Tankebevægelsen foregaaer i Troens
Sphære, da Begreberne, der analyseres, ere Troesbegreber,
saa vil det være en af de store Opgaver, at udvikle og ved
Fortolkning af de enkelte Evangelieberetninger tydeliggjøre det
Forhold imellem Historisk og Overhistorisk, Forestillen og
Tænken, Totalanskuelse og Indtryk af det Individuelle, der
væsentlig tilfredsstiller Andagten, idet Andagtens Begreb med
det Samme udvikles.

Spørge vi til Slutning, om her i en Bevidsthed, der
skal kjende og erkjende baade Troens og Videns Princip, da
ikke er et særligt psychologisk Medium, hvori de „to absolut
ueensartede Principer“ kunne mødes og skilles og paa en
Maade træde i Underhandling med hinanden, svare vi:
Jo! her er ganske rigtigt et saadant Medium. Men hvad
skulle vi vel kalde dette Medium? Viden kunne vi da ikke
kalde det; thi Viden hører til Videnskaben og er ikke neutralt
Gebet. Tro kunne vi heller ikke kalde det, thi Troeserkjendelsen
udelukker Videnskaben, og frembyder altsaa heller
ikke noget Neutralt. Vi maae da kalde det \emph{Mening}.
Mening er det rette Navn, den kategoriske Betegnelse for det
søgte Medium. Den vældige Modsætning mellem Oldtidens
og Nutidens Anskuelse af Himmel og Jord, en Modsætning,
der ikke blot angaaer Phænomenerne i Tid og Rum, men
Phænomenernes Aarsager, Betingelser og Væsenhed, en Modsætning,
der fra Culturbevidsthedens Standpunkt maa synes
at undergrave det Gamle og det Nye Testamentes og dermed
Aabenbaringens guddommelige Autoritet, vil, ikke forsvinde,
men lade sig udligne ved fortsatte Forhandlinger paa
Meningens neutrale Grund.

% --- p. 63 ---
\opage{63}\setcounter{footnote}{0}%
Den væsentlige Forskjel imellem \emph{Viden} og \emph{Mening}
blev allerede Grækerne klar; men den principielle Forskjel
imellem Tro og Mening synes endnu ikke at være bleven
Theologerne klar. Ligesom Forskjellen mellem Viden og
Menen % printer's defect: "Menen" (p. 63 l. 5), for "Mening"; transcribed as printed
kun kan klares ved en Videnskritik, saaledes kan
Forskjellen mellem Tro og Mening kun klares i en Troeskritik.
Et Exempel. Hvad der i første Mosebog fortælles
om Paradiset, om Livets Træ, Kundskabens Træ o. s. v.,
er ganske vist fortalt i den Mening, at det Altsammen hører
under det udvortes Virkelige, at Livets Træ og Kundskabens
Træ ligesaavist have havt deres Rødder i Havens, sandselige
materielle Jordbund, som enten Viintræet eller Figentræet.
I denne Mening have vore luthersk-orthodoxe Fædre endnu
for hundrede Aar siden ganske vist opfattet Fortællingen om
Paradiis. Men i denne Mening kunne vi paa Culturbevidsthedens
nuværende Trin umulig opfatte den. Dersom
der nu ingen væsentlig Forskjel var mellem Tro og Mening,
saa vilde en Forsoning af Religion og Cultur ganske vist
være umulig; men er her en Forskjel, en principiel Forskjel,
saa er Forsoningen mulig.

Men at her virkelig er en principiel Forskjel mellem
Tro og Mening, skal kjendes derpaa, at vi kunne bevare
Fædrenes Tro og forkaste deres Mening. Til Troens ufravigelige
Forudsætning hører, at Fortællingen om Paradiset
med alle deri indeholdte Charakteertræk tillægges væsentlig og
sand Realitet, væsentlig og sand Virkelighed. Men hvilken
Virkelighed? En ydre eller en indre? Ja, derom er der
forskjellige Meninger. Men saa er det jo kun en Mythe;
eller rettere: hvorvidt det er Mythe eller Aabenbaring, vi
have for os i dette Sagn, derom kan der dog vel være forskjellig
Mening? Nei, Sagnet om Paradiis kan, saa sandt
det i rene, klare Træk for Troens Blik afslører et af de
dybeste Villiesmysterier, umulig være Mythe. Men hvorledes
skal man da forestille sig Forholdet mellem det Indre
% --- p. 64 ---
\opage{64}\setcounter{footnote}{0}%
og det Ydre, mellem det som skeer legemligt og det som foregaaer
i Aandens Verden? Ja, derom er der forskjellige
Meninger.

At det, naar virkelige Analyser foretages, maa paa Meningernes
Omraade komme til mangfoldige Conflicter og
Forhandlinger mellem Tro og Viden, vil af det her Anførte
forhaabentlig være klart. Vi have alt i længere Tid bearbeidet
disse Forudsætninger og ved deres Anvendelse paa
Fortolkningen af Kjendsgjerningerne i den hellige Historie
øinet en Udsigt til, at de gamle Stridigheder om den hellige
Skrifts Troværdighed, maatte, vistnok ikke ende, men dog
indtræde i en ny Phase. Fra dette Synspunkt var det da,
saa langt fra, at Nogen i Menighedens Navn behøvede at
tage Forargelse enten af den naturlige eller den mythiske
Fortolkning, at han meget mere maatte paaskjønne den Hjælp
og Tjeneste, Videnskaben ved sin gjennemførte Kritik middelbart
har ydet Troen. Havde skarpsindige Mænd som Paulus,
Eichhorn og Strauß ikke ved aaben, conseqvent og resolut
Fremgangsmaade udviist de paa Videns Vei charakteristiske
Stadier, skulde det endnu have faldet Evangelietroen vanskeligt
nok at finde tilbage til sit oprindelige Udgangspunkt og
forvisse sig om sit eiendommelige Princip.

Disse bibelske Studier, disse Forberedelser, ak, disse
skjønne Forhaabninger maae forsvinde som Dunst, naar Biskop
Martensen nu gaaer hen og opretter sit theologiske Universalmonarchie.
Det vil for Resten koste endeel Arbeide,
inden et saadant Monarchie bliver oprettet. Føler Biskop
Martensen virkelig Aandens Kald og Naadegave, og det maa
han vel føle, siden han melder sig med et saa stort Forehavende,
til at fuldbyrde, hvad „Skriften tydelig siger“ ---
med udtrykkeligt Hensyn paa Videnskaben? --- at Alt, baade
det Synlige og Usynlige skal“ --- videnskabeligt? --- „indsamles
under Christus som Hovedet (ἀνακεφαλαιώσασθαι
τὰ πάντα ἐν χριστᾷ Eph. 1), % printer's defect: the Greek has "χριστᾷ" (alpha with circumflex and iota subscript) for "χριστῷ" (Eph. 1:10); transcribed as printed
at Alt er System“ ---
% --- p. 65 ---
\opage{65}\setcounter{footnote}{0}%
videnskabeligt System! --- „i ham (πάντα ἐν αὐτῷ συνέ\-%
στηκε Coll. 1)“ (Leilighedsskr. S. 80): da maa Biskoppen
i Sandhed være beredt paa en kraftig Optræden mod Videns
mange Afguder: en anden Elias, en „philotheoretisk“ Thesbiter,
maa han fra sit speculative Karmel lade Ild regne ned
over „Autonomien“, den ugudelige Autonomie, der bedaarer
Philosophien, frister Theologien, boler med „Heteronomien“ og
offrer „Theonomien“ til Kritikens Moloch.

Her maa paa Historiens Vegne ryddes op, og det tilgavns,
ikke blot udenfor, men ogsaa indenfor Theologiens
Leir. Det er nemlig ikke nok, at „Antichrister“ som Strauß
og Renan lyses i Band; nei, Theologiens eget Land maa
indtil videre belægges med Interdict. Her er jo Theologer,
store Theologer, Mænd som Semler, Eichhorn, Paulus fra
Heidelberg, Schleiermacher, de Wette, Daub, Marheineke
m. fl., om hvilke det kan bevises, at de alle have været med
at smage paa Kundskabens Frugt, alle have viist en afgjort
Forkjærlighed for det Logisk-Physiske, alle have bøiet Knæ for
den rationelle Nødvendighed og offret til Kritikens Moloch.
De endnu levende Theologer kunne maaskee omvende sig; men
det kunne de døde jo ikke. Der er ved en Antinomie i Dogmatiken
udtrykkelig gjort Indsigelse imod de Dødes Omvendelse.
Der skulde til Antinomiens Løsning dog vel aldrig
gjøres en videnskabelig Undtagelse med Theologerne; saa at
Theologer, der medens de levede havde forsømt at skifte Standpunkt
i rette Tid, fik Lov til at skifte, naar de vare døde? Dersom
dette ikke gaaer an, maa det især gaae ud over de Døde. Her
vil ved saa mange og betydelige Theologers Excommunication,
blive et Hul, som vi, da de Excommunicerede ved deres Anseelse
og Mængde repræsentere en heel Tidsalder, maaskee kunne
kalde Frafaldets, det store Frafalds Hul i Theologiens Historie.
Det vil, derom ere vi overbeviste, gjøre Biskop
Martensen ondt for mange af disse Hædersmænd, ikke for
Eichhorn, ikke for Paulus, men for Daub, for Marheineke
% --- p. 66 ---
\opage{66}\setcounter{footnote}{0}%
og især for Schleiermacher. Gjerne havde Theologien bevaret
Schleiermacher som et Sandhedsvidne, et ægte Sandhedsvidne;
men, ved det kongelige Princip! det er umuligt.
En Theolog, hvis Kritik ikke engang lader Evangelierne være
ægte, hvorledes skulde han selv være ægte?

Hvad i al Verden kan det hjælpe at skjende paa den
Viden, der sublimerer Naturalisme og Rationalisme til speculativ
Theologie, aldenstund Naturalismen er Theologiens
egen Naturalisme, Rationalismen Theologiens egen Rationalisme,
aldenstund det endnu fra det theologiske Høidepunkt
lyder ned til os, at det høieste Værende tillige er det dybeste
Speculative? Hvordan, om vi maae spørge, ere de Herrer
Rationalister med Et komne ud over deres Rationalisme?
Er det ved at sige: „historisk“ og „noget Historisk“ og „et vist
Grundfactum“, eller ved at declamere om „en Historie, der
netop som historisk Virkelighed er bleven Vendepunktet i den
religiøse Verdenshistorie“? Eller er der maaskee opdaget et nyt
Amalgam af Orthodoxie og Rationalisme? Hvad kan det
hjælpe at ville lægge Dølgsmaal paa, hvad dog fra alle
Sider kan sees, at Straußes „Jesu Liv“ er et theologisk
Speil, hvori Tidsalderens Aand kan betragte sin egen Skikkelse.
„Gives der“, siger Baur\footnote{Jfr. \textit{Dr.} H. N. Clausen: Fortolkning af de synoptiske Evangel.
1ste Deel S. 74.}, „noget Værk, der saavidt
det er muligt, alene er bygget paa Forgængernes Arbeider,
alene vil sammenstille Resultater af en Række, af mange
Andre for længe siden førte Undersøgelser, alene uddrage den
sidste Conseqvens af Præmisser, om hvilke man forud er
bleven enig, saa er det dette. Naar man lægger den
Straußiske Kritik Negativiteten af alle dens Resultater til
Last, saa sige man dog, hvad andet Resultat end et blot
negativt der lod sig uddrage af de daværende Undersøgelser
om Evangeliernes Oprindelse og deres indbyrdes Forhold!

% --- p. 67 ---
\opage{67}\setcounter{footnote}{0}%
Idet Mening stod overfor Mening, og alle Meninger til\-%
sammen modsagde og ophævede hverandre indbyrdes, faldt
alle sammen i en Indifferens, hvor man ikke havde noget
Fast og Sikkert at holde sig til.“ Skal Strauß excom\-%
municeres, da er der Mange, der maae excommuniceres; er
Strauß „en Antichrist“, da er der i Leiren mange „Anti\-%
christer“. Skal Riget oprettes, da maae Straffedomme gaae
forud; da maa Elias komme med Ildsluer, i Vredens
Storm, med gloende Vogn --- dogmatisk gloende Heste.

Det er, for at blive ved Forbillederne, til Helligdom\-%
mens Rensning ikke engang nok, at Biskop Martensen sidder
dogmatisk i en gloende Vogn og er Elias; han maa, hvis
det hele Forehavende ikke skal ende med „en literær Skan\-%
dale (\textit{parturiunt montes})“, nødvendigviis ogsaa være en
Judas Makkabæus: kun en Judas Makkabæus kan rense
Helligdommen og feire den Tempelvielsesfest, som her skal
feires. Men for at rense Helligdommen og feire den
store Fest er der eet Punkt, rigtignok kun eet eneste, men
et vigtigt Punkt, han absolut maa klare. \emph{Hvorledes for\-%
holder det Logisk-Physiske sig til det høiere Vi\-%
denskabelige?} Paa dette Spørgsmaal maa der svares
reent og tydeligt. Uden en tydelig Besvarelse af dette Spørgs\-%
maal vide vi jo i Grunden hverken ud eller ind; og naar
vi midt i al vor dogmatiske Viden hverken vide ud eller ind,
hvad kan det saa hjælpe, at vi anstrenge os med at være
enten Elias eller Makkabæus? Endnu engang: \emph{Hvorledes
forholder det Logisk-Physiske sig til det høiere
Videnskabelige?} hvor er Grændsen?

Biskop Martensens Udtalelser angaaende dette vigtige
Punkt ere jo vistnok, hver for sig, meget bestemte; men idet
den ene Udtalelse saa bestemt gaaer i een Retning, den anden
i en anden, bliver „Totalanskuelsen“ paa Grund af lutter
Bestemthed tilsidst ubestemt.

Biskop Martensen maa, det siger han meget bestemt,
% --- p. 68 ---
\opage{68}\setcounter{footnote}{0}%
vedblive „fremdeles at betragte Dogmatikerens Forhold til
Aabenbaringen paa tilsvarende Maade som Naturforskerens
til Naturen“. Men Biskop Martensen maa, det siger han
ligesaa bestemt (Leilighedsskr. S. 93), erkjende, „at med det
ethiske Gudsbegreb indtræder et andet Vidensbegreb end der,
hvor der kun er Tale om en blot logisk og physisk Viden.
Thi er Gud den personlige, villende Gud, da kunne vi
hverken erkjende Mere eller Mindre af ham, end han selv
vil aabenbare os“. For nu ret paa det Bestemteste at styrke
denne sidste Sætning til uforsonlig Strid med den første,
tilføies: „Apriori construere ham (den ethiske Gud) kunne vi
ikke, thi kun det logisk Nødvendige, men ikke den frie Aand
lader sig apriori construere.... \emph{Tvinge} den ethiske Gud til
at aabenbare sig for os, som vi ved physicalske Experimenter
kunne tvinge Naturen til at røbe sine Hemmeligheder, kunne
vi heller ikke“; med andre Ord: „betragte Dogmatike\-%
rens Forhold til Aabenbaringen paa tilsvarende Maade som
Naturforskerens til Naturen“, kunne vi heller ikke. „Reent fri\-%
villigt maa det Godes Gud meddele sig til os, og vi vilde
ikke engang kunne søge ham, dersom denne Søgen ikke stam\-%
mede fra ham selv“. Altsaa: ligerviis som Naturforskerens
Forhold til Naturen beroer paa det Logisk-Physiske, forsaa\-%
vidt det beroer paa apriorisk-empirisk Fornuftnødvendighed,
saaledes beroer Dogmatikerens Forhold til Aabenbaringen
\emph{ikke} paa det Logisk-Physiske, forsaavidt det ikke beroer paa
apriorisk-empirisk Nødvendighed; ganske vist er her et Liger\-%
viis, men et Ligerviis af denne Natur minder stærkt om det
Bekjendte: ligerviis som Loven er et glubsk Dyr, saa skulle vi
og vandre udi et nyt Levnet.

Det nytter ikke med Talemaader og Totalanskuelser.
Skulle vi gaae ind i Universalmonarchiet og ved \textit{credo, ut
intelligam} gjøre „den gode Villie“ til Herre over Viden\-%
skaben, da maa den gode Villie virkelig have den Godhed at
udstede sine Ordrer i bestemt Form, give os tydelig Besked
og ikke afspise os med Totalanskuelser.

% --- p. 69 ---
\opage{69}\setcounter{footnote}{0}%
Dette bemærker jeg dog ikke for min egen Skyld, thi
jeg agter ikke at gaae ind i Universalmonarchiet, men for
deres Skyld, som enten allerede ere derinde eller dog maatte
ønske og begjære at komme derind; navnlig tænker jeg igjen
med Bekymring paa Facultetets synoptiske Apologetik, thi den
synoptiske Apologetik er virkelig --- midt paa Middelveien ---
i Forlegenhed med Grændsen for det Logisk-Physiske.

Det er ikke Lærdom, ikke Skarpsindighed, ikke kritisk
Takt og Fiinhed, der mangler den synoptiske Apologetik; der
mangler den kun Eet: et Princip. Paa en principløs Maade
har den, for at gjentage et af Biskoppens Yndlingsudtryk,
ved „\emph{Concordatet}“ imellem Tro og Viden indladt sig paa
Forhandlinger mellem det Physiske og det Hyperphysiske, der
ikke kunne Andet end halvere dens Bevidsthed og paaføre den
halverede Bevidsthed mangfoldige logiske Ulemper.

Man kaste engang et Blik paa et af Evangelierne f. Ex.
det om Bespiisningen af de fem Tusinde i Ørken. For den
exegetiske Betragtning, det være sagt til Troens Trøst, har dette
Under „saa fuldstændig udvortes Hjemmel som nogen af de
evangeliske Undergjerninger“. Men naar Exegeten nu --- for at
fyldestgjøre Videns Fordring --- skal, ifølge Biskop Marten\-%
sens dogmatiske Anviisning, forholde sig til denne hellige
Kjendsgjerning „paa tilsvarende Maade som Naturforskeren
til Naturen“, hvad seer han saa? Han seer „en Magthand\-%
ling, i absolut Ubegribelighed, af den skabende Guddoms\-%
kraft, et τέρας, som tillukker sig for enhver Tanke\-%
anstrengelse“: nu er Bevidstheden halveret. Bevidsthedens
Halvering er her saa slaaende, at man næsten med anatomisk
Nøiagtighed kan paavise, hvorledes den ene Halvdeel holder
paa Troen, den anden paa Viden.

Fortolkeren bør, saa lyder Fordringen fra den troende
Halvdeel, hæve sig over det Logisk-Physiske, bør i Alvor
fastholde „det ethiske Gudsbegreb“ og trøste sig med, at efterdi
Gud er „den personlige, villende Gud, saa kunne vi hverken
% --- p. 70 ---
\opage{70}\setcounter{footnote}{0}%
erkjende Mere eller Mindre af ham, end han selv vil aaben\-%
bare os“; nu har „den personlige, villende Gud i det fore\-%
liggende Evangelium blot villet aabenbare os, at, men ikke
\emph{hvorledes}, han kan mætte fem Tusinde med fem Brød og
to Fiske“. Vi skulle, med andre Ord, holde os til det
Ethiske uden at spørge om det Physiske. Herved er Biskoppens
ene Anviisning efterkommet.

Men --- mærkeligt nok, hvor Alt slaaer til i Theo\-%
logien --- ligerviis som den halverede Bevidsthed trænger til
en halveret Anviisning, en halv for Troen og en halv for
Viden, saaledes har ogsaa Biskop Martensen givet Apologe\-%
tiken en, i to halve faldende, heel Anviisning: en halv, der
ifølge „det ethiske Gudsbegreb“, passer for Troen, og en halv,
der i Lighed med Naturforskerens Forhold til Naturen, passer
for Viden. Fra sidste Halve lyder da den Fordring, at
Fortolkeren skal som theologisk Naturforsker bestræbe sig for
at trænge ind i Begivenhedens Indre. „Den ægte Natur\-%
forsker opgiver ikke Naturen, fordi han ikke seer sig istand til at
forklare ethvert af dens Phænomener, men opgiver dog heller
ikke at trænge ind i Naturens Indre, saavidt han formaaer
det.“ Den ægte Theolog opgiver ikke Miraklet, fordi han
ikke seer sig istand til at forklare det --- hvem giver sig vel
af med Mirakelforklaringer? Men han søger dog, saavidt
han formaaer, at trænge ind i Miraklets Indre. Den ægte
Theolog kan ikke undlade at spørge om det Physiske: hvis
han oversaae det Physiske i det Hyperphysiske og tog til
Takke med det Ethiske, var han jo ingen ægte Theolog;
han maa som ægte Theolog da spørge, „om Levnets\-%
midlernes Forøgelse skal tænkes at være foregaaet i Jesu
Hænder eller i Disciplenes og hvorledes foregaaet: om ved
Udvoxning (men Talen er her ikke om Naturlegemer, der
endnu staae under en organisk Udviklingslov) eller ved
Mangfoldiggjørelse (Orig. \textit{in Matth.} 1, 11: τῷ λόγῳ
καὶ εὐλογίᾳ αὔξων καὶ πληθύνων. Clem. Strom.
% --- p. 71 ---
\opage{71}\setcounter{footnote}{0}%
VI, 11): ved Brødene en Voxen, ved Fiskene en Opstandelse
o. s. v.“\footnote{\textit{Dr.} N. H. Clausen Fortolkn. af synopt. Evang. 1ste Deel S.
347. Anm.% printed "N. H." (inverted initials; H. N. Clausen, as on p. 66), checked at 300 ppi; transcribed as printed
}

Naar nu den ægte Theolog ved at spørge videnskabeligt
har forvisset sig om, at her ikke kan gives et eneste fornuftigt
Svar paa hans mange fornuftige Spørgsmaal, ja, forvisset
sig om, at et i sig urimeligt Mirakel, „et τέρας, der til\-%
lukker sig for enhver Tankeanstrengelse, afviser ethvert Spørgs\-%
maal som ubesvarligt, staaer for os“: hvad gjør han saa?
Opgiver han maaskee et saadant Mirakel eller fordrer det
henregnet blandt apokryphiske Bestanddele, der have indsneget
sig i det Nye Testamente? Ingenlunde! Det vilde jo være
stik imod „Concordatet“, et uforsvarligt Overgreb af Viden.
Kun hos en Strauß og de med ham Ligesindede finde vi slige
Overgreb, slige iøinefaldende Misgreb! Kun en Strauß vil sige
det reent ud, at „dette Underværk hører til dem, der kun
saalænge kan synes paa nogen Maade troligt, som man veed
at holde det i en ubestemt Forestillings Halvmørke: saasnart
man trækker det frem for Lyset og nøie vil undersøge alle
dets Dele, opløser det sig i en Taageskikkelse. Brød, der i
de Uddelendes Hænder opsvulme som Svampe, man har
vædet; stegte Fisk, hvoraf de afbrukne Dele voxe efterhaanden,
som Saxen man river af en levende Krebs, høre aabenbart
ikke til i det Virkeliges Rige, men i et ganske andet.“\footnote{Jesu Liv § 102.}

Men derfor er ogsaa Strauß, hvad man seer af denne
Yderlighed, ikke nogen ægte Theolog; der mangler ham
Noget: han har en Viden, men han mangler Troen. At
mangle Troen er i theologisk Forstand at mangle Bevidst\-%
hedens ene, og vistnok bedste, Halvdeel; at mangle Troen
er at miskjende „Concordatet“.

% --- p. 72 ---
\opage{72}\setcounter{footnote}{0}%
Det vilde være et Brud, et syndigt Brud paa „Con\-%
cordatet“, om nogen Theolog opgav Troen (\textit{fides historica})
paa Virkeligheden af den mirakuløse Bespiisning, „eftersom
dette Under for den exegetiske Betragtning har saa fuldstændig
udvortes Hjemmel som nogen af de evangeliske Undergjer\-%
ninger“. Vi maae bestandig huske paa, at det er den \emph{histo\-%
riske Tro}, der ifølge „Concordatet“ afgiver det rette exegetisk
paalidelige Grundlag for den religiøse. Skulde den ægte Theo\-%
log. opgive Troen paa det Historiske i Bespiisningens Under,% sic: a full stop after "Theolog" in the print (both witnesses and the 300 ppi image); transcribed as printed
maatte han jo opgive Troen paa det Historiske i hele den hellige
Historie. Men, spørge vi, hvad er da nu i denne Bespiis\-%
ningshistorie det egentlig Historiske? Derpaa maa den ægte
Theolog give os det bestemte Svar: \emph{det veed man ikke}!

Dersom de ægte Theologer endda kunde blive staaende
ved en Halvering af Bevidstheden: det fik at være; de kunde
da troe paa Bespiisningsmiraklet med den ene Halvdeel og
vide, at Miraklet er utroligt, med den anden. Men det
gaaer paa ingen Maade an; det vil Biskop Martensen
ikke tillade. Enhver Fornuftig maa give Biskoppen Ret
i, at en „Dualisme mellem det Theoretiske og det Prak\-%
tiske“ som den her beskrevne ikke kan løses i „et selv\-%
modsigende Concordat“. Men hvorledes vil Biskop Mar\-%
tensen nu i Virkeligheden --- thi ved almindelige Talemaader
udrettes Intet --- bære sig ad med at hjælpe den dybt i
„Concordatet“ stikkende, mellem Tro og Viden uophørlig vrik\-%
kende Apologetik ud af Halveringen?

Ja, er her i nogen videnskabelig Bevidsthed Tegn paa
dualistisk Halvering, da maa det være i den synoptisk-apolo\-%
getiske. Man tænke sig en Conflict mellem Tro og Viden
som denne: Troen fordrer, at jeg i Bespiisningsmiraklet skal
erkjende det historisk Virkelige; thi hvis jeg fornegter det, for\-%
negter jeg den evangeliske Historie. Viden viser mig, at jeg
umulig kan enten tænke eller forestille mig det i Bespiis\-%
ningsmiraklet historisk Virkelige. Jeg skal altsaa troe paa
% --- p. 73 ---
\opage{73}\setcounter{footnote}{0}%
Virkeligheden af det, hvis Virkelighed jeg hverken kan tænke
eller forestille mig. Hvad Troen fordrer, er umuligt for
Viden; hvad Viden fordrer, er umuligt for Troen: der\-%
som dette ikke kan kaldes Dualisme, hvad kan da kaldes
Dualisme?

Da nu Intet er Universalmonarchiet saa farligt som
Dualisme og selvmodsigende Concordater, saa maa Biskop
Martensen, der vil oprette Monarchiet, jo vide Udvei til at
hæve den synoptisk-apologetiske Dualisme, det synoptisk-apo\-%
logetisk selvmodsigende Concordat.

Men hvordan skal Dualismen hæves? At gjøre en
kritisk-subtil Adskillelse mellem Sagen i sig og dens Ind\-%
klædning, give den evangeliske Fremstilling til Priis som „en
mythiserende“ Opfattelse og beraabe sig paa et ubekjendt og
uerkjendeligt „Grundfactum“ som paa en kantisk Ding an
sich, gaaer paa ingen Maade an. For det Første strider
det imod den Agtelse, vi skylde Evangelisternes Blik for det
historisk Virkelige. Staaer nemlig den fremstillede Begiven\-%
hed paa samme Trin af Troværdighed som en hvilkensom\-%
helst anden evangelisk Undergjerning, hvorledes vil man da
kunne drage Fremstillingens historiske Charakteer i Tvivl
uden at drage Troværdigheden af den hele evangeliske Historie
i Tvivl? Tvivl om Evangelisternes Evne til historisk
Fremstilling er \textit{eo ipso} Tvivl om den af dem fremstillede
Historie. Dernæst er Adskillelsen imellem et i sig uerkjende\-%
ligt Grundfactum og et med dette Grundfactum ueensartet
Virkelighedsphænomen, et Virkelighedsphænomen, der oveni\-%
kjøbet er ligesaa uerkjendeligt som det tilsvarende Grund\-%
factum, aldeles stridende imod Historiens, den hellige Histories
Væsen. Lader os for al Ting ikke glemme, hvad „den
christelige Dogmatik“ har lært os om den hellige Historie,
at det er en Historie, „hvor Guds Ord saaledes indfatter
sig i det menneskelige Ord, at det sidste bliver det rene
Organ for det første“, en Historie, „hvor Guds Gjerning
% --- p. 74 ---
\opage{74}\setcounter{footnote}{0}%
saaledes indfatter sig i Menneskets Gjerning, at den sidste
bliver fuldkommen gjennemsigtig for den første“; med andre
Ord: en Historie, hvor ethvert Grundfactum kan erkjendes i
dets Phænomen, idet Phænomenet bliver gjennemsigtigt, „fuld\-%
kommen gjennemsigtigt“, for sit Grundfactum.

Men hvordan skal Dualismen hæves? At beraabe sig
paa Takt og Sands, en eiendommelig Forening af religiøs
Sands og kritisk Takt, kan vist heller ikke føre til Noget.
Forstaaer sig, det kan jo vistnok føre til en synoptisk-exegetisk-%
apologetisk Taktik; men denne Taktik er af meget betænkelig
Art. Taktiken gaaer ud paa at gjøre Forskjel imellem de
evangeliske Mirakler, ikke med Hensyn paa Troværdigheden
af det Historiske --- thi det „Historiske“, det „Grundhistoriske“
maa holdes fast --- men med Hensyn paa deres, om vi saa
maae sige, religiøse Frugtbarhed. I denne Henseende ind\-%
deler man, ikke principielt, men antydningsviis, Mirak\-%
lerne i to Hovedklasser med en passende Mellemklasse. Til
første Hovedklasse høre Mirakler af religiøs Frugtbarhed,
saadanne som ret ere tilgængelige for den dogmatiserende,
ethiserende Tænkning; vi kunne kalde dem de \emph{rimelige,
næsten fornuftige Mirakler}. Til anden Hovedklasse
høre de i religiøs Forstand ufrugtbare Mirakler (Miraklet f. Ex.
med Gadarenernes Sviin); vi kunne betegne dem som de
for den sunde Menneskeforstand \emph{urimelige, næsten ufor\-%
nuftige Mirakler}.

Mon Dualismen nu kan hæves ved at indføre en
Mellemklasse og anbringe Gradationer? Ingenlunde! Det
er jo forudsat, at alle Mirakler, mod hvis historiske Tro\-%
værdighed Kritiken, den apologetiske Kritik, Intet har at
indvende, uden Forskjel skulle troes. Men Troen paa „et
urimeligt“, for Forstanden forargeligt Mirakel, er jo en
Tro paa det Absurde, en Tro, der staaer i skarp Modsæt\-%
ning til Viden, medens Troen paa et rimeligt, for den
religiøse og videnskabelige Tænkning tilgængeligt Mirakel, er
% --- p. 75 ---
\opage{75}\setcounter{footnote}{0}%
en fornuftig, med den theologiske Viden forsonet Tro. Ved
denne Taktik bliver da Tro absolut ueensartet med Tro;
Dualismen farer i Troen, som Dæmonerne fordum fore i
Svinene.

Men hvordan skal Dualismen hæves? Gaadens Løsning
kan umulig bestaae deri, at Troen og Viden hver fra sin Side
ophøre at være Principer. Vistnok bleve vi paa disse Vilkaar
frie for at plages med to absolut ueensartede Principer; ja vi
bleve endogsaa frie for at plages med noget Princip: Gaa\-%
dens Løsning var da den, at det Hele løftes op i en fuld\-%
stændig Principløshed. Men, ved Monarchiets kongelige
Princip! en fuldstændig Principløshed er Theologiens fuld\-%
stændige Ruin.

Hvordan skal da Dualismen hæves? Biskop Martensen
maa vide det; jeg veed det ikke.

\section*{III.\\ Philosophiens Stilling i Monarchiet.}
\addcontentsline{toc}{section}{III. Philosophiens Stilling i Monarchiet}
\markboth{Philosophiens Stilling i Monarchiet}{Philosophiens Stilling i Monarchiet}

En Vurdering af den Stilling, Philosophien vil komme til
at indtage i det store, theologiske Monarchie, kan maaskee bedst
indledes derved, at vi henvende Opmærksomheden paa, hvad
dog for Resten kun er af forbigaaende og psychologisk Inter\-%
esse, at Biskop Martensen har anmeldt sit monarchiske
Princip og begyndt de nye Erobringer med at skrive et heelt
Leilighedsskrift, fuldt af mange forskjellige og interessante
Oplysninger om mange og forskjellige philosophiske Concor\-%
dater og Theorier, imod min Philosophie. „Naar Professor
Nielsen“, siger Biskoppen (Leilighedssk. S. 8---9) „har stillet
sig den store og omfattende Opgave at nedlægge en Protest
imod al Theologie, har jeg i nærværende Skrift stillet mig
en langt mere underordnet Opgave, nemlig, at nedlægge en
Protest, ikke mod al Philosophie, men kun mod den R.
% --- p. 76 ---
\opage{76}\setcounter{footnote}{0}%
Nielsenske.“ Vi skulle i det Følgende forsøge at vise, hvor\-%
ledes Biskop Martensen er kommen til at løse sin vistnok i
og for sig underordnede Opgave paa en saa overordnet
Maade, at han dog virkelig --- maaskee imod sit Vidende ---
kan siges at have nedlagt en Protest imod al Philosophie.

Idet vi fra et mere underordnet, blot logisk-physisk
Standpunkt skulle see opad mod det Overordnede i Biskop
Martensens overlogisk-overphysiske Philosophie, ville vi for
en Ordens Skyld endnu engang, maaskee for sidste Gang,
begynde med at undersøge, hvorledes det staaer sig med Phi\-%
losophien, navnlig Logiken, i vor „christelige Dogmatik“. Vi
ville, da Standpunkterne ere saa forskjellige, begynde med et
Blik paa deraf flydende Vanskeligheder.

Det er fra Philosophiens logisk-physiske Synspunkt
naturligviis altid en vanskelig Sag at skulle indlade sig i
videnskabelige Forhandlinger med en Theolog; jo mere ud\-%
mærket Theologen er, desto større er Vanskeligheden. Hører
nu, hvad vi vel tør forudsætte som givet, Biskop Martensen
til Nutidens mest udmærkede Theologer, og er vor opstillede
Forudsætning rigtig, saa følger jo heraf, at philosophiske
Forhandlinger med Biskop Martensen maae høre blandt de
allervanskeligste. Store Vanskeligheder have imidlertid, ifølge
det gamle Ord: \textit{per ardua ad astra}, noget Fristende; heri
søge man Grunden til, at jeg saa ofte, maaskee altfor ofte,
er bleven fristet. For atten Aar siden, da Fristelsen første
Gang besøgte mig, saa at jeg i en, som mig syntes, temmelig
objectiv Anmeldelse vovede at sammenstille \textit{Dr.} Martensens
Dogmatik med S. Kierkegaards Johannes Climacus, blev
jeg ramt af Dommens Ord: „Salige ere de, der ikke ville
paanøde Andre deres Viisdom!“ det er udlagt: Salige ere
de, der ikke ville kritisere „min Dogmatik“! Og jeg Usalige,
der nu til Straf, fordi jeg har været Dommen overhørig, er
kommen saaledes i Fart med at kritisere „min Dogmatik“,
at jeg neppe kan standse igjen: det være et advarende Ex\-%
% --- p. 77 ---
\opage{77}\setcounter{footnote}{0}%
empel for yngre Medphilosopherende, at de betænke sig vel,
før de røre ved et vist Noget, der kaldes „min Dogmatik“;
thi at røre ved „Dogmatiken“ er i Philosophien, hvad det er
i Eventyret at røre ved „Guldgaasen“.

Vi høre Biskop Martensens efter en syttenaarig Over\-%
veielse veloverveiede Erklæring (Leilighedsskrift S. 7): „Men
naar jeg for mit Vedkommende nu optræder imod Prof. R.
Nielsen, da vente man dog ingenlunde, at jeg paa nogen
Maade vil indlade mig paa at forsvare min Dogmatik mod
de af ham rettede Angreb. Naar der er nedlagt en viden\-%
skabelig Protest mod „al“ Theologie (Fortalen til Grund\-%
ideernes Logik II), naar altsaa selve Theologiens Mulighed
under alle Skikkelser benægtes, kan det visselig ikke falde
mig ind at fremkomme med et Forsvar for mit eget Forsøg
paa en Fremstilling af de christelige Troeslærdomme.“

Er dette nu ikke til Skræk og Advarsel? Jeg gaaer
i min Troskyldighed ud fra den Forudsætning, at Biskop
Martensens Fremstilling af de christelige Troeslærdomme
er saa omtrent den bedste videnskabelige Fremstilling, der
i vor Tid kan gives\footnote{Dette er træffende udtalt af en af de yngre Deeltagere i disse
Forhandlinger: „Naar Professor Nielsen i sine Angreb paa Theo\-%
logien stadig har sigtet paa Martensens Dogmatik, saa er dette skeet
af samme Grund, ifølge hvilken man ved et Bombardement sigter
efter Kirketaarnene.“ Rudolf Schmidt i Dansk Kirketidende. 1866.
Nr. 43 Side 679.}, at altsaa det Angreb, der væsent\-%
lig rammer denne Theologie, maa ramme al Theologie; og
see! saa har Biskop Martensen den uforudseelige Beskedenhed
at nedsætte sin Fremstilling til et, maaskee endog underordnet,
Forsøg, og det uden engang saameget som at give mig et
Vink om, hvor jeg da kan finde et betydeligere Forsøg, en
værdigere Gjenstand for min Kritik.

% --- p. 78 ---
\opage{78}\setcounter{footnote}{0}%
Men hvorfor kan det --- lader os høre det endnu en
Gang --- ikke falde Biskoppen ind at fremkomme med et
Forsvar for sin egen, af mig ved saamange og bestemte
Indsigelser angrebne Dogmatik, medens det efter en saa
mangeaarig Taushed dog omsider har kunnet falde ham ind
at udsende et Leilighedsskrift med en Protest imod min Phi\-%
losophie? Grunden, siger han, er den, at „selve Theologiens
Mulighed under alle Skikkelser benægtes“ i mine Angreb.
Dette Argument er typisk theologisk; det har nemlig den
mærkværdige Egenskab, at naar man siger \textit{credo, ut intel\-%
ligam} og seer det i en Totalanskuelse, saa voxer det op til
Størrelse og Vægt af et Bjerg; siger man derimod ikke \textit{credo,
ut intelligam}, men underkaster det en logisk Analyse, kryber
det ind og forsvinder.

Altsaa med \textit{credo, ut intelligam}: Ingen kan med
Eftertryk angribe Theologien uden at angribe den i en vis
bestemt Skikkelse; sin bestemte Skikkelse faaer Theologien hos
de enkelte Theologer i de enkelte theologiske „Forsøg“. Det
enkelte „Forsøg“ er ikke Theologien; Theologien som Theo\-%
logie maa søges i alle „virkelige og mulige Forsøg“ hos alle
virkelige og mulige Theologer. Paa denne Maade er det en
Umulighed at angribe Theologien. Hver Gang en Theologie
er angreben, selv om Angrebet var saa slaaende, at Theo\-%
logerne ikke kunde sige et Ord derimod, saa har man dog
den Trøst, at det kun er en Theologie, ikke Theologien, der
er angreben. Om ogsaa en philosopherende Kritiker havde
en saa overmenneskelig Productivitet, at han hver Dag i et
Tidsrum af halvtredssindstyve Aar kunde gjøre det af med
en enkelt Theologie: han vilde dog langtfra blive færdig
med alle virkelige Theologier; og selv om han blev færdig
med alle de virkelige, saa havde han dog endnu alle de
mulige tilbage. Kan nu ovenikjøbet hver enkelt Theolog,
efter Biskop Martensens Exempel, ignorere ethvert Angreb paa
Theologien i sin Theologie, saalænge den Angribende ikke er
% --- p. 79 ---
\opage{79}\setcounter{footnote}{0}%
bleven færdig med alle andre mulige og virkelige Theologier, saa
sees jo heraf, at Theologerne ligesaa lidt behøve at bryde sig
om Angreb paa de enkelte Theologier som paa Theologien.
Stærkere Position kan ingen Videnskab indtage, at sige, naar
den holder fast ved \textit{credo, ut intelligam}.

Men uden \textit{credo, ut intelligam}, hvad bliver det saa
til med Totalanskuelsen? Den enkelte Theolog maa være
uhyre beskeden, dersom han virkelig skal ansee sit theologiske
Forsøg for saa usselt, at et Angreb godt kan ramme det
uden i nogen Maade at afficere Theologien selv. En saadan
Beskedenhed er unaturlig for en Theolog, unaturlig for en
theologisk Biskop, aldeles unaturlig for en Biskop Martensen.
Jo betydeligere det enkelte theologiske Forsøg er, desto større
Betydning maa ikke blot et Angreb paa, men ogsaa et For\-%
svar for et saadant Forsøg kunne have for Theologien som
Theologie.

Ved ikke at ville imødegaae et endog temmelig udførligt
og vedholdende Angreb paa hans egen Theologie af den
Grund, at Angrebet falder sammen med en Protest imod al
Theologie, har Biskop Martensen neppe beviist enten sig selv
eller Theologien i det Hele nogen stor Tjeneste.

Havde Biskoppen, hvad der jo maatte være ham med
sit „kongelige Princip“ en let Sag, slaaet et Slag til
Forsvar for sit „Forsøg“, da vilde den vundne Seier
aabenbart have kastet Glands, ikke blot paa hans Theo\-%
logie, men paa Theologien i det Hele. Ikke-Theologer,
Saadanne, som allerede have begyndt at vakle i Troen,
vilde da have sagt, enhver til sin Næste: „Der kan man see,
hvad det er for blind Allarm med den idelige Raaben paa,
at Theologien maa bort; man har villet angribe al Theo\-%
logie og see nu, hvor ynkelig man er kommen fra det ved at
angribe Biskop Martensens Theologie; jo, jo, Theologien ---
det tænkte vi nok --- er min Tro dog en Videnskab.“ Nu
derimod, da Biskop Martensen unddrager sig Forsvaret under
% --- p. 80 ---
\opage{80}\setcounter{footnote}{0}%
Paaskud af, at det er al Theologie, der er angreben, vaagner
let en Mistro hos dem, der ikke ere faste i Troen; de ville
sige: „Det var da kjedeligt, at Biskop Martensen, den Eneste
her hjemme, der kan slaae Slaget, har ladet saa god en
Leilighed til en eclatant Seiervinding slippe sig af Hænderne:
det skulde da aldrig være galt fat med den gamle Theo\-%
logie?“ Det er naturligviis Vantro, vi indrømme det, den
pure Vantro; men saadanne ere nu Menneskene: Verden
ligger i det Onde.

Ikke alene for Theologiens Skyld, men ogsaa for sin
egen Skyld havde Biskop Martensen, som mig synes, burdet
indlade sig paa at forsvare sit „Forsøg paa en Fremstilling
af de christelige Troeslærdomme.“ Sæt, at det fra Mod\-%
partens Side virkelig var en Overilelse med denne Nedlæggen
Protest imod al Theologie, fulgte saa uden videre deraf, at
det ogsaa maatte være en Overilelse med at nedlægge Pro\-%
test imod Videnskabeligheden i Biskop Martensens Dogmatik?
Der skal et stærkt \textit{credo} til, naar Nogen skal holde fast ved
en Slutning som denne: at kritisere Biskop Martensens
Dogmatik er at nedlægge Protest imod al Theologie; at
nedlægge Protest imod al Theologie er usaligt: salige ere de,
som ikke ville nedlægge Protest imod al Theologie; salige ere
de, som ikke ville kritisere Biskop Martensens Dogmatik.
Et saadant \textit{credo} vilde jo være uhyggeligt for Dogmatikeren
selv og staae i Strid med hans Beskedenhed.

Men tør et saadant \textit{credo} ikke lægges til Grund for
vore Betragtninger af Sagen: hvordan kan det da falde
Biskoppen ind at benytte den Omstændighed, at der i For\-%
talen til Grundideernes Logik II er nedlagt Protest imod al
Theologie, som Paaskud til at trække sig tilbage fra at for\-%
svare sig imod den Protest, der er nedlagt imod hans egen
Theologie? Her maa stikke Noget under.

Jeg har det. Grunden, hvorfor Biskop Martensen
ikke vil forsvare sig imod mit Angreb, er, hvad jeg jo
% --- p. 81 ---
\opage{81}\setcounter{footnote}{0}%
nok kunde formode, slet ikke den, at jeg har nedlagt Pro\-%
test imod al Theologie; nei Grunden er en ganske anden.
I Leilighedsskriftet (S. 106) er Grunden angiven; Grunden
er, at man i Grunden ikke kan indlade sig med mig. „Det
hører til Grundpræget i Prof. Nielsens Polemik, at han
idelig tillægger sine Modstandere Intensioner, som de ikke
vedkjende sig, og som han er uberettiget til at paabyrde dem;
og det er denne hans Polemiks Charakteer, der umuliggjør
enhver Forhandling med ham.“

Her have vi da en ordentlig, og hvis Sagen forholder
sig som Biskop Martensen siger, tilstrækkelig Grund, en
Grund, der vel kan forklare Læseverdenen, hvorfor Biskoppen
ikke ønsker at forsøge noget Forsvar. Men det kunde
han jo saa have sagt os lige strax. I en videnskabe\-%
lig Forhandling om Tro og Viden seer det da ganske be\-%
synderligt, næsten fruentimmeragtigt ud, først at fixere os
med en Grund, der ikke er en Grund, for siden, naar
vi mindst vente det, at overraske os med den rette og egent\-%
lige Grund. Altsaa --- \textit{hic Rhodus hic salta!} --- min
Polemiks Charakteer umuliggjør enhver Forhandling med mig;
den Sag maa undersøges.

Heldigviis kan Undersøgelsen skee ganske objectivt, efter\-%
som det under de forhaandenværende Omstændigheder ikke vil
være vanskeligt at fremlægge saamange Data, at den tænkende
Læser selv kan følge med og see, hvor vi ere.

Det er ganske rigtigt, hvad Biskop Martensen anfører,
at jeg i mit Svar til \textit{Dr.} Zeuthen har kaldt Theologien
en Mirakelvidenskab, en Videnskab, der forklarer Mirakler.
Hvad har nu Biskoppen at indvende herimod? Hans Af\-%
viisning er høist besynderlig: „Theologien“, siger han, „er
Aabenbaringsvidenskab, og til Aabenbaringen hører det gud\-%
dommelige Under. Men ingenlunde er Theologien en til\-%
fældig Samling af Mirakler. Heller ikke giver den sig af
med Mirakelforklaringer.“ Den Forestilling, der især synes
% --- p. 82 ---
\opage{82}\setcounter{footnote}{0}%
at vække Biskoppens Uvillie, er Forestillingen om, at der til
Mirakelforklaringer absolut skulde fordres, at Theologien blev
en Videnskab om „en tilfældig Samling af Mirakler“, og at
det saa kom an paa at „forklare selve den guddommelige
Virkemaade eller Underets Hvorledes“. Men herved tillægger
han mig „Forudsætninger og Intensioner“, som jeg ikke ved\-%
kjender mig, og som han er uberettiget til at paabyrde mig.
Thi jeg har virkelig ingensteds opstillet den urimelige For\-%
dring, at en Mirakelforklaring skulde bestaae i at forklare
„selve den guddommelige Virkemaade“, endnu mindre kunde
det falde mig ind, at Theologien kun blev Mirakelvidenskab
ved at blive Videnskab om „en tilfældig Samling af Mi\-%
rakler“. Jeg har slet og ret paastaaet, at Theologien, for
at være Videnskab, maa være en Mirakelvidenskab, en Viden\-%
skab, der forklarer Miraklerne, d. v. s. faaer noget Viden\-%
skabeligt ud af Miraklerne.

Hvad vil da Biskop Martensen have, at her skal staae
istedenfor Mirakelvidenskab? Aabenbaringsvidenskab! „Theo\-%
logien er Aabenbaringsvidenskab, og til Aabenbaringen hører
det guddommelige Under.“ Hvor charakteristisk! Naar Theo\-%
logen definerer sin Videnskab, bliver Definitionen som af sig
selv en Tvetydighed, en Gaade; thi en Sætning som den:
til Aabenbaringen hører det guddommelige Under, er jo tve\-%
tydig som en mørk Tale, gaadefuld som et Orakelsvar. Til
Aabenbaringen hører det guddommelige Under: hvad vil det
sige? Skulde Meningen maaskee være den, at Aabenbarin\-% p. 82: a faint raised dot between "den," and "at" read as a speck, not a sort
gen er Et og Underet (ved Underet tør vi dog vel forstaae
Miraklet) et Andet; ja det var jo ganske snildt: paa denne
Maade kunde Theologien blive en Videnskab om Aaben\-%
baringen uden at blive en Videnskab om Miraklet, en Aaben\-%
baringsvidenskab, der ikke var Mirakelvidenskab.

Men hvordan skulle vi ulykkelige Mennesker, der nu
saa gjerne ville forstaae, hvad Biskop Martensen forstaaer
ved Theologie, komme ud af det med at tænke os en Aaben\-%
baring, der ikke er Mirakel, og et Mirakel, nemlig „et
% --- p. 83 ---
\opage{83}\setcounter{footnote}{0}%
guddommeligt“ Mirakel, der vel hører til Aabenbaringen,
men ikke gaaer i Eet med Aabenbaringen? Thi, jeg vil op\-%
rigtig tilstaae det: jeg havde hidtil i min utheologiske En\-%
foldighed bildt mig ind, at naar der taltes om Under, om
Underet, særdeles og med Eftertryk om „det guddommelige
Under“, saa var Underet Aabenbaring, og Aabenbaringen
Under. Men det gaaer ikke an: til Aabenbaringen hører det
guddommelige Under, det vil sige: til Aabenbaringen hører
Aabenbaringen; nei, ved Hermeneutiken, det gaaer ikke an.
Der er da intet Andet for end at forstaae det efter Bogstaven:
Theologien er en Aabenbaringsvidenskab, men ingen Mirakel\-%
videnskab, d. v. s.: Theologien kan ikke faae noget Viden\-%
skabeligt ud af Miraklerne.

Men naar Theologien omsider har besluttet sig til at
lade Miraklerne, ja Miraklet selv, staae udenfor Viden\-%
skaben, kan den saa nok faae Aabenbaringen ind i Viden\-%
skaben? Nei, det kan da heller ikke være Meningen. Hvad
er da Meningen? Jeg kan ikke finde Meningen; jeg for\-%
klarer og forklarer, men jo længere jeg forklarer, desto ufor\-%
klarligere bliver det Hele. Jeg maa begynde fra hvilken
Ende jeg vil, det ender bestandig paa samme Maade; det
ender med Rahbeks Omkvæd: „Enten man ta'er det paa
langs eller tvers, Taabeligt er det som Lüttikaus Vers.“

Men, tænker den billige Læser, man bør da heller ikke
hænge sig i et enkelt Ord, en enkelt Sætning, men smukt
forklare det Enkelte af det Hele og alvorlig bestræbe sig for
at forstaae Biskop Martensen udaf hans „Totalanskuelse“.
Velan! vi ville forsøge det med Totalanskuelsen. „Jeg ved\-%
bliver“ --- dette er jo Biskop Martensens Totalanskuelse
(Dogm. Oplysning. S. 55) --- „fremdeles at betragte Dog\-%
matikerens Forhold til Aabenbaringen paa tilsvarende Maade
som Naturforskerens til Naturen.“ Altsaa: saalidt som
Naturforskeren kan adskille Naturen fra Naturens Kjends\-%
gjerninger, ligesaa lidt kan Dogmatikeren adskille Aabenba\-%
% --- p. 84 ---
\opage{84}\setcounter{footnote}{0}%
ringen fra Aabenbaringens Kjendsgjerninger. Men til Aa\-%
benbaringens Kjendsgjerninger høre jo Miraklerne. Altsaa:
ligerviis som Naturforskerens Forklaring af Naturen falder
sammen med Forklaringen af dens Kjendsgjerninger, saaledes
falder ogsaa Dogmatikerens Forklaring af Aabenbaringen
sammen med hans Forklaring af Aabenbaringens Kjends\-%
gjerninger ɔ: af Miraklerne. Totalanskuelsen taler, som man
seer, udtrykkelig for, at Theologien maa være en Mirakel\-%
videnskab, en Videnskab, som dog paa en eller anden Maade
giver sig af med at forklare Mirakler.

Men mod Totalanskuelsen strider Specialanskuelsen:
„heller ikke giver den (Theologien) sig af med Mirakelfor\-%
klaringer“, paa det Allerbestemteste.

Den eneste Udvei, her er tilbage, er at undersøge, ikke
de forskjellige Betydninger af Ordene: forklare og Forklaring,
thi det fører ikke til Noget, men den Behandling af Mirak\-%
lerne, som vor Dogmatiker selv har anvendt, idet han i Gjer\-%
ningen har viist os, hvad der skal forstaaes ved de gaade\-%
fulde Ord: „Theologien er en Aabenbaringsvidenskab, og til
Aabenbaringen hører det guddommelige Under.“

For ikke at tillægge mine Modstandere „Forudsætninger,
som de ikke vedkjende sig“, har jeg i mit Svar til \textit{Dr.}
Zeuthen udtrykkelig henviist til „den christelige Dogmatik“.
Mine Ord (S. 6) ere disse: Hvad har da vor „christelige
Dogmatik“ bragt frem til Forstaaelse af Aabenbaringens tre
Grundundere: \emph{Skabelsen}, \emph{Incarnationen} og \emph{In\-%
spirationen}? Med faa Træk og i saa skarpe Omrids
som muligt har jeg her søgt at vise, hvad Dogmatiken med
sine Forklaringer har opnaaet. Den Indvending, at man
ikke kan tage Hensyn til en Charakteristik, der er saa kort,
at den paa fire Octavsider omfatter tre Grundundere, altsaa
i Grunden den hele Dogmatik, gjælder ikke i Videnskaben;
thi Videnskaben maa ligesaa fuldt lægge an paa at reducere
som paa at deducere.

% --- p. 85 ---
\opage{85}\setcounter{footnote}{0}%
Havde Biskop Martensen, hvad den angivne Reduction
angaaer, virkelig havt Noget at klage over, Noget, der i nogen
Maade kunde berettige ham til en saa personlig nærgaaende
Yttring som den, at jeg tillægger mine Modstandere „For\-%
udsætninger og Intensioner, som de ikke vedkjende sig“, og at
det er denne min Polemiks Charakteer, der „umuliggjør enhver
Forhandling“ med mig: da havde han her den bedste Lei\-%
lighed til at gribe mig paa fersk Gjerning. Jeg, der havde
den Formastelighed at ville ved en comprimerende Reduction
knuse „den christelige Dogmatik“ og med den „al Theologie“,
var da selv bleven knust. Thi jeg skal ikke fragaae, at der\-%
som det virkelig kunde bevises, at Reductionen indeholdt
enten Fordreielser eller væsentlige Misforstaaelser, vilde Theo\-%
logerne jo saa smaat kunne begynde at illuminere.

Hvorfor har Biskop Martensen ikke taget fat paa
Reductionen? Hvorfor har han ladet sig nøie med bare
Beskyldninger, hvor det dog vist havde været ham let at
skaffe en Smule Beviis tilveie? Istedenfor at vove sig
med en fremmed Logik ud paa Ideernes aabne Hav
imellem krydsende Reflexioners Malstrømme, hvor en lille
Fiskerbaad saa let kan kæntre, var det jo dog langt sikkrere
at holde sig i de velbekjendte Farvande, blive ved det dog\-%
matiske Kystland og see at værne lidt om sit Eget. At
Biskop Martensen endnu har havt Leilighed til at læse
Grundideernes Logik, tør vi, efter hvad vi kunne skjønne af
hans Citater i Leilighedsskriftet, vistnok ikke forudsætte; der\-%
imod tør vi nok forudsætte, at han har havt Leilighed til at
læse sin egen Dogmatik, ligesom vi heller ikke troe at gjøre for
store Fordringer, naar vi antage, at han ved Siden af at
gjennemgaae sin Dogmatik kunde have fundet Leilighed til at
gjennemgaae de fire Sider med Reductionen i mit Svar til
\textit{Dr.} Zeuthen.

Det vilde meget have fremmet Forhandlingerne, om
Biskop Martensen paa en eller anden Maade, ved et eller
% --- p. 86 ---
\opage{86}\setcounter{footnote}{0}%
andet slaaende Exempel havde viist mig, hvori min Mis\-%
kjendelse, af hvad jeg maa kalde hans dogmatiske Mirakel\-%
forklaringer, da egentlig bestaaer.

Er her i den „christelige Dogmatik“ maaskee Intet, der
kunde pege hen paa, at „Aabenbaringsvidenskaben“ alligevel,
hvor den seer sit Snit, praktiserer saa smaat i Mirakelfor\-%
klaringer? For til en Forandring at begynde med det tredie
Grundmirakel er her Pintsemiraklet. Tør „den christelige
Dogmatik“ nu række sine tre Fingre iveiret og sværge, at
den veed sig fri, at den, hvad Pintsemiraklet angaaer, ikke
har givet sig af med Noget, der kunde ligne en Mirakel\-%
forklaring? tør den sige „ved Gud“ paa det? Jeg troer
det ikke. Der staaer jo dog virkelig i Dogmatiken --- hvis
Biskop Martensen ikke saa nøie husker det, kan han selv slaae
det efter S. 347 § 186. „Hiin Talen i Tunger var alt\-%
saa Udtrykket for den første, endnu regelløst frembrydende og
ligesom overstrømmende Begeistring, i hvilken Apostlene følte
hele deres naturlige Tilværelse gjennemzittrede af Guds Aand
som af et Lyn, følte sig overvældede af den overvættes Kraft“
o. s. v. Her er jo uden Omsvøb en videnskabelig Mirakel\-%
forklaring, og det ovenikjøbet en Forklaring, som, naar
Ordene „regelløst frembrydende og ligesom overstrømmende“
behørigt accentueres, ikke er langt fra at forklare „selve den
guddommelige Virkemaade eller Underets Hvorledes.“ Ja,
hvor alvorlig det er meent med Underets Hvorledes, kan
sees deraf, at Dogmatiken, for at kunne forklare „selve den
guddommelige Virkemaade eller Underets Hvorledes“, ikke tager
i Betænkning at endevende hele Miraklet; thi der staaer med
rene Ord (S. 137 § 186): „Naar derimod de andre Til\-% p. 86: "S. 137" as printed (cf. "S. 347 § 186" above)
stedeværende hørte dem tale om Guds store Gjerninger Hver
i sit Tungemaal, da er det“ --- thi at Lukas Apst. Gjerng.
2, 3, udtrykkelig siger, at Apostlene begyndte at tale med
\emph{andre Tungemaal} (ἑτέραις γλώσσαις) kommer her, hvor
der spørges om det høiere theoretiske \emph{Hvorledes}, ikke prak\-%
% --- p. 87 ---
\opage{87}\setcounter{footnote}{0}%
tisk i Betragtning --- „ikke fornødent her at tænke paa for\-%
skjellige Sprog, men derpaa, at de Hver for sig følte sig
tiltalte i deres egen inderste \emph{Eiendommelighed}, at det
indvortes Menneske saaledes“ --- sig saa, at her ikke er et
\emph{Hvorledes} --- „tiltaltes af dette Henrykkelsens Tungemaal,
at Sjælen“ --- her sees jo „selve den guddommelige Virke\-%
maade“ --- følte sig udløst af den vante, den naturlige Ind\-% p. 87: printer's defect? no „ reopens the quotation before "følte" (its close “ follows "aabenbar."); dash printed short, as "--"
skrænkning og paa underfuld Maade blev sig selv aabenbar.“

Læseren vil heraf see, at Biskop Martensen, idet han
vil være uenig med mig, bliver uenig med sin egen Dog\-%
matik. Giver Theologien sig af med Mirakelforklaringer?
Biskop Martensen svarer med et bestemt Nei! Men Dog\-%
matiken og jeg svare med et ligesaa bestemt Ja! Mod mig
kan Biskop Martensen jo sagtens staae sig; men kan han
nu ogsaa staae sig imod sin egen Dogmatik?

Man vil maaskee for at mægle Forlig imellem Biskoppen
og hans Dogmatik, lægge en ganske overordentlig Vægt paa
den Omstændighed, at Forklaringens \emph{Hvorledes} ikke egent\-%
lig er metaphysisk, men blot psychologisk; man vil angaaende
de psychologiske Forklaringer maaskee anvende paa „den chri\-%
stelige Dogmatik“, hvad en troskyldig Examenscandidat, der
blev spurgt, hvorfor der i hans Læsebog stod \textit{mon amie} og
ikke \textit{ma amie}, engang skal have sagt om den franske Gram\-%
matik: „de Franske tage det ikke saa nøie“; men det er en
Forhaanelse imod Dogmatiken. Til Beviis for, at Dogma\-%
tiken tager det nøie selv med Miraklets metaphysiske Hvor\-%
ledes, anføres ordret, hvad der om det andet Grundmirakel,
Incarnationen, er sagt i Svaret til \textit{Dr.} Zeuthen.

„Naar man“, siger Dogmatiken, „har fundet det selvmod\-%
sigende, at den evige Logos bliver Menneske „fordi den Evige og
Allestedsnærværende ikke kan lade sig føde midt i Tiden“, da kan
det vistnok ikke være Opgaven at tænke, at den evige Logos
med Incarnationen skulde ophøre at existere i sin almindelige
Verdensaabenbaring, eller at Logos som selvbevidst personligt
% --- p. 88 ---
\opage{88}\setcounter{footnote}{0}%
Væsen skulde være indesluttet i Modersliv, fødes som Barn,
tage til i Kundskab o. s. v., thi denne Forestilling ophæver
selve Fødselens Begreb. Fødsel i Tid medfører nødvendigt
Fremgang fra det Ubevidste til det Bevidste, fra Mulighed
til Virkelighed, fra Sædekorn og Spire til fuldbaaren Or\-%
ganisation. At den guddommelige Logos lader sig føde, vil
derfor sige, at han indsænker sig i Menneskehedens Skjød
som Mulighed, som hellig Sæd, for at kunne opstige i
Menneskeheden i midlende og frelsende Menneskeaabenbarelse“.
Her have vi altsaa Forklaringen. Men skulde vi ikke let
blive enige om, at denne Forklaring selv trænger til en
Forklaring? Thi hvad er det den guddommelige Logos
egentlig har „indsænket“ i „Menneskehedens Skjød“, sig
selv eller ikke \emph{sig selv}? Er det ikke sit egentlige Selv,
saa er det jo heller ikke Logos selv, der er incarneret. Det
Væsen, hvis egentlige Selv ikke kan „indsænkes“ som Mulig\-%
hed i Modersliv, kan ikke incarneres. Men er det virkelig
sig selv, sit mulige, væsentlige Selv, men dog ikke det fra
Evighed af selvbevidste personlige Selv, hvorom her tales:
saa maa den guddommelige Logos jo have et dobbelt Selv,
et, der kan incarneres, og et, der ikke kan incarneres. Den
Logos, der kan incarneres, er altsaa i sig selv, i sit egent\-%
lige Selv, en Anden end den evige personlige Logos selv:
med andre Ord: den evige personlige Logos selv er ifølge
„den christelige Dogmatiks“ Forklaring slet ikke bleven in\-%
carneret.“% p. 88: this “ closes the quotation from the Svar til Dr. Zeuthen, which opened on p. 87 with the same „ as „Naar man“ (no separate opening mark printed)

Da Biskop Martensen, hvad Citatet af hans Dogmatik
(S. 269) angaaer, ikke har paaviist og ikke nogensinde vil
kunne paavise nogensomhelst enten frivillig eller ufrivillig
Forvanskning: kan her naturligviis ikke reise sig Strid om
Texten selv, men alene om dens Fortolkning. Fra hvilket
Standpunkt vil Dogmatikeren her have sin dogmatiske For\-%
klaring opfattet? Seet objectivt, ɔ: under Synspunktet af
det Logisk-Physiske, indeholder Forklaringen den himmel\-%
% --- p. 89 ---
\opage{89}\setcounter{footnote}{0}%
skrigende Modsigelse, at det Selv, som i Mennesket Jesus
er født i Tiden af Jomfru Maria, og det evige Logos-%
Selv skal være baade det samme og dog ikke det samme
Selv. Denne Modsigelse lader sig løse paa to aldeles for\-%
skjellige Maader, eftersom den løses enten i Viden eller i Troen.

Skal Modsigelsen løses i Viden, da maae alle Tvetydig\-%
heder bort; dette er den første Fordring. Til de iøinefal\-%
dende Tvetydigheder høre Sætninger som denne: „At den
guddommelige Logos lader sig føde, vil sige, at han ned\-%
sænker sig i Menneskehedens Skjød som Mulighed, som hellig
Sæd“ o. s. v. Men naar alle Tvetydigheder hæves, hvad
bliver der saa af Grundmiraklet? En Grundmodsigelse, som
under den videnskabelige Analyse udvikler sig til en Vrimmel
af uopløselige Modsigelser! Skulle de mange i Videnskaben
utaalelige Modsigelser fjernes, da maa Grundmodsigelsen selv
hæves; men seet i Videns Lys er Grundmodsigelsen Eet med
Grundmiraklet. Idet Videnskaben skiller sig af med Mod\-%
sigelsen, skiller den sig af med Miraklet; Mirakelforklaringen
ender i Mirakelfornegtelse: dette er Knudens videnskabelige
Løsning.

Skal Modsigelsen løses i Troen, da bliver det ingen\-%
lunde Opgaven at bestemme, hvorledes Incarnationens Meta\-%
physik skal enten tænkes eller ikke tænkes, thi Troesprincipet
hæver sig kun over det Logisk-Physiske derved, at det ude\-%
lukker al objectiv-\emph{logisk-physisk} og da med det Samme% p. 89: only "logisk-physisk" is letterspaced, not "objectiv-"
al \emph{metaphysisk} Tænkning. I Christi Person forudsættes
Eenheden, den oprindelige, ved Undfangelsen givne Eenhed
af guddommelig og menneskelig Natur, guddommelig og men\-%
neskelig Villie. I Lys af denne Eenhed opfatter og bestemmer
den religiøse Tænkning Alt, hvad Evangeliet lærer os om
Jesu Liv, men Eenheden selv gjennemtænker den ikke: Een\-%
heden er og bliver et \emph{absolut Mysterium}. Det Første,
den religiøse Tænkning har at gjøre, er at skille sig af med
alle dogmatiske Incarnationsproblemer. For den troende
% --- p. 90 ---
\opage{90}\setcounter{footnote}{0}%
Bevidsthed er Grundmiraklets Modsigelse ganske vist løst;
men ikke derved, at den selv \emph{tænker} Løsningen, nei derved,
at den \emph{troer}, at Løsningen er tænkt i Mysteriet.

Dersom vi turde tage Biskop Martensen paa Ordet,
naar han afviser alle Mirakelforklaringer, saa turde vi jo
ogsaa gjøre Regning paa hans Bistand, naar vi alvorlig
lægge an paa at afvise „den christelige Dogmatiks“ halv\-%
videnskabelige Qvaklerier med Incarnationsproblemets Meta\-%
physik. Incarnationsproblemet er og bliver et Mirakelpro\-%
blem; et Mirakelproblem kan ikke løses, end ikke til en vis
Grad, uden at man, idetmindste til en vis Grad, giver sig
af med en Mirakelforklaring. Er det nu virkelig Biskop
Martensens Alvor, og en saa alvorlig Sag maa jo tages
alvorligt, at nedlægge Protest imod alle Mirakelforklaringer
og derved mod al Mirakelvidenskab: da ere vi ganske enige,
og naar vi To blot ere enige i dette ene Punkt, skulle vi,
som jeg haaber, nok faae Bugt med Dogmatiken. Hvad vil
nemlig Dogmatiken med sine Mirakelforklaringer? Den vil
--- lader os kun sige det, for at bevise det --- den vil, reent
ud sagt, \emph{vrøvle}. Idet vi da, stolende paa et biskoppeligt
Ord: ingen Mirakelforklaringer! agte at gaae „Dogmatiken“
paa Klingen og ligefrem sigte den for \emph{Vrøvl}, maae vi dog
strax bemærke, at vi derfor ingenlunde glemme, at der er
Forskjel imellem dogmatisk Vrøvl og andet Vrøvl. For at
adskille dogmatisk Vrøvl fra andet Vrøvl, ville vi om det
dogmatiske ikke bruge Ordet Vrøvl, men hellere indføre en
ny Kategorie: det \emph{Phlyariske}, eller, som vi ogsaa kunde
kalde det, det Phlyaretiske, men ikke Phlyaristiske (af
φλυαρέω); og saa spørge: \emph{hvori kan den „christelige
Dogmatiks“ Phlyariske eller Phlyaretiske som
saadant nu egentlig siges at have sin Bestaaen}?
Det kan siges at have sin Bestaaen i en tilsyneladende dyb\-%
sindig, men i Grunden selvmodsigende, afsindig og menings\-%
løs Forening af Tro og Viden. Det Phlyariske i Foreningen
% --- p. 91 ---
\opage{91}\setcounter{footnote}{0}%
træder frem og lægger sig for Dagen i en stor Mængde saa\-%
vel for Videnskaben som for Religionen usunde, yderst usunde
\emph{Mirakelforklaringer}. Blandt Exempler paa de mange
Forklaringer fortjener Forklaringen af Incarnationens Mi\-%
rakel en særegen Plads. Til nærmere Forstaaelse sættes det
Phlyariske først i sit \emph{negative} og dernæst i sit \emph{positive}
Moment; Mediationen af det negative og positive Moment
giver os saa tilsidst den speculative Eenhed.

1) \emph{Det Phlyariske i sit negative Moment.}
Troen forudsættes, og saa gaaer det løs paa Viden: \textit{credo,
ut intelligam}. Det Eiendommelige i den dogmatiske Viden
maa naturligviis bestemmes ved det Eiendommelige i den
dogmatiske Tænkning. Den dogmatiske Tænkning er meget
eiendommelig; den charakteriserer sig fra først til sidst der\-%
ved, at den efter at have forkastet det Tænkelige ender med
at tænke det Utænkelige. Incarnationsmetaphysiken afgiver
et fortræffeligt Paradigme. Incarnationen maa nemlig ikke
tænkes saaledes, at den evige Logos har sænket sig selv, sit
evige, uforanderlige, altid bevidste Selv i Menneskehedens
Skjød: „At Guds Søn ikke som selvbevidst Guddomsjeg,
men som ufuldbaaret Foster er i Modersliv, antyder ogsaa
Skriften“ o. s. v. (Christl. Dogm. § 132). Men Incarna\-%
tionen maa heller ikke tænkes saaledes, at Selvet, det Selv,
der som ufuldbaaret Foster begynder i Modersliv, ikke skulde
være det evige Logos-Selv. „Skjøndt Christus vidner: Jeg
og Faderen ere Eet, siger han dog aldrig, jeg og Logos ere
Eet. Thi han er den menneskelige \emph{Selv}aabenbaring af den% p. 91: only "Selv" is letterspaced
guddommelige Logos.“ (Dogm. § 132). At det Phlyariske
her er sat i sit negative Moment, er kjendeligt paa Resul\-%
tatet, thi Resultatet er complet meningsløst. Det sidste Glimt
af mulig Mening svandt med den Bemærkning, at Christus,
der vidner: Jeg og Faderen ere Eet, aldrig siger: Jeg og
Logos ere Eet. Havde Christus blot een eneste Gang i det
johanneiske Sprog sagt: Jeg og Logos ere Eet, saa havde
% --- p. 92 ---
\opage{92}\setcounter{footnote}{0}%
Incarnationsmetaphysiken dog havt en fjern Udsigt til en
Slags Mening. Man kunde da idetmindste have forsøgt at
tænke sig en Identitet af det Selv, der undfanges og dannes
i Modersliv, med det evige Logos-Selv, der svarede til
Faderens Identitet med Sønnen, og paa denne Maade presse
en halv Tanke ind i Dogmatikens Ord: „Men det Hellige,
der fødes af Maria, og som under sin Fremvæxt bliver sig
bevidst som et menneskeligt Jeg, bliver sig i samme Maal sin
Guddom bevidst; veed sig som guddommeligt Jeg, fordi Gud\-%
domsfylden er \emph{Livsgrunden} til hans menneskelige Leven;
veed sig ikke blot som deelagtig i den guddommelige Logos,
men som den gudmenneskelige Fortsættelse af det evige Gud\-%
domsliv fra Begyndelsen“. (Dogm. § 132). Denne sidste
Udvei er nu aldeles spærret. Christus kan jo ikke sige: Jeg
og Logos ere Eet. Kan han virkelig ikke? Nei, det veed
den store Dogmatik, han ikke kan! Og hvorfor ikke? Fordi
Christi menneskelige Jeg, det Jeg, der har en Begyndelse
i Tiden, har udviklet sig fra et Ubevidst til et Bevidst ---
og det selvbevidste Guddomsjeg, det evige Logos-Jeg, det
Jeg, der ikke har Begyndelse i Tiden, ikke har udviklet sig
fra et Ubevidst til et Bevidst --- ikke ere to identiske Jeger,
men bogstavelig eet og samme Jeg. Yderligere kan dog vel
ingen Tænkning gaae --- i det phlyaretisk Phlyariske.

2) \emph{Det Phlyariske i sit positive Moment.} Det
Eiendommelige ved den dogmatiske Tænkning er, at den
tænker videnskabeligt i Troen: \textit{credo, ut intelligam}. At
der videnskabeligt tænkes i Troen, vil sige, at der ikke tænkes,
men forestilles, anskues, phantaseres paa metaphysisk Maade
til Opklaring af Mysterierne, til Bedste for Aabenbaringen,
til Glæde for Troen. Dogmatiken, der objectivt bygger paa
Skriftens Ord, seer i Skriftstederne noget ganske Andet og
Dybere, end enten Christus selv eller hans Apostle nogensinde
have seet: den seer det Metaphysiske. „Han vilde“, hedder
det (Dogm. S. 133) om Christi Selvfornedrelse, „ikke blot% p. 92: "S. 133" as printed (for "§"?); the KB copy has a hand stroke over the S and a "§" in pencil/ink below
% --- p. 93 ---
\opage{93}\setcounter{footnote}{0}%
leve i den rene Guddomsherlighed, i metaphysisk Majestæt;
men han udtømte sin Fylde i den ringe Tjenerskikkelse for at
kunne blive Hovedet i Kjærlighedens Rige, Forsoneren og
Forløseren. Men denne Selvfornedrelse maa tillige“ --- her
kommer den rette positive Mirakelforklaring --- „sees som
hans Selvfuldendelse“, Selvfuldendelsen nemlig af den i sig
evig fuldendte Logos. Idet han, den ene og samme Logos,
„lever paa menneskelig Viis, og som Menneskens Søn kun
besidder sin Guddom i den menneskelige Individualitets Ind\-%
skrænkning, i den menneskelige Bevidstheds begrændsede Former:
maa han vistnok siges at leve i Fornedrelse og Fattigdom,
efterdi han har givet Afkald paa den majestætiske Herlighed,
i hvilken han som den allestedsnærværende Logos gjennem\-%
lyser Alskabningen.“ (Dogm. § 133). Lægge vi nu Mærke
til de to Udtryk: „han har givet Afkald“, og „han gjennem\-%
lyser“, og erindre, at det er den samme, den selvsamme Lo\-%
gos, det ene og selvsamme Logos-Jeg, der her har givet Af\-%
kald paa den majestætiske Herlighed, hvorpaa han dog ikke et
eneste Øieblik har givet Afkald, hvorledes skulle vi da kunne
tvivle om, at det netop er det høiere Phlyariske, det dybere
Phlyaretiske, som ved denne Forklaring bliver sat i sit posi\-%
tive Moment.

3) \emph{Det Phlyariske i speculativ Fuldendelse.}
Mediationen af Tro og Viden lader sig naturligviis kun
gjennemføre paa den Betingelse, at begge Parterne nærme sig
hinanden med al mulig diplomatisk Forekommenhed d. v. s. i% p. 93: a raised dot between "diplomatisk" and "Forekommenhed" read as a speck (word space both sides)
al Oprigtighed, Uegennyttighed, kun besjælede af de varmeste
Ønsker om Fred og god Forstaaelse. Under den gjensidige
Tilnærmelse foregaaer der en lykkelig Forandring med dem
begge: Troen bliver gemytlig, Viden godmodig. At Troen
bliver gemytlig, vil sige, at den opgiver sin gammeldags or\-%
thodoxe Prikkenhed, undseer sig ved hvad der er Jøder en
Forargelse og Grækere en Daarskab, smiler til Philosophien
og forelsker sig i Speculationen. Det kan naturligviis ikke
% --- p. 94 ---
\opage{94}\setcounter{footnote}{0}%
være nogen almindelig Tro, ikke den folkelige. enfoldige Asyl-% p. 94: printer's defect? a full stop (or a comma that has lost its tail) after "folkelige"; transcribed as printed
Tro, der teer sig saa charmant; nei, det er en dannet, vel\-%
opdragen Tro, en Tro, der kjender sine Mystikere og har hørt
Forelæsninger i Institutet; kun en saadan Tro, en subjectiv
Tro med et objectivt Hjerte, en Tro, der ikke vilde elske
praktisk, dersom den ikke elskede theoretisk, kan aabne sin Favn,
sin dybe, dogmatiske Favn, for den høiere Viden. At den
Viden, der uden Qvalme skal kunne give sig hen i en saa
dannet, saa institutdannet Kjærlighed, absolut maa være \emph{god\-%
modig}, er vel af sig selv indlysende. Her er nemlig en
Forskjel, en meget betydelig Forskjel imellem den godmodige
og den ikke godmodige Viden. Den ikke godmodige Viden
kunne vi med Grund kalde den drillende, nemlig den dog\-%
matisk drillende, fordi den altid har sin Glæde af at drille
Dogmatiken. Drilleriet bestaaer ganske simpelt deri, at den
med ubarmhjertig Conseqvens forfølger enhver Modsigelse,
hvormed Dogmatiken i sin Troskyldighed har indladt sig. Af
et Par uskyldige: „Maa-ikke-tænkes“, fremmaner den ved
Hjælp af det Logisk-Physiske en dæmonisk Vrimmel af „Maa-%
ikke-tænkes“; den faaer ved kritisk dialektiske Hexekunster de
dogmatiske Modsigelser til, om vi saa maae sige, at yngle.
Ja, hvad der er det mest Oprørende, just det dybeste Posi\-%
tive, det, hvorom Dogmatiken udtrykkelig siger, at \emph{det maa
tænkes}, forvandler den til et uroligt Negativt, et i Bund
og Grund Utænkeligt, et „Kan-ikke-tænkes“: det er ondskabs\-%
fuldt. Den ondskabsfulde Viden, den Viden, der ikke vil
slippe det Logisk-Physiske, er en syndig Viden, en Viden, der
ikke kan sige \textit{credo, ut intelligam}, en \emph{vantro} Viden. An\-%
derledes med den godmodige Viden: salige ere de, der ikke
ville kritisere den godmodige Viden. Den godmodige Viden
seer vel undertiden noget tankeløs ud; men, ved Dogma\-%
tiken, den er ikke tankeløs, thi den gaaer i Tanker. Hvor\-%
for gaaer den i Tanker? Den er forelsket, subjectiv forelsket
i Troens objective Hjerte: det er jo en troende Viden. Den
% --- p. 95 ---
\opage{95}\setcounter{footnote}{0}%
troende Viden er en Viden, der længes --- ud over det
Logiske ind i det Overlogiske: Mirakelmetaphysik er dens
Længsels Maal. Naaer den ikke sin Længsels Maal, saa
døer den; men naaer den sin Længsels Maal, saa lever den,
lever af Troens Positive, blæser ad Kritikens Negative og
neddysser alle Modsigelser i absolut Glæde over de relative
Mysterier. Skulde en eller anden Modsigelse være saa
ustyrlig at skrige ud fra Mysteriernes Rede, hvad gjør saa
det? „De Franske tage det ikke saa nøie“. Den troende
Viden er godmodig: godmodig synker den til Dogmatikens
Bryst, godmodig hviler den ved Troens Hjerte. Praktisk er
Godmodigheden Kjærlighed, theoretisk er Kjærligheden God\-%
modighed. Maa det i det Praktiske siges, at Kjærlighed
skjuler Syndernes Mangfoldighed, da maa det her, hvor det
dybeste Praktiske tillige er det høieste Theoretiske, siges, at
Godmodighed, den troende Videns utrolige Godmodighed,
skjuler Modsigelsernes Mangfoldighed. Paa disse Vilkaar
er Foreningen af Tro og Viden fuldbragt; det Phlyariske er
sat i speculativ Fuldendelse.

Ved dette Resultat tør jeg vel mene at have Biskop
Martensen med mig; jeg stoler paa hans Ord, at Dogma\-%
tiken maa høre op med sine Mirakelforklaringer; jeg er vis
paa, at han vil vedkjende sig alle af hans Forudsætninger
flydende Conseqvenser, saa vis, som om jeg allerede havde
dem sikkrede ved biskoppelig Confirmation. Kun een Skrupel
er endnu tilbage, den nemlig, at jeg ikke har forstaaet hans
Ord aldeles correct. Thi det lader sig jo tænke, at Biskop
Martensen ikke vil have sin Afviisning af de theologiske Mi\-%
rakelforklaringer forstaaet saaledes, at Theologien ligefrem
giver Afkald paa de Mirakelforklaringer, den virkelig opgiver, men
saaledes, at den i sin Fornedrelse og Fattigdom „giver Afkald paa“
de Mirakelforklaringer, hvormed den i sin „metaphysiske Her\-%
lighed“ desuagtet vedbliver fremdeles at „gjennemlyse“ Dog\-%
matiken. Er dette Meningen, saa er det en anden Sag; saa
% --- p. 96 ---
\opage{96}\setcounter{footnote}{0}%
tør vi da ikke vente nogen Confirmation, men maae, hvad
end Udfaldet bliver, vove os ud paa egen Haand. Vi gaae
nu over til en Behandling af vor sidste, Systemets første,
store Mirakelforklaring, Forklaringen af \emph{Skabelsesmi\-%
raklet}.

Om min Opfattelse af Skabelsesbegrebet udtaler Biskop
Martensen sig i sit Leilighedsskrift (S. 59---60) paa følgende
Maade. „Vi skulle ikke opholde os ved Forfatterens Frem\-%
stilling af det theologiske Skabelsesbegreb, der er meget ufuld\-%
stændigt og ucorrect. Kun skulle vi i Almindelighed be\-%
mærke, at Skabelsen ubetinget forudsætter det ethiske Guds\-%
begreb, og at det uden dette er aldeles urimeligt at tale om
Skabelse. Skabelsesbegrebet fremstilles derfor aldeles urigtigt,
naar man, som Professor Nielsen paa mange Steder ynder
at fremstille det, bestemmer det som en reen Vilkaarligheds\-%
act af den guddommelige Almagt, hvorved man kun bliver
staaende ved en physisk Bestemmelse, medens der ved Ska\-%
belsesbegrebet kun er og kan være Tale om den ved Kjær\-%
ligheden og Viisdommen bestemte Almagt. Vil man derfor
kritisere Skabelsesbegrebet, da kritisere man først og fremmest
selve det ethiske Gudsbegreb og vise, at det er umuligt, at
Almagten kan være ethisk bestemmet. Har man beviist dette
--- men hidtil er Beviset udeblevet --- kan man spare sig
Uleiligheden saavel med Skabelsesbegrebet som med alle de
øvrige Dogmer. Angriber man Skabelsesdogmet ved at
polemisere mod den nøgne Almagt og en blot Vilkaarlig\-%
hedsact, gjør man sig Sagen vistnok meget let, men fægter
ogsaa kun i Luften til sin egen og de Uvidendes For\-%
nøielse.“

Efter dette Vidnesbyrd at dømme staaer min Sag kun
maadeligt. Mit Haab, mit skjønne sangvinske Haab, at Bi\-%
skop Martensen selv vilde have hjulpet mig mod Dogma\-%
tiken, er nu saa godt som svundet; kun et Glimt er
tilbage; det glimter i en Anmærkning (Leilighedsskr. S. 60):
% --- p. 97 ---
\opage{97}\setcounter{footnote}{0}%
„Theologien saavelsom den christelige Philosophie og Theoso\-%
phie (J. Bøhme, Fr. Baader) har til alle Tider i Skabelsen
erkjendt et uudgrundeligt Mysterium, og har aldrig udgivet
sig for at kunne udgrunde Skabelsens Hvorledes“. En meget
vakker lille Anmærkning, ret som greben ud af Grund\-%
ideernes Logik; jeg støtter mig til den og mener fremdeles,
at Dogmatiken bør have Da-Da, naar den er saa uartig at
ville udgrunde det Uudgrundelige, udgive sig for at kunne ---
naturligviis kun til en vis Grad --- „udgrunde Skabelsens
Hvorledes“. Man høre blot det høie Maal, hvori Dogma\-%
tiken begynder sin § 59 om Skabelsen: „At Gud skaber,
indeholder, at han frembringer det, som ikke er Gud, det,
hvis Væsen er forskjelligt fra hans eget, en fri Endelighed,
hvilken han vil opfylde med sin Fylde“. Dersom det var et
med Vidensbegrebet absolut ueensartet Troesbegreb, hvortil
her sigtedes, da skulde vi fra vort Synspunkt ikke have Noget
at indvende; men hvorledes skulde „den høieste Videnskab“
med sin „gode Villie“ vel kunne nedlade sig til at adskille
Troesbegreber fra Vidensbegreber? Hvis Dogmatiken indrøm\-%
mede Gyldigheden af en med en hvilkensomhelst Udvikling af
Vidensbegreber absolut ueensartet Udvikling af Troesbegreber,
vilde den dogmatiske Christus, skjøndt han selv er baade
Theos og Logos, jo ikke kunne blive Theologus, altsaa heller
ikke Summus Theologus, altsaa heller ikke en Videnskabernes
Konge. Nei, ved det kongelige Princip, Dogmatiken vil
forstaae at hævde og bevare sit høie Sigte.

At det dogmatiske Skabelsesbegreb uden videre opfattes
som et vel begrundet Videnskabsbegreb, kan sees af kategoriske
Bestemmelser som: „Just fordi Gud er Kjærlighed, kan han
ikke afslutte sig i sig selv som Ideernes Gud, men maa bestemme
sig som Aand“ o. s. v.; fremdeles: „Naar der kan siges, at
Gud skaber Verden for at tilfredsstille et Savn i sig selv, da maa
dette ifølge Kjærlighedens Begreb forstaaes saaledes, at denne
% --- p. 98 ---
\opage{98}\setcounter{footnote}{0}%
Mangel ligesaameget er en \emph{Overflødighed}. Thi denne
Mangel i Gud er ikke som i Pantheismens Gud den blinde
Hunger og Tørst efter Tilværelsen, men er Eet med Fri\-%
hedens overvættes Rigdom, der ikke kan Andet, end at ville
aabenbare sig“. Allerede disse faa Sætninger, hvori Logisk
og Ulogisk, Physisk og Hyperphysisk, Nødvendighed og Kjær\-%
lighed, Mangel og Overflødighed røres sammen i speculative
Tvetydigheder, ere tilstrækkelige til at vise, at Videnskaben er
bleven troende, og Troen videnskabelig.

Idet „den christelige Dogmatik“ efter bedste Evne slaaer
Skabelsesbegrebet sammen med det guddommelige Kjærligheds\-%
begreb, gererer den sig som en velstuderet Coquette. Den
vil ikke sige det lige ud, at Skabelsesacten er en Nødvendig\-%
hedsact; thi den coquetterer med Troen. Den vil ikke sige det
ligefrem, at Skabelsesacten er en ubetinget Frihedsact; thi
den coquetterer med Viden. Hvad siger den saa? Den siger,
at Skabelsesacten er en guddommelig Kjærlighedsact. I
Kjærligheden kan Troen nu indlægge ligesaa megen Frihed,
den vil; thi Kjærligheden er jo fri: dens Savn er Tilfreds\-%
stillelse, dens Mangel Overflødighed. „Naar der kan siges,
at Gud skaber Verden for at tilfredsstille et Savn i sig selv,
da maa dette ifølge Kjærlighedens Begreb forstaaes saaledes,
at denne Mangel ligesaa meget er en \emph{Overflødighed}“. I
Kjærligheden kan Viden indlægge saamegen Nødvendighed,
den vil; thi Kjærligheden bevæges af hellig Nødvendighed:
dens Tilfredsstillelse er Savn, dens Overflødighed Mangel.
„Just fordi Gud er Kjærlighed kan han ikke“ --- mærk det
vel: han \emph{kan} ikke --- „afslutte sig selv som Ideernes Gud,
men maa“ --- mærk det vel: han \emph{maa} --- „bestemme sig“
o. s. v. „Den christelige Dogmatik“ tør ikke sige ligefrem, at
Verden er skabt af Guds Væsen\footnote{Det tør en Mand som Jacob Bøhme: „Vel have mange Skri\-%
benter skrevet, at Himmelen og Jorden ere skabte af Intet; men
% p. 98 | p. 99 (note continues)
jeg undrer mig over, at der blandt saa mange fortræffelige Mænd
ikke har været en Eneste, der har kunnet beskrive den rette Grund,
da dog den samme Gud, som nu er, har været fra Evighed af.
Hvor der ikke er Noget, bliver der heller ikke Noget; enhver Ting
maa have en Rod, ellers voxer der Intet; havde ikke den gud\-%
dommelige Naturs syv Aander været fra Evighed af, saa var der
ikke blevet nogen Engel, eller Himmel, og ei heller nogen Jord. ---
Da Gud har skabt baade denne Verden og alt Andet, har han
ikke havt andet Stof, hvoraf han kunde skabe det, end sit eget
Væsen, sig selv“. Den christel. Troesl. af D. F. Strauß, overs.
af H. Brøchner. 1ste Bd. S. 543.}, thi derved vilde den
% --- p. 99 ---
\opage{99}\setcounter{footnote}{0}%
krænke Troen; men den coquetterer jo med Troen. Hvad
siger den saa? Den siger, at Verden paa en Maade er
skabt af „det Ikke-Værende“, af den guddommelige Villies
„evige Muligheder“, altsaa af den Mangel, der er Kjærlig\-%
hedens Overflødighed. Og nu kan Troen lægge al den Intet\-%
hed, den behøver til en \textit{creatio ex nihilo} ind i dette Ikke-%
Værende, i disse Muligheder, i denne Mangel. Men „den
christelige Dogmatik“ tør heller ikke sige ligefrem, at Verden
er skabt af Intet; thi derved vilde den krænke Viden; men
den coquetterer jo med Viden. Hvad siger den saa? Saa
siger den, for ikke at compromittere „den nøgne Almagt“, at
det Intet, hvoraf Verden er skabt, naturligviis ikke er = 0,
„mod hvilken Forestilling det gjælder: \textit{ex nihilo nihil fit}“;
men at det er et \emph{positivt} Intet, et indholdsrigt og fyldigt
Intet, et Intet af „magiske Virkeligheder“, „plastiske For\-%
billeder“, „en himmelsk Hærskare af Syner“, den Mangel i
Kjærlighedens Overflødighed, hvormed den skabende Gud alene
kan tilfredsstille „et Savn i sig selv“.

Hvis nogen Læser drager i Tvivl, at den her beskrevne
Mediation skulde være i kategorisk Overeensstemmelse med den
dogmatiske Text: da bede vi ham sammenligne, hvad her under
forskjellige Omskrivninger er gjengivet, med hvad der bog\-%
stavelig staaer i Dogmatikens § 57, § 59, § 60, § 61. Man
% --- p. 100 ---
\opage{100}\setcounter{footnote}{0}%
skal ikke kunne sige, at Skabelsesbegrebets ueensartede Ele\-%
menter blive enten logisk eller ontologisk sammenarbeidede;
de blive kun i al Tvetydighed sammenrørte\footnote{Da Biskop Martensen (Leilighedsskr. S. 114) i Anledning af mine
Angreb paa de dogmatiske Mirakelforklaringer henviser mig til
Strauß, som „med større Klarhed og mindre Lidenskabelighed“
fremsætter sine Indvendinger, vil jeg i denne Sammenhæng be\-%
nytte Leiligheden til at minde om, hvilket Skudsmaal Strauß har
givet den Mediation, vi her betegne som en Sammenrøren af ueens\-%
artede Bestanddele. „Ikke Enhver“, siger Strauß (Den christelige
Troeslære, oversat af H. Brøchner. 1ste Bind S. 61), „besidder det
Apparat eller den Udholdenhed, med hvilken Schleiermacher, for at
iværksætte en Blanding, pulveriserede Christendommen og Spinozismen
saa fiint, at der hører et skarpt Øie til at skjælne de sammenblandede
Bestanddele fra hinanden; hos Mange er det christelige og det mo\-%
derne Element ikke bedre forenede end sammenrystet Olie og Vand,
der kun saalænge vise sig forenede, som man vedbliver at ryste
dem; medens atter andre, og tildeels ikke uberømte, Dogmatiker
endog ligne en Pølsemasse (Billedet er ikke uædlere end Sagen),
i hvilken den orthodoxe Kirkelære udgjør Kjødet, den Schleier\-%
macherske Theologie Flæsket, den Hegelske Philosophie Krydderiet.
Det er disse Blandinger, i hvilke det Bedærvede ved allehaande
Tilsætninger igjen skal gjøres velsmagende, som allerede Lessing
fandt saa ækle, saa modbydelige, saa opstødende; det er den saa\-%
kaldte fornuftige Christendom, om hvilken han tilstod, at han
hverken vidste, hvor dens Fornuft eller hvor dens Christendom
stak; den moderne Theologie, for hvilken han foretrak den gamle
orthodoxe af den Grund, at denne aabenbart streed mod den sunde
Menneskeforstand, medens hiin hellere vilde bestikke den“.}. Man kan
neppe engang sige, at Mediationen bestaaer i en Art Overeens\-%
komst eller Concordat; thi Tvetydighederne ere i en altfor
flydende Tilstand til at kunne bevare endog blot et Omrids
af noget Concordatagtigt; snarere kunde man, naar man ab\-%
straherer fra Fremstillingens qvasiæsthetiske Form, sige, at det
% --- p. 101 ---
\opage{101}\setcounter{footnote}{0}%
Sammenrørte medieres i en halvklar Vælling, en ethisk-%
hyperphysisk Kjærlighedsvælling.

Ved en saadan Mediation maa det da vel blive ind\-%
lysende, at den dogmatiske Fremstilling, trods dens Anstren\-%
gelser med at give sig et videnskabeligt Udseende, er ueens\-%
artet med Videnskab; thi Vælling og Videnskab ere ueens\-%
artede. Kjærlighedsbegrebet, det høieste ethiske Gudsbegreb, er
et Troesbegreb, som Dogmatiken mishandler i Videnskaben.
For at gjøre Mishandlingen kjendelig, have vi ved at an\-%
gribe Skabelsesdogmet polemiseret, ikke mod „den nøgne Al\-%
magt“, men mod den sminkede Afmagt. Man behøver kun
at løfte en Flig af Tvetydighedens Slør for at opdage den
Vrimmel af Modsigelser, der skrige bag Sløret. At dysse
disse Modsigelser i Søvn ved at kaste Kjærlighedskaaben over
dem, hjælper ikke, saa lidt som det hjælper at give dem
Vælling. Der er kun Eet, der hjælper, og det er: at rydde
Reden. Det kan gjøre Faderhjertet ondt --- det gjør jo
ogsaa mig ondt ---, men bort med utidig Skaansel! Dog\-%
matiken er blødagtig; Dogmatiken er forkjælet; Dogmatiken
maa i Skole; Dogmatiken maa have Riis, indtil den kommer
af Vane med at speculere i Tvetydigheder, pukke paa Philo\-%
sophie og mediere i Vælling.

Biskop Martensen vil ikke opholde sig ved min Frem\-%
stilling af det theologiske Skabelsesbegreb, der er „meget ufuld\-%
stændigt og ucorrect“. Hvorfor har da Biskoppen ikke villet
vise mig, hvori det Ufuldstændige og Ucorrecte bestaaer?
Naar ingen bedre Underrettet vil tage sig af min dogmatiske
Opdragelse, er det jo intet Under, om jeg skifter og skifter
uden at træffe det Rette.

Det bebreides mig, at jeg opfatter Skabelsen som en
Vilkaarlighedsact, et Værk af „den nøgne Almagt“; men
denne Bebreidelse rammer jo ikke mig: den rammer de gamle
Kirkelærere, den rammer Middelalderens Scholastikere, den
rammer alle vore protestantisk-orthodoxe Dogmatikere. Skal
% --- p. 102 ---
\opage{102}\setcounter{footnote}{0}%
den Almagt, der frembringer en Verden af Intet, kaldes
„den nøgne Almagt“, da kan jeg jo kalde Fædrene --- lige\-%
fra Irenæus\footnote{Man læse Irenæus \textit{adv. hæreses II, 10, 4: homines quidem de
nihilo non possunt aliquid facere, sed de materia subjacente.
Deus autem materiam fabricationis suæ, cum ante non esset,
ipse adinvenit.} Man læse Augustinus. \textit{Confess. XII, 7: Fecisti
coelum et terram non de te; nam esset æquale unigenito tuo
--- et aliud præter te non erat, unde faceres: et ideo de nihilo
fecisti coelum et terram.}} til Qvenstedt, og deriblandt En, der dog vist
vil gjøre et Indtryk paa Biskop Martensen, nemlig den
hellige Thomas\footnote{Man sammenligne \textit{Summa I, 45, 2: Antiqui philosophi non
consideraverunt nisi emanationem effectuum particularium a
causis particularibus, quas necesse est præsupponere aliquid in
sua actione. Et secundum hoc erat communis eorum opinio,
ex nihilo nihil fieri. Sed tamen hoc locum non habet in prima
emanatione ab universali rerum principio \dots} med Efterfølgende:
\textit{Si ergo Deus non ageret nisi ex aliquo præsupposito, seque\-%
retur quod illud præsuppositum non esset causatum ab ipso.
Ostensum est autem supra, quod nihil potest esse in entibus,
quod non sit a Deo, qui est causa universalis totius esse. Unde
necesse est dicere quod Deus ex nihilo res in esse produxit;}
man vil da ikke kunne betvivle, at det \textit{universale rerum princi\-%
pium}, hvorom Scholastikeren taler, netop er det \textit{liberrimæ volun\-%
tatis imperium}, hvorom Qvenstedt taler.} --- til Vidner paa, at det ikke er mit,
men Biskop Martensens Skabelsesbegreb, der er „meget ufuld\-%
stændigt og ucorrect“. Naar en Qvenstedt, hvad jeg netop
har anført som Exempel (Grundideernes Logik II, S. 150),
definerer Skabelsen som: \textit{Actio Dei triuni externa, qua is
res omnes visibiles ex nihilo \dots\ solo liberrimæ volun\-%
tatis suæ imperio omnipotenter et sapienter produxit},
lægger han saa ikke den afgjørende Vægt paa, at Skabelsen
er en ubetinget Frihedsact, et Værk af Almagten? At Kjær\-%
ligheden og Viisdommen maae antages at virke med i Alt,
% --- p. 103 ---
\opage{103}\setcounter{footnote}{0}%
hvad Almagten gjør, er en Selvfølge; men derfor staaer det
ligefuldt fast, at det er Almagten, ikke Kjærligheden, den
ubetinget frie Villiesact, ikke Viisdommen, hvorpaa Eftertrykket
falder, naar Talen er om en Bestemmelse af Skabelsens Begreb.
Det er, hvad enhver theologisk Candidat kan vide, aldeles
bagvendt at ville aflede Skabelsesbegrebet af det ethiske Guds\-%
begreb istedenfor at udvikle det af Almagtsbegrebet.

Naar Biskop Martensen vil presse Skabelsesproblemet
ind i Spørgsmaalet om Muligheden af, „at Almagten kan
være ethisk bestemmet“, og det med en Knaldeffect som denne:
„Angriber man Skabelsesdogmet ved at polemisere mod den
nøgne Almagt og en blot Vilkaarlighedsact, gjør man sig
Sagen vistnok meget let, men fægter ogsaa kun i Luften til
sin egen og de Uvidendes Fornøielse“, da vil selv den mindre
Indviede mærke Uraad og faae en Anelse om, at Dogma\-%
tikeren nok vil bruge sit „ethiske Gudsbegreb“ til Skjold for at
hugge sig igjennem og komme bort fra det farlige Problem.

I Grundideernes Logik II, S. 149---166 er Skabelsesbe\-%
grebet behandlet, vistnok udelukkende fra Logikens Standpunkt, men
dog med udtrykkeligt Hensyn saavel til den orthodoxe Lære som
til den Martensenske Speculation. Det bliver her gjennem
en Række Analyser viist, hvor tvetunget, selvmodsigende og
vakkelvorn den dogmatisk-speculative Behandling af Ska\-%
belsesdogmet er fra Først til Sidst. Her er i disse Ana\-%
lyser Punkter nok at tage fat paa. Hvis Biskop Martensen
i hele dette Parti har fundet enten Fordreielser eller Forbi\-%
gaaelse af væsentlige Punkter eller andre Feilgreb: hvorfor
har han da ikke udtrykkelig paatalt det? Har Biskop Mar\-%
tensen virkelig den Tillid til sit ethiske Gudsbegreb, at det i
Dogmatiken kan løse os Skabelsesproblemet, og at det kun
er logisk Forblindelse, der har ledet mig til „at fægte i
Luften“: hvorfor har han da ikke paa bestemt Maade af\-%
væbnet Logikens bestemte Indsigelser ved en aldeles bestemt
Anvendelse af det ethiske Gudsbegreb?
% p. 103 ends short (last quarter of the page blank); p. 104 opens flush left, no indent: paragraph break kept

% --- p. 104 ---
\opage{104}\setcounter{footnote}{0}%
Jeg vil ikke trætte Læseren med enten at udskrive eller om\-%
skrive de i Logiken gjorte Indsigelser, men overlade det til
Vedkommende selv at slaae de anførte Steder efter, saavel i
Biskop Martensens Dogmatik som i min Logik, og endnu kun
til Slutning gjøre opmærksom paa det typisk Theologiske i
at ville ved almindelige Betragtninger over Facit og Resultat
unddrage sig det egentlige Regnskab og retfærdiggjøre slige
Udflugter ved det høitravende Ord: Totalanskuelse.

Efter disse Forberedelser gaae vi nu over til en Be\-%
tragtning af den Stilling, Philosophien vil komme til at
indtage i Monarchiet. Da Ordet Philosophie er fleertydigt,
maae vi strax bemærke, at det er Philosophien som \emph{Viden\-%
skab}, udelukkende som Videnskab, hvorom her tales. Som
Videnskab er Philosophien med Hensyn paa sit absolute Ind\-%
hold Ideelære, og som Ideelære en eiendommelig, fra de
særskilte Videnskaber, de saakaldte Fagvidenskaber, forskjellig
Videnskab.

Men med al Eiendommelighed og indbyrdes Forskjellig\-%
hed maae alle Videnskaber uden Forskjel dog hvile paa et
almeenvidenskabeligt Grundlag. Det almeenvidenskabelige
Grundlag er bestemt ved den rationelle Nødvendighed og naaer
ikke ud over det Logisk-Physiske. Paa Grund af Grundlaget
er det umuligt at tænke sig absolut ueensartede Videnskaber.
Saalidt som der kan gives en med de exacte og positive
Videnskaber absolut ueensartet Philosophie, ligesaa lidt kan
der gives to Slags Philosophie, af hvilke den ene skulde være
absolut ueensartet med den anden. Men naar Philosophien
umuligt kan spaltes i to absolut ueensartede Philosophier,
naar den umuligt kan udvikle sig af to absolut ueensartede
Principer, saa sees allerede heraf, at dens Stilling i Mon\-%
archiet ikke kan blive misundelsesværdig.

Ved at henvise os til Paradiis, til „\emph{de tvende abso\-%
lut ueensartede Træer}, af hvilke det enes Frugter vare
Mennesket forbudne“, henviser Biskop Martensen os til
% --- p. 105 ---
\opage{105}\setcounter{footnote}{0}%
tvende absolut ueensartede Videns- og Videnskabsprinciper,
af hvilke det enes Udvikling absolut maa være os Mennesker
forbuden. Her have vi igjen det Elias-, om man vil,
Makkabæusagtige; vi have „den gode Villie“.

Det gjælder en Renselse af Philosophiens Helligdom.
Allerførst maa den gamle græske Philosophie feies ud; thi det
er vist, at hverken Plato eller Aristoteles nogensinde have
tænkt paa at gaae ud over det Logisk-Physiske; det er vist,
at hverken Plato eller Aristoteles ville vedkjende sig en Viden,
der bestemmes, ikke med rationel Nødvendighed, men ved „den gode
Villie“; det er vist, at den græske Viden er autonomisk, ikke theo\-%
nomisk; det er vist, at den græske Philosophie aldrig nogensinde vil
kunne omdannes til enten selv at blive en Troesphilosophie eller
anden Gang at bøie sig under Theologien som en Tjeneste\-%
qvinde under sin Herskerinde, som en Hagar under en Sara.
Det var kun i Middelalderen, at Plato var „der große
Pfaff“; det var kun i Middelalderen, at Aristoteles gjaldt
for en \textit{christianus ante Christum}; det var kun i Middel\-%
alderen, at Theologien var Scholastik, og at Scholastiken
sammenskruede det absolut Ueensartede. Skulle vi vel kalde
Middelalderen tilbage? Nei, den klassiske Philosophie maa
ud af Monarchiet.

Men den nyere Philosophie fra Cartesius indtil Hegel
maa, ved Makkabæus! ogsaa ud af Monarchiet. Den car\-%
tesianske Philosophie begynder jo netop med en Grundsæt\-%
ning, der staaer i absolut Strid med \textit{credo, ut intelligam},
nemlig: \textit{de omnibus dubitandum est}. En Philosophie, der
begynder med Tvivl, og en Philosophie, der begynder med
Tro, ere uforenelige. Havde det været muligt at forene en
Troesphilosophie med en Philosophie, der har Tvivlen til
methodisk Udgangspunkt, da havde vi jo nu i vor egen Lite\-%
ratur havt det første Grundlag for en saadan Forening i
J. L. Heibergs Udvikling af den hegelske Philosophie. „Den
hegelske Philosophie“, siger Heiberg (Perseus Nr. 1 1837,
% --- p. 106 ---
\opage{106}\setcounter{footnote}{0}%
S. 40): „hviler ikke blot paa Guds Aabenbaring i Almin\-%
delighed (thi det Samme gjøre alle de andre Philosophier),
men den hviler paa den christelige Aabenbaring isærdeles\-%
hed, og den har denne Forudsætning med en Bevidsthed, der
kommer fra Slutningen til Begyndelsen, og netop derfor kan
absolut ignorere sig selv, uden at tabe eller glemme sig.“

Heiberg gjør netop i den anførte Afhandling en velbe\-%
grundet Indsigelse imod det Martensenske Project med et
eget philosophisk Troesprincip. Han siger: „Den Fordring, % KB copy: manuscript note in the left margin here ("Pers. Skr. (61) II 51"?), not transcribed
som Hr. Martensen gjør til en Troes-Philosophie, synes
allerede at være opfyldt i den hegelske. Thi at denne ikke
desmindre i Begyndelsen af sit System gjør sig til Igno\-%
rant af hvad den veed, og til Tvivler, ja endog Fornægter
af den Gud, som den erkjender, skeer kun for at opfylde den
samtidige Fordring: forudsætningsløs Begyndelse, og kan
neppe blive anderledes ved nogen Philosophie, der virkelig
skal svare til sit Navn, og altsaa ikke være Religion. \dots\
Skulde Hr. Martensen indvende herimod, at om end Guds
Aand er den hegelske Philosophies sidste Resultat, saa frem\-%
stilles dette kun logisk, ikke med den christelige Aabenbarings
hele Fylde for Hjertet eller Følelsen, da maa jeg herpaa
svare, at just deri bestaaer den evige Forskjel mellem Philo\-%
sophie og Religion, thi hvis denne Forskjel ikke var evig,
vilde Philosophien jo blive Alt, og Religionen, saavelsom Kun\-%
sten, overflødig. Men før jeg indlader mig herpaa, vil det
være hensigtsmæssigt at oppebie den nærmere Udvikling, som
Hr. Martensen har givet os Haab om at skjænke os.“

Den belovede Udvikling er nu skjænket os; vi have dens
modne Frugter i „vor christelige Dogmatik“. Den Mar\-%
tensenske Dogmatik er et Exempel paa, hvad der kommer ud
af at gaae Cartesius glat forbi og slutte sig til Anselmus, at
forkaste Philosophens \textit{de omnibus dubitandum} og holde
fast ved Scholastikerens \textit{credo, ut intelligam}. Det er
slet ikke Andet end at ville gaae ud over det Logisk-Physiske,
% --- p. 107 ---
\opage{107}\setcounter{footnote}{0}%
at ville bort fra den rationelle Nødvendighed, at ville op\-%
føre en videnskabelig Lærebygning udenfor Videns Grændser,
altsaa udenfor det videnskabelige Grundlag. Men viden\-%
skabelige Lærebygninger uden videnskabeligt Grundlag kalder
man Luftkasteller.

Hvorfor begynder vel Philosophien med en Tvivlen om
Alt? For Undersøgelsens, Prøvelsens og Forvisningens
Skyld, at den rationelle Nødvendighed kan komme til at
fortrænge Tvivlen! Videnskaben kjender ingen Vished, agter
ingen Vished uden den, som paa logisk-physisk Maade paa\-%
tvinges Bevidstheden. En Sætning, jeg baade kan antage og
ikke antage, er uden videnskabeligt Værd; kun den Sætning,
jeg umulig kan forkaste uden at forkaste Videnskaben selv,
er videnskabelig. „Men“, spørger Biskop Martensen (Lei\-%
lighedsskr. S. 102), „skal kun det kaldes Videnskab, der kan
paatvinges med uimodstaaelig Nødvendighed, da synker Viden\-%
skaben ned til noget meget Tomt og Værdiløst, som det ikke
er Umagen værdt at gjøre store Ophævelser over. Kun den
formelle Logik og Mathematik blive da Videnskaber i stren\-%
gere Forstand.“ Denne Erklæring røber en saa forunderlig
stor Miskjendelse af det, som giver Videnskab og Viden\-%
skabelighed under alle Skikkelser Værd og Betydning, at der
blandt Dyrkerne af de allerforskjelligste Videnskaber neppe
vil være Nogen, der ikke umiddelbart og ved første Blik
seer Misforstaaelsen i hele dens Udstrækning. Thi hvad er
det vel, hvorefter en Historiker maa stræbe ligesaa ivrig som
en Mathematiker, en Botaniker ligesaa ivrig som en Mecha\-%
niker, en Anthropolog ligesaa ivrig som en Astronom? Det
er \emph{Vished, uimodsigelig Vished}! Hvorledes og paa % the "!" is a deformed sort (or overprinted by a pen stroke); read as "!"
hvilke Betingelser naaes denne Vished? Den naaes i de
forskjellige Videnskaber ad forskjellige Veie og paa forskjellige
Vilkaar; men Eet er fælles for al videnskabelig Vished,
dette Ene, at den beroer og maa beroe paa almeengyldige,
% --- p. 108 ---
\opage{108}\setcounter{footnote}{0}%
uimodsigelige Grunde, at den maa „paatvinge sig med uimod\-%
staaelig Nødvendighed.“

Til Understøttelse for sin Erklæring bemærker Biskoppen:
„Allerede en speculativ Logik gaaer ud over det blot Tvin\-%
gende; thi Ingen vil f. Ex. kunne fatte det Store i Hegels
Logik eller i et Værk som Kants Fornuftkritik uden archi\-%
tektonisk Phantasie“. Denne Bemærkning er intetsigende.
Thi der er heller Ingen, der vil kunne fatte det Store
i Newtons \textit{Principia} eller i et Værk som Laplaces \textit{Mécanique
céleste} „uden architektonisk Phantasie“. Den Forsikkring,
at der om saadanne Værker som Hegels Logik og Kants
Fornuftkritik kan disputeres i det Uendelige, og at uendelige
Misforstaaelser her ere mulige, er til Lykke for Videnskaben
greben ud af Luften. Kunde der blandt Sagkyndige dis\-%
puteres om Hegels Logik og Kants Fornuftkritik i det Uende\-%
lige, uden at sikkre, uimodsigelige Resultater lode sig opnaae:
hvad vilde saa det bevise? Det vilde bevise, at Hegels
Logik og Kants Fornuftkritik vare uden videnskabeligt Værd.

Staaer det kun fast, at ingen Philosophie --- her tales
om Philosophien som Videnskab --- nogensinde har forsøgt
eller nogensinde vil forsøge at gaae udenfor Grændsen af
det Logisk-Physiske eller hæve sig op over det rationelt Nød\-%
vendige, og denne Hovedsætning skal vel staae fast: da er
det mig nu, hvad enhver tænkende Læser kan sige sig selv,
saare let at bevise, at \emph{den Protest, Biskop Martensen
i sit Leilighedsskrift har nedlagt mod min Philo\-%
sophie, er en Protest imod al Philosophie}.

Ved at skue hen til Paradiis og stadig have de tvende
\emph{absolut ueensartede Træer} for Øie har Biskop Mar\-%
tensen seet sig istand til at holde to absolut ueensartede
Vidensprinciper imod hinanden. Der gives da efter hans
Synsmaade en vantro og en troende Viden, en vantro og
en troende Videnskab: den vantro Viden har „den onde Villie“,
den troende Viden „den gode Villie“ til Princip. Heri sees
% --- p. 109 ---
\opage{109}\setcounter{footnote}{0}%
det Originale. Det er Opstillingen af de tvende absolut
ueensartede Træer i Forening med de tvende absolut ueens\-%
artede Villier indenfor den objective Videns Synskreds,
hvorved den Martensenske Philosophie bliver epochegjørende i
Videnskabernes Historie.

Vel er det ikke saameget Hensyn til Philosophien som
til Theologien, der her bevirker Gjennembrudet; dog haabe
vi, at den fra Træerne og Villierne udgaaede Revolution skal
vorde til Held og Lykke, om ikke for Theologien, saa dog
for Philosophien.

Som Aandsvidenskab har Philosophien, endog i den
nyere Tid, havt en Tilbøielighed til at assimilere sig Troens
Forudsætninger og blive lidt theologisk. Heibergs Ord, at
den hegelske Philosophie hviler „ikke blot paa Guds Aaben\-%
baring i Almindelighed men paa den christelige Aabenbaring
i Særdeleshed“, ere i denne Henseende charakteristiske. Det
er betænkeligt. Endskjøndt den nøiere Tilslutning af Philo\-%
sophien til Theologien jo vistnok i det Hele har havt til Følge,
at Theologien (hvad man seer paa Dogmatikere som Daub
og Marheineke) kom til at gaae i Philosophiens Ledebaand,
ikke omvendt: saa har det Theologiske dog kunnet inficere
Philosophien selv i den Grad, at Hegelianere som Weiße have
drevet Treenighedslæren indtil Karikatur. I velbegrundet
Indignation over en saadan Karikatur, udbryder Strauß:
„Altsaa derfor, gode Hegel, har Du fundet dine dybsindige
Kategorier: Bliven til Andet og Udtræden af sig (Entäuße\-%
rung), Negation og Negationens Negation, Momenternes
Ophævelse i et Høiere, at en raa Phantasies crudeste Mis\-%
fostre, ved hvilke den afsindigste Gnostiker ikke vilde have
skammet sig, skulde lade sig indfatte i dem? Tre Per\-%
soner, udtrykkeligt bestemte som forskjellige Bevidsthedsmiddel\-%
punkter, skulle dog være absolute; to saadanne Personer skulle
i en tredie Person forenes til en Eenhed, der er mere end
en blot Overeensstemmelse; en Person og tilmed en ab\-%
% --- p. 110 ---
\opage{110}\setcounter{footnote}{0}%
solut, skal en Tidlang blive borte for sig selv, blive uper\-%
sonlig, splittes ad i en uendelig Mængde endelige Personlig\-%
heder, ligesom en Kartoffelstamme i Knolde, men efterat en
vis Tid er forløben igjen erholde sin Personlighed, hvorhos
tillige de i den mellemliggende Tid opstaaede creaturlige Per\-%
sonligheder uden at ophøre at være mange og fra hinanden
adskilte, skulle flyde sammen i den guddommelige Aands ene
Personlighed. --- Hvor er \textit{Symbolum Quicunque?} Rækker
mig det! før vil jeg ti Gange sværge paa det, før jeg blot
een eneste Gang vil kalde vor Philosophs Sætninger andet
end Afsindighed.“\footnote{Troeslære, overs. af H. Brøchner. 1ste Bd. S. 435.}
Denne Fristelse til at være med i
Troen og explicere sig theologisk vil for Philosophiens Ved\-%
kommende naturligviis falde bort, naar Philosophien nu bliver
dreven ud af Monarchiet. Og det vil skee kategorisk. Phi\-%
losophien kan umulig skille sig ad i to absolut ueensartede
Philosophier, og endda hænge sammen saaledes, at den ene
subordineres den anden. Skal en Philosophie med \textit{credo, ut
intelligam} gjælde i Monarchiet, maa Philosophien som Vi\-%
denskab ud af Monarchiet.

Ved at anerkjende Theologien, vel som en Videnskab,
der trængte til en grundig Reform, men dog som en Viden\-%
skab, blev Philosophien, endog som Naturphilosophie, for\-%
ført til at overfare de virkelige Kjendsgjerninger paa samme
geniale Maade som Mystikerne og de speculative Theologer
altid have gjort. I Hegels Naturphilosophie --- om Schel\-%
lings Phantasieanskuelser ville vi slet ikke tale --- findes
Opfattelser, Forklaringer og Overblik, der ingenlunde ere en
Biskop Martensen uværdige. Ogsaa Hegel antog, „at Natur\-%
videnskaberne, skjøndt de høre til Nødvendighedens Sphære,
paa ingen Maade kunne kaldes Videnskab i strengere For\-%
stand“; ogsaa Hegel kunde mene, at Naturvidenskaben var
en underordnet Videnskab, fordi den „ingen andre Grunde
% --- p. 111 ---
\opage{111}\setcounter{footnote}{0}%
havde at støtte sig til end Induction og Analogie“. Denne
daarlige Indbildning var nær kommen Philosophien dyrt at
staae. Den schelling-hegelske Philosophie kom, just ikke ved
at troe paa Mirakler, men ved at troe paa sin egen absolut
adæqvate Viden af det Absolute, i den Grad bort fra Erfaring
og Virkelighed, at den philosophiske Teleologie stod i Begreb
med at blive ligesaa huul og phantastisk som den theologiske.
Ved den fra Træerne udgaaede Revolution er Fristelsen
endt: Philosophien som Videnskab gaaer ud af Monarchiet.

Det har i en Række af Aar været et Maal for mine
Bestræbelser at sikkre de philosophiske Undersøgelser et virkeligt
Grundlag og faae Forholdet mellem Philosophien og Fag\-%
videnskaberne nogenlunde klaret. Jeg maa ved denne Leilig\-%
hed endnu engang minde om det i Indledningen til Grund\-%
ideernes Logik I fastholdte Synspunkt: „Skal en „Begrebets“
Philosophie kunne sikkre sig Begrebets Realitet, maa den dog
vel fremfor Alt være opmærksom paa, hvorledes et Real\-%
begreb dannes; men, i hvilken Grad den hegelske Philosophie,
netop hvad dette væsentlige Punkt angaaer, har forfeilet Be\-%
grebsdannelsen, er iøinefaldende, naar man kaster et Blik
paa Naturphilosophien. At Hegel ikke har dannet sine første,
elementære Naturbegreber under Veiledning af nogen mathe\-%
matisk Physik, men paa reent logisk speculativ Maade, er
i alle Punkter indlysende. Det sees ikke blot deraf, at
han udtrykkelig bekæmper den mathematiske Physiks Beteg\-%
nelser og Forestillinger; men det sees fornemmelig deraf, at
han ikke engang har fundet det nødvendigt at sikkre sig en
authentisk Opfattelse af de Kategorier og Synsmaader, han
hos Mathematikerne vil bekæmpe.\footnote{Jvfr. Forelæsninger over „Phil. Propædeut. fra Universitetsaar
1861--62.“ S. 49--75.}
De Misligheder, her
paapeges, angaae ikke Feiltagelser i dette eller hiint Enkelte;
% --- p. 112 ---
\opage{112}\setcounter{footnote}{0}%
slige Feiltagelser ere i en philosophisk Fremstilling af væsent\-%
lig Gehalt altid mere eller mindre uvæsentlige; --- at gjøre
Væsen af det Uvæsentlige maa (Jvfr. Begrebet: Nøiagtighed
S.3--11) naturligviis forbeholdes Pedanter og Rabulister ---;
nei, de her paapegede Misligheder angaae Methoden selv, og
forraade en Svaghed, der, saa at sige, er Systemet medfødt.
Det er Maaden, hvorpaa Naturkategorierne fra først af op\-%
fattes og behandles, som her er forfeilet. Naturbegreberne
ere Virkelighedsbegreber; men for at Virkelighedsbegreberne
kunne dannes, maae de forskjellige Fagvidenskaber, hver med sine
Iagttagelser, Analyser og Beregninger, gaae foran og vise
Vei. Her er i hele Naturriget ikke en eneste Begrebssphære,
hvorom det gjælder, at den philosophiske Tænker, uden særlig
Forhandling med den paagjældende Fagvidenskab, selv skulde
være istand til at vise Kategoriernes Udspring og finde deres
Rødder i Realitetens Jordbund.“

For at bestemme Forholdet imellem Philosophien og
Fagvidenskaberne har jeg i mine propædeutiske Undersøgelser
ikke kunnet indskrænke mig til almindelige Forhandlinger om
det Empiriske og det Aprioriske, men har, saavidt Kræfterne
tillode det, maattet gjøre encyclopædiske Studier i et ikke
ganske ringe Omfang. At her endnu staaer meget tilbage,
indrømmer jeg; men jeg meente dog at have bragt det saa\-%
vidt, at Enhver, der fulgte med, vilde betænke sig to Gange,
inden han paany opvartede Læseverdenen med Opkog af den
schelling-hegelske Viisdom, som: „at Naturvidenskaberne, skjøndt
de høre til Nødvendighedens Sphære, ingenlunde kunne kaldes
Videnskab i strengere Forstand.“

At Biskop Martensen, naar han engang i et Leiligheds\-%
skrift fik Leilighed til at nedlægge Protest imod min Philo\-%
sophie, vilde tage mig ordentlig op i Propædeutik, kommer
mig ikke uventet. Har jeg givet Dogmatiken op udenom,
bør jeg naturligviis examineres skrapt i Propædeutiken. Een
Ting kan jeg dog ikke ret forstaae. Maaskee jeg tager feil;
% --- p. 113 ---
\opage{113}\setcounter{footnote}{0}%
men det forekommer mig, som om den høie Mine, hvormed
Biskop Martensen vil fungere som Examinator, staaer i Mis\-%
forhold til den ringe Grad af Forberedelse, hvormed han er
gaaet til Examinationen. „Man nødes“, siger Biskoppen
(Leilighedssk. S. 97), „til at spørge om“ min „Syslen med
Dogmatiken“ ikke har for mig medført „en ikke ringe For\-%
sømmelse af Philosophien“; og jeg nødes ved denne Leilighed til
at spørge, om Hs. Hvdhds megen Syslen med „Concordater“
ikke har for ham selv medført en ikke ringe Forsømmelse af
Propædeutiken, særlig af § 13, § 14, § 15 og § 16. Thi
den Bispeexamen i philosophisk Propædeutik, hvortil jeg stedes,
er anlagt paa en Besvarelse af det Spørgsmaal: \emph{Hvor\-%
ledes indføre vi en saa grundig Subordination i
Ideeriget, at Theologiens „gode Villie“ kan blive
anerkjendt for Selvhersker over alle Videnskaber?}
Dette Spørgsmaal kan jeg ikke besvare. Jeg falder igjennem,
men tager i Faldet den Trøst med mig, at jeg falder som
Offer for „den gode Villie“.% a small raised stroke after the full stop (spacing material or blemish), not transcribed

Gaaer det mig nu saaledes i Propædeutiken, saa kan
Læseren sige sig selv, hvordan det maa gaae mig i Logiken.
Det er en streng Kritik, Biskop Martensen har ladet udgaae
over Grundideernes Logik, en Kritik, der i Sandhed er „det
kongelige Princip“ værdig. Vi maae nemlig lægge Mærke
til, at den monarchiske, fra de tvende absolut ueensartede
Træer udspringende Philosophie har forvandlet os ikke blot
Logiken, men ogsaa Kritiken. Ligerviis som der i Paradiis
gaves tvende absolut ueensartede Træer, saaledes gives der
nu i det theologiske Monarchie tvende absolut ueensartede
Logiker, en theologisk og en subordineret, tvende absolut ueens\-%
artede Kritiker, en theologisk og en subordineret.

Hvor trykkende Philosophiens Stilling i Monarchiet
maa blive, naar en theologisk overordnet Kritik skal behandle
en subordineret Logik, er let at indsee. Den gode Villie har
% --- p. 114 ---
\opage{114}\setcounter{footnote}{0}%
den til det Gode svarende Theorie; den onde Villie har den
til det Onde svarende Theorie: hvad Theorie har nu Lo\-%
giken? Dette Spørgsmaal kan kun besvares fra den theo\-%
logiske Kritiks overordnede Standpunkt. Paa en almeen\-%
menneskelig Videns Standpunkt veed man nemlig ikke, om
man helst skal begynde med Theorien og sige: naar Theo\-%
rien er god, saa er Villien god, thi en god Theorie kan kun
fremgaae af en god Villie; eller omvendt begynde med
Villien og sige: naar Villien er god, saa er Theorien god,
thi en god Villie maa nødvendig frembringe en god Theorie.
Vanskeligheden stiger, naar Talen er om Logiken som Theorie,
om den gode og den onde Villie i Logiken. Hvilken Villie
er der nu at opdage i en aristotelisk Logik, en god eller
ond? Vi svare: Villien i den aristoteliske Logik er hverken
god eller ond; den er blot logisk. Og hvilken Villie er der
vel i en hegelsk Logik, en god eller en ond? Vi svare:
Villien i en hegelsk Logik er hverken en god eller en ond,
men logisk. Hvilken Villie bør her da være i en Grund\-%
ideernes Logik, en god eller en ond? Vi svare: Villien i en
Grundideernes Logik bør hverken være god eller ond, men logisk.

Fra dette Synspunkt er man hidtil gaaet ud ved Be\-%
dømmelsen af de forskjellige Logiker i de forskjellige Philo\-%
sophier. Men denne gamle, forslidte Anskuelse vil nu ved
Monarchiets Oprettelse komme til at vige Pladsen for en,
„der funkler af Nyhed“. Begyndelsen er skeet med den
Revision, som Biskop Martensen i sit Leilighedsskrift har
underkastet Grundideernes Logik. Til en Prøve anføre vi
et Par af de kraftigste Steder.

Biskoppen siger (S. 32): „Prof. Nielsen kan nu eengang
ikke blive Theologien qvit. Han har den ikke blot udenfor
sig, idet han stadigt maa bekæmpe den christelige Dogmatik,
uden nogensinde at kunne finde Enden paa disse Angreb,
fordi han, som i det Følgende skal vises, endnu ikke har
kunnet finde Begyndelsen; men han har den tillige indeni
% --- p. 115 ---
\opage{115}\setcounter{footnote}{0}%
sig selv, og maa nu producere en speculativ Theologie, op\-%
stille en Lære om Gud, om Logos, om Triniteten, om
Skabelsen og de sidste Ting o. s. v., vistnok den christelige
Dogmatik saa modsat som muligt, men reent philosophisk,
reent rationelt-dogmatisk.“

Her er det Overordnede! Jeg, som paa mit under\-%
ordnede Standpunkt bildte mig ind at have bøiet begge
Ender sammen paa „den christelige Dogmatik“ og svunget
Kritikens Riis, staaer nu her --- med løftet Riis --- og kan
ikke „finde Enden“. Det var da kjedeligt!

Men, hvorledes gaaer det nu egentlig til, at Theologien
kommer op i mig, uden at jeg selv veed af det? Det gaaer
overordentlig til; thi det gaaer til paa den Maade, at den
overordnede Kritik, der, saavidt muligt, nøies med at kritisere
mine Fortaler og sjeldent kan faae Tid til at læse en Fortale
heelt igjennem, foretrækker at gribe et enkelt Stykke og kaste
mig det i Ansigtet.

I Fortalen til Grundideernes Logik II har Biskop
Martensen desværre ikke kunnet holde ud at læse længere
end til S. IX; havde han blot faaet S. X og S. XI med,
saa havde han, før han paany hexede Theologien ind i mig,
faaet en lille Opgave at løse. Det maa i denne Anledning
være nok at henlede Læserens Opmærksomhed paa følgende,
fra min Side, som mig synes, utvetydige Udtalelse (S. XI):
„Istedenfor at indlade sig med umiddelbare Gudsforestillinger
og qvakle med Raisonnementsbeviser for Guds Tilværelse,
begynder Philosophien med at forudsætte, hvad der ikke er et
Hiinsidigt eller Fraværende, men et for den almene Viden
væsentligt Nærværende, nemlig Tilværelsen selv med sit uende\-%
ligt rige Indhold. Dette Vidensindhold bestemmes da efter
dets indholdsrige Eenhed med sig selv som det Absolute, og
Opgaven bliver at udvikle og klare de for Værens absolute
Indhold adæqvate Vidensformer. Denne Opgave er viden\-%
% --- p. 116 ---
\opage{116}\setcounter{footnote}{0}%
skabelig, thi slige Former kunne udvikles, revideres og cor\-%
rigeres med en Nøiagtighed, der i Dialektiken ikke staaer til\-%
bage for det Exacte i Mathematiken.\footnote{Jvfr. Mathematik og Dialektik S. 1. 2.}
For Formernes
Udvikling er Methoden saaledes anlagt, at Indholdet som
ved et Ideens Høidetryk strømmer til af sig selv: er blot
Formen sand, saa er enhver Tvivl om Indholdets Realitet
med det Samme hævet. Som den fuldkomne Form for det
absolute Indhold bliver da Gudsbegrebet udviklet; i det ud\-%
viklede Begrebs Eenhed med sit Indhold erkjendes Guds\-%
ideen; i Gudsideen erkjendes Gud; men denne Erkjendelse
er ikke theologisk.“

Det er, hvad enhver tænkende Læser maa sige sig selv,
naturligviis ikke det \emph{ethiske}, men det \emph{logiske} Gudsbegreb,
der skal udvikles i Logiken. For imidlertid at forebygge
Misforstaaelse i denne Retning har jeg i den her omhandlede
Fortale udtrykkelig mindet om Logikens Forhold til Reli\-%
gionsphilosophien (S. VIII): „Ihvorvel en Grundideernes
Logik ikke i noget Parti kan falde umiddelbart sammen med
Religionsphilosophien, saa er det dog aldeles nødvendigt, at
den, medmindre den da skulde kunne hjælpe sig med et
atheistisk Princip, maa udvikle saavel Gudsideens som Ver\-%
densideens logiske Form, maa logisk bestemme Forholdet
mellem disse Ideer og saaledes afgive Grundlaget for en
tilsvarende Religionsphilosophie.“ Denne Antydning er i
nærværende Tilfælde ikke kommen mig tilgode.

Forvisset om, at her er Intet i Veien med Theologien,
har Biskop Martensen følt sig tilskyndet til at hjælpe lidt
paa Philosophiens forfaldne Sager ved at nedlægge Protest
imod min Philosophie. Jeg har kun opstillet et Vidensideal,
men Biskoppen fordrer et Viisdomsideal og beraaber sig paa
Plato. „Det Høieste og Dybeste hos Plato“, siger han S.
105, „findes som oftest paa de Steder, \dots\ hvor han taler
% --- p. 117 ---
\opage{117}\setcounter{footnote}{0}%
mythisk, idet han ligesom visionært berører, hvad ikke kan
naaes ad den blot dianoetiske Vei, hvad vistnok heller ikke lader
sig bringe ind i en afsluttet Theorie, men ikke desto mindre
hører med til Philosophien, ja til dens Høieste, efterdi Phi\-%
losophien væsentligt er Stræben efter Viisdom, og derfor
ikke en afsluttet Theorie. Men netop herved føres vi tilbage
til det Ønskelige i, at Prof. Nielsen vilde have skjænket Phi\-%
losophien noget større Opmærksomhed, hvorved som sagt hans
Angreb paa Theologien vilde have vundet i Grundighed og
Vægt, og at han navnlig vilde have indladt sig paa at
undersøge Forskjellen mellem et Vidensideal og et Viisdoms\-%
ideal.“

Altsaa ved et Viisdomsideal og ved at gjøre opmærksom
paa Forskjellen mellem et Vidensideal og et Viisdomsideal
vil Biskop Martensen nu tage sig af Philosophien. Vi ere
ham paa Philosophiens Vegne meget forbundne; kun maae
vi bemærke, at den, der vil tage sig af Philosophien, maa
allerførst see at gjøre sig lidt bekjendt med dens Encyclo\-%
pædie. Uden Kjendskab til Encyclopædien udsætter man sig
for at løbe sur i Disciplinerne og f. Ex. lede efter Viis\-%
domsidealer, hvor man, ifølge Sagens Natur, ikke kan vente
at finde Andet end Vidensidealer. I sin Iver for at ile
Philosophien til Hjælp har Biskoppen taget Encyclopædien
i eet Spring, og er ved at springe lige lukt ind i Viisdoms\-%
idealet i den stærke Fart kommen til at blande Logik,
Ethik og Religionsphilosophie paa en høist original Maade
imellem hverandre. Der er ved Blandingen indtraadt en
fuldstændig Confusion i „Totalanskuelsen“.

Hvor gjennemgaaende Confusionen er, kan sees af de
høist besynderlige Indvendinger imod Theismen i Grund\-%
ideernes Logik, hvormed Biskop Martensen for en Deel har
fyldt sit Leilighedsskrift. Disse ethiske, thetiske, æsthetiske,
pathetiske Indvendinger ere rigtignok kun Variationer, hvori
een og samme Misforstaaelse udtværes; men de kunne dog,
% --- p. 118 ---
\opage{118}\setcounter{footnote}{0}%
naar man forsøger at bringe nogen Orden i dem, henføres
under følgende Synspunkter.

1) \emph{Theismen i Grundideernes Logik er kun
en omskreven Pantheisme.} % the numeral "1)" is not letterspaced

„Forfatteren maa“, siger Biskop Martensen (Leilighedsskr.
S. 33), „enten have taget Theisme i en heel anden Betydning,
end det i Almindelighed tages, eller han maa her have havt en
Fristelse i Retning af Mediation. Thi i Virkeligheden er den
Viden, han vil indføre, Pantheisme eller om man hellere vil,
blot og bar Rationalisme og Naturalisme, hvad ogsaa stemmer
med hans øvrige Udtalelser, idet han ikke bliver træt af at
gjentage, at der ikke gives andre Aarsager end blotte Fornuft\-%
aarsager og Naturaarsager, hvilket da tillige maa gjælde om
den høieste Aarsag, og til Overflod forsikkrer, at ethvert
„overnaturligt Princip“ bør udelukkes af Videnskaben.“ Paa
saadan Maade nedlægger Biskop Martensen Protest imod
Grundideernes Logik.

Hovedindvendingen er, som man seer, den, at Logikens
Gudsbegreb kun er logisk, ikke ethisk. Ved Protesten gaaer
da Biskoppen øiensynlig ud paa at indføre ikke „en grundig
Subordination“, men en grundig Confusion i Ideeriget. Thi her
er i denne Sag Noget, ethvert fornuftigt Menneske strax kan
forstaae, og det er Følgende: Enten tør Gudsideen slet ikke
optages i en Ideernes Logik, eller, hvis Ideen tør optages i
en Logik, da maa Gudsbegrebet udvikles logisk. En Udvik\-%
ling af det logiske Gudsbegreb er saa langt fra at beroe paa
nogen Fornegtelse af det ethiske, at den tvertimod indeholder
alle for det ethiske Gudsbegreb nødvendige Momenter. Vel
er Logik ikke Ethik, og de logiske Begreber ikke ethiske; men
derfor maa Ethiken ligefuldt have sin Logik, og de ethiske
Begreber reflectere sig i den logiske Idee. Den Viden, jeg
vil indføre, er, det paastaaer Biskoppen, „i Virkeligheden Pan\-%
theisme“; hvorfor? Fordi Logik skal være Ethik: en ret in\-%
teressant Confusion.
% --- p. 119 ---
\opage{119}\setcounter{footnote}{0}%
„Hvad der“ (Leilighedsskr. S. 33), „kan frembringe
Skinnet af Theisme er dette, at det Princip, der ifølge denne
Philosophie bærer Tilværelsen og dermed ogsaa det philoso\-%
phiske System, er et Subjectivitetsprincip. „Er det sandt“,
siger Prof. N., „at Subjectiviteten maa være Princip ikke blot
for sig selv, men for al Objectivitet, og kun kan være hiint
formedelst den absolute Viden, dette formedelst den absolute
Magt og begge Dele som en Modsætningseenhed af Viden
og Magt: da vil det allerede heraf være klart, at den logiske
Bestemmelse af Grundforholdet mellem Gud og Verdensideen
maa blive theistisk“. (Grundideernes Logik II. Fortale, S. IX).
Men denne Slutning er altfor overilet og kan ingenlunde uden
videre indrømmes“.

Hvorfor kan Slutningen ikke indrømmes? „Man maatte“,
siger Biskoppen, „jo være vildfremmed i Philosophiens Hi\-%
storie for ikke at vide, at ogsaa Phantheismen har sit Sub\-%
jectivitetsprincip, der snart er bleven bestemt som et Videns\-%
princip, snart som et Magtprincip og snart som en Eenhed
af begge“. En ret interessant Confusion.

Det Subjectivitetsprincip, der udvikles i Grundideernes
Logik, er under Udviklingen Skridt for Skridt, ja Punkt for
Punkt, adskilt fra og retfærdiggjort ved sin Modsætning til
de i Philosophiens Historie fremragende Subjectivitetsprin\-%
ciper. Kun en med den videnskabelige absolut ueensartet
Kritik kan her mene at udrette Noget ved en Paastand som
den (Leilighedsskr. S. 48), at jeg „sætter Hegels logiske Sub\-%
jectivitet sammen med et evigt Magtprincip“; kun en med
den videnskabelige absolut ueensartet Kritik kan mene, at en
tænkende Læser ikke ved sig selv skulde opdage en saa tyndt
tilsløret Modsigelse som denne: først paastaaes der, at Sub\-%
jectiviteten i min Logik er Hegels logiske Subjectivitet; saa
indrømmes der i samme Aandedræt, at jeg har et ganske
andet Begreb om den logiske Subjectivitet end Hegel. Man
skulde vel knap troe det, dersom man ikke kunde see det med
% --- p. 120 ---
\opage{120}\setcounter{footnote}{0}%
egne Øine. Men Confusionen er et Factum, ja „et
Grundfactum“.

Den confuse Paastand med den ligesaa confuse Ind\-%
rømmelse staaer at læse i Leilighedsskriftet S. 48--49, og
lyder i sin Sammenhæng saaledes: „Den virkelige Tanke\-%
gehalt i den Hegelske Lære om „Begrebet“, den evige νοῦς,
er den, som Heiberg i sin Fremstilling af Hegels Logik frem\-%
sætter paa følgende Maade: „Subjectiviteten udtrykkes ved
Jeg; men dette er endnu ikke Bevidsthed; det er kun den for
Naturens og Aandens Skikkelse fælles (altsaa logiske) Sub\-%
jectivitet og saaledes \emph{den Rod, hvoraf Bevidstheden
udvikler sig, skjøndt endnu ikke denne}“. (Pros.
Skrifter 1 B. S. 276--77). Dette er klar og forstaaelig
Tale. Men den er ikke theistisk og udgiver sig ikke for at være
det. Naar nu Prof. Nielsen sætter Hegels logiske Subjec\-%
tivitet sammen med et evigt Magtprincip, saa indsees ikke,
hvad der i theistisk Henseende, i Henseende til selve Subjec\-%
tivitetsprincipet er vundet. Hegel lader nu den uendelige
Subjectivitet, dette Vidensprincip, der er Bevidsthedens Rod,
men ikke den selv, komme til Bevidsthed i Mennesket ved den
bekjendte Sætning, at Menneskets Viden om Gud er Guds
Viden om sig selv. Denne Sætning forkaster Prof. Nielsen
og med al Føie, paa Grund af vor psychologiske Bevidstheds
uimodsigelige Indskrænkning, Dualismen mellem Viden og
Magt. Den absolute Viden maa da findes i den ontologiske
Subjectivitet selv, uafhængigt af den menneskelige Bevidsthed,
der kun participerer deri. Men naar vi da spørge hvor\-%
ledes dette skal tænkes, da faae vi vistnok en Mangfoldighed
af dialektiske Reflexionsbestemmelser om Vexelforholdet mellem
det Subjective og det Objective, men maae særligt holde os
den Formaning efterrettelige, at vi ikke tør maale den gud\-%
dommelige Viden efter vor egen „psychologiske Alen“. Men
hvad vil dette sige Andet, end at vi her ikke tør personificere
% --- p. 121 ---
\opage{121}\setcounter{footnote}{0}%
og kun have med en Idee, en \textit{res cogitans} at gjøre, der er
uden virkelig Selvbevidsthed?“

Disse ere Ordene. Vi bede om Aandens Bistand: at vi
maae fatte disse Ord! at vi ret maae belæres af „det kongelige
Princip“! at Fichte og Hegel, Heiberg og Kierkegaard maae
smile til os fra Logikens Himmel! at det maa forundes os
at nyde Skatten, den rige Skat af logiske, psychologiske,
ethiske, pathetiske Confusioner, hvormed Biskop Martensen
ved denne Leilighed saa gavmildt har regaleret Logiken!

\emph{Første Confusion}. Her siges om Heibergs Tale,
der ganske vist er klar og forstaaelig, at den hverkeu er theistisk% printer's defect: "hverkeu" (turned n) for "hverken", p. 121 l. 11; transcribed as printed
ei heller udgiver sig for at være det. Denne besynderlige
Yttring har sit Udspring fra en ligesaa besynderlig Sam\-%
menblanding af logisk og theologisk Theisme. I Ideens
Navn have naturligviis baade Hegel og Heiberg paa det
Kraftigste værget sig imod at faae deres logiske Theisme
confunderet med den vulgære i dogmatiske Lærebøger figu\-%
rerende Forestillingstheisme. De havde i denne Henseende
begge et lige skarpt Blik for Theologernes velbekjendte For\-%
vexling af Sag og Billede, Begreb og Forestilling. Men
dersom Nogen heraf lader sig forlede til den Antagelse, at
„den speculative Logik“ ikke skulde være anlagt paa at give
os en „logisk Theisme“, da tager han mærkelig feil. At
Biskop Martensen med al sin dogmatiske Speculation just
var synderlig stiv enten i platonisk eller hegelsk Logik, har
jeg rigtignok aldrig formodet. Hvad Plato angaaer, har Bi\-%
skoppen naturligviis vraget den Propædeutikens „magre Kost“,
hvormed en Dialog som „Parmenides“ opvarter; han har natur\-%
ligviis, for at forøge sin Samling af philosophiske Læsefrugter,
udsøgt sig „det Høieste og det Dybeste“, nemlig de Steder,
hvor Plato „taler mystisk, idet han ligesom visionært berører,
hvad ikke kan naaes ad den blot dianoetiske Vei“. Hvad
Hegel angaaer, da forekommer det mig, som om jeg engang
havde hørt, at „den store Logik“, hvorover Saamange i sin
% --- p. 122 ---
\opage{122}\setcounter{footnote}{0}%
Tid fik Hovedpine, blev erklæret for en Art ufrugtbar
Formalisme; at en dogmatisk Genius godt kunde nøies med
det lille Udtog i Encyclopædien og stole paa sit Blik for det
Positive. Men at Biskop Martensen i den Grad skulde
have udslettet ethvert Minde om J. L. Heibergs logisk-dialektiske
Virksomhed og i den Grad staae udenfor Forstaaelsen af det
Logiske i „den speculative Logik“, at han virkelig kan mene
at sige noget Træffende, naar han siger, at Heibergs Tale
om den logiske Subjectivitet hverken er eller vil være theistisk,
det har dog overrasket mig.

Enhver, der nogensinde har læst og forstaaet Heibergs
philosophiske Skrifter, kan vidne med mig, at Heiberg hverken
kan eller vil adskille sit logiske Gudsbegreb fra det christelige
Gudsbegreb, uagtet han baade kan og vil erkjende „den
evige Forskjel mellem Philosophie og Religion“. Det christe\-%
lige Gudsbegreb er theistisk; det theistiske Gudsbegreb er en
Concretion af det logiske Subjectivitetsbegreb. „Idet Hegel“,
siger Heiberg (Perseus Nr. 1. S. 39--40), „har reduceret
den philosophiske Forudsætning til en saa abstract Bestemmelse
af Gud, at den i systematisk Henseende ingen Forudsætning
er, og altsaa heller ingen Hindring for Systemets absolute Be\-%
gyndelse, saa kan den ikke være tilfældigen optagen af den indre
Erfaring, saaledes som Tanken eller Substansen eller Monaden.
Naar man nu veed, at denne Philosophie ender med Er\-%
kjendelsen af Gud, at den fremstiller al Udvikling som et
Kredsløb, hvori Endepunktet falder sammen med Begyndelses\-%
punktet .... saa vil man indsee, at Systemet ikke kunde be\-%
gynde med Guds \emph{abstracteste} Bestemmelse, ifald det ikke
endte med den \emph{concreteste}“. Og saa skulde Heibergs klare
og forstaaelige Tale hverken være theistisk eller udgive sig for
at være det!

\emph{Anden Confusion}. Naar Theismen i Grundideernes
Logik betegnes som \emph{udpræget} Theisme, da skeer det for at
adskille den fra en Theisme, der ikke er udpræget, men \emph{dunkel}
% --- p. 123 ---
\opage{123}\setcounter{footnote}{0}%
og tvetydig, nemlig den hegelske. Forskjellen imellem dunkel
og udpræget Theisme maa fra Grunden af søges i en for\-%
skjellig Bestemmelse af Subjectivitetens ontologiske Væsen. I
denne Henseende blander nu Biskop Martensen Hegelsk, Ikke-%
Hegelsk, Logisk, Ethisk, Selvhed, Person og Personification paa
en besynderlig Maade imellem hverandre. Det siges reentud,
at jeg sætter „Hegels logiske Subjectivitet sammen med et
evigt Magtprincip“; Subjectiviteten i min Logik skal altsaa
ikke være forskjellig fra Subjectiviteten i den hegelske. Men
saa siges der igjen, at der er en Evigheds Forskjel imellem
mit Subjectivitetsprincip og det hegelske: hos Hegel kan nemlig
dette Princip først komme til Bevidsthed i Mennesket, hos mig
maa den absolute Viden findes i den ontologiske Subjectivitet,
uafhængig af den menneskelige Bevidsthed, der kun participerer
deri. Og saa siges der igjen, at denne Forskjel ikke er nogen
Forskjel; og hvorfor ikke? Fordi vi ikke tør maale den gud\-%
dommelige Viden efter vor egen psychologiske Alen. Naar
tør vi da, ifølge Biskoppens Anviisning, maale den gud\-%
dommelige Viden efter vor egen psychologiske Alen? Naar
vi paa en passende Maade blande det Ethiske i det Logiske,
indføre „Kjærligheden“ og „den gode Villie“! Istedenfor det
vidende Selv behøve vi blot at sætte det villende Selv, og
den ontologiske Subjectivitet kommer strax til Bevidsthed; thi i
Villien er jo Kjærligheden. Og Kjærligheden? Ja Kjærlig\-%
heden! Kjærligheden i Logiken og Kjærligheden i Ethiken og
Kjærligheden i Dogmatiken, Kjærligheden i de tre Flammer,
den guddommelige Kjærlighed, ja det er en anden Sag: den
guddommelige Kjærlighed tør vi naturligviis nok maale efter
vor egen psychologiske Alen!

Altsaa: hvoraf vide vi, at den i Dogmatiken \emph{villende}
Gud ikke er en Personification, en \textit{res volens}, men et selv\-%
bevidst personligt Selv? Vi \emph{maae} vel lade Biskop Martensen
svare: deraf, at den guddommelige Villie er den menneske\-%
lige saa fuldkommen lig, at den kan maales efter vor egen
% --- p. 124 ---
\opage{124}\setcounter{footnote}{0}%
psychologiske Alen. Og hvoraf vide vi, at en i Logiken
\emph{vidende} Gud er en Personification, en \textit{res cogitans}, ikke
et selvbevidst personligt Selv? Vi \emph{maae} vel lade Biskop
Martensen svare: deraf, at den guddommelige Viden er
den menneskelige saa fuldkommen ulig, at vi ikke tør maale
den efter vor egen psychologiske Alen. Og hvoraf vide vi
saa, at den guddommelige Villie i Dogmatiken og den gud\-%
dommelige Viden i Logiken absolut maae være hinanden saa
ulige, at Villien kan maales, men Viden ikke maales efter
vor egen psychologiske Alen? Hvoraf vide vi, at der vel kan
tænkes en Personification af en \textit{res cogitans}, men ikke af
en \textit{res volens}? Vi \emph{maae} vel lade Biskop Martensen svare:
deraf, at her er tilveiebragt en Confusion, en grundig Con\-%
fusion i Ideeriget!

\emph{Tredie Confusion}. Der er i vor psychologiske Be\-%
vidsthed to ueensartede Elementer, et menneskeligt og et gud\-%
dommeligt, et forkrænkeligt og et evigt. Blandes nu de to
Elementer saaledes sammen, at man snart optager begge,
hvor man dog kun skulde optage det ene, snart igjen udskyder
begge, hvor man dog kun skulde udskyde det ene: da er For\-%
virringen i fuld Gang. Det er en saadan Forvirring, der
indtræder med hvad vi her kunne kalde Biskop Martensens
tredie Confusion. Ligesom han i første Confusion har sammen\-%
blandet Begrebstheisme og Forestillingstheisme, og med anden
Confusion løber vild i personificeret og ikke personificeret
Subjectivitet, saaledes kommer han i tredie Confusion til at
blande det Uforkrænkelige med det Forkrænkelige, det Gud\-%
dommelige med det Menneskelige i den psychologiske Bevidst\-%
hed. Tingen er, at Forfatteren af Leilighedsskriftet har været
saa gridsk paa at finde sit \emph{Viisdomsideal} i Logiken, at
han ikke har havt Leilighed til at assimilere sig det logiske
Vidensideal.

Om det Princip, hvorefter Grundideernes Logik anvender
den psychologiske Bevidsthed ved Udviklingen af Vidensidealet
% --- p. 125 ---
\opage{125}\setcounter{footnote}{0}%
og den vidende Subjectivitet, handles der udførligt i første
Deels første Afsnit. Vi udhæve Følgende (S. 230): „At
det i Logiken som i enhver anden Videnskab væsentlig kommer an
paa tilforladelige Udgangspunkter og paalidelige Kjendemærker,
naar ellers Undersøgelser skulle kunne anstilles, forstaaer sig af
sig selv. Det er den psychologiske Subjectivitet med dens For\-%
udsætninger og Betingelser, hvorfra vi i Grundideernes Logik
bestandig maae gaae ud, og hvortil vi igjen maae vende til\-%
bage. Det kommer ved de logiske Analyser Altsammen an
paa, hvorledes de som Kjendsgjerninger givne Udgangspunkter
benyttes: at blive i Endeligheden tilfredsstiller ikke; at abstra\-%
here fra alt Endeligt og lade Tankerne svæve i det tomme
Uendelige, tilfredsstiller heller ikke. Her er ved Behandlingen
af Forholdet mellem den fuldkomne Viden og Videns psycho\-%
logiske Factorer i Logiken Noget, der minder om Behand\-%
lingen af de saakaldte „ubestemte Former“ i Mathematiken;
ingen af de endelige Forudsætninger kan ligefrem beholdes
som Kjendsgjerning, men enhver Kjendsgjerning kan og maa
benyttes som Moment: Opgaven er som at regne med uende\-%
lige Størrelser. Umiddelbart og ligefrem kan Ingen sige,
hvormeget 0.\,$\infty$ gjælder\footnote{Jvfr. Mathematik og Dialektik. Kjøbhvn. 1859. S. 65--90.}, eller hvilken Værdi man maa
tillægge $\infty-\infty$; ligesaa kan heller Ingen umiddelbart og
ligefrem afgjøre, hvorledes et Realbegreb gaaer op i sin
Existens, og hvilke Bevægelser Tankehandlingerne maae frem\-%
kalde i den rene Idee. Ikke ved et Spring, men ved en til
Functionernes Eiendommelighed svarende Overgang maa den
begribende Subjectivitet i hvert Tilfælde bevæge sig fra det
Endelige til det Uendelige, dersom det ellers skal lykkes den at
finde Tilbageveien fra det Uendelige til det Endelige.“

Havde Biskop Martensen fra sin monarchiske Høide kunnet
nedlade sig til at overveie, hvad heri ligger, og hvad heraf
følger, da havde han sikkerlig forskaanet os for den Art Cruditeter,
% --- p. 126 ---
\opage{126}\setcounter{footnote}{0}%
hvormed han i Leilighedsskriftet saa salvelsesfuldt udbreder sig
over Forskjellen imellem personificeret og personlig Subjec\-%
tivitet. Det er ikke ved at sige „Kjærlighed“ og „det Ethiske“
og „tre Flammer“, at man i Videnskaben forvisser sig om
Subjectivitetens absolute Gyldighed; her fordres en alsidig,
grundig, logisk-dialektisk Analyse af Vexelbestemmelsen imellem
Subjectivitetens psychologiske og ontologiske Momenter. Men
ved en saadan Analyse vil man da ogsaa kunne forvisse sig
om, at en ontologisk Personification ligesaa lidt kan erstatte
os den \emph{vidende} som den \emph{villende} Guddom, at vi altsaa
--- fra Videns Synspunkt --- have samme, d.\,v.\,s. ligesaa
stor Vished for en guddommelig Selvbevidsthed i Logiken
som i Ethiken.

Men, indvender maaskee den tænkende Læser, i Logiken
er dog Alt saa abstract; „Fylden“, det rette Concrete kommer
først, naar vi komme til Dogmatiken og til Ethiken. Til en
Opklaring af herhen hørende Misforstaaelser vil Forfatteren
af Leilighedsskriftet maaskee bedst kunne veilede os, idet han
til Bedste for Philosophien befinder sig i en stadig Fremgang
og nu staaer i Begreb med at træde ind i sin:

\emph{Fjerde Confusion.} I et Leilighedsskrift som nær\-%
værende Martensenske er Alt naturligviis System --- „thi“,
siger Biskop Martensen S. 142, „et ἀσύστατον, et ikke i
sig Sammenbestaaende, hvis Mangfoldighed manglede Eenhed,
vilde være i Modstrid med sig selv“ ---; men naar Alt i et
Skrift er systematisk, og samme systematiske Skrift er temmelig
rigt paa Confusioner, saa maae Confusionerne jo ogsaa være
systematiske. Skulde Systematiken her udvikles fuldstændig,
maatte der naturligviis tages Hensyn ikke blot til Leddenes
indbyrdes Sammenhæng, men ogsaa til deres Antal. Hvad
Antallet angaaer, er det nærved at gaae os med Confu\-%
sionerne i Leilighedsskriftet, som det nok gaaer Astronomerne
med Stjernerne i Mælkeveien: jo længere vi fortsætte
vore Observationer, desto flere bliver der. Vi havde rigtignok
% --- p. 127 ---
\opage{127}\setcounter{footnote}{0}%
tænkt paa ved streng Oekonomie at reducere det hele
Antal speculativt til Tolv --- Tallet Tolv er i speculative
Systemer meget at anbefale, og det af forskjellige Grunde:
for det Første passer det ligesaa godt til en typisk Tredeling
som til en fremherskende Fiirdeling, og frembyder altsaa to
store Sider for „en Totalanskuelse“; for det Andet har
Tallet Tolv en fleersidig symbolsk Betydning: Tolv Himmel\-%
tegn, tolv Stammer i Israel, tolv Sognekirker i Han-Herred\footnote{Er der flere, maae de overflødige fjernes ved en Abstraction; er
der ikke saamange, maae de manglende tilføies i en Totalanskuelse;
thi Totalanskuelsen er det Vigtigste.}
o.\,s.\,v. For det Tredie .... --- Men det overgaaer vore
Kræfter; det bliver for vidtløftigt.

Altsaa til Sammenhængen! Det Principielle i Confu\-%
sionerne beroer paa en uophørlig Forvexling af Begreb med
umiddelbar Forestilling. Usikkerhed paa hvad der er Begreb,
og hvad der er Forestilling maa især være piinlig for den,
der vil hjælpe paa Philosophien. Af Usikkerhed paa Begreb
fulgte Forvexlingen af logisk og theologisk Theisme: det
var første Confusion. Af Usikkerhed paa Begreb fulgte
Usikkerhed paa personlig og personificeret Subjectivitet med
overvættes Tillid til dogmatisk „Kjærlighed i tre Flammer“:
det var anden Confusion. Af Usikkerhed paa Begreb fulgte
Usikkerhed paa Guddommeligt og Menneskeligt i den psycho\-%
logiske Subjectivitet, forbunden med en charakteristisk Afsmag
for Alt, hvad der hedder Reflexion og Reflexionsbestemmelse,
altsaa med en charakteristisk Afsmag for det Logiske i Logiken:
det var tredie Confusion. Af disse ingenlunde subordinerede,
men indbyrdes coordinerede Confusioner vil man nu let kunne
forklare sig den fjerde og sidste, vi her ville nævne: Confu\-%
sionen med det Abstracte og det Concrete.

Idet Grundideernes Logik lægger an paa at udvikle
Subjectiviteten som \emph{vidende}, vil den naturligviis ikke over\-%
% --- p. 128 ---
\opage{128}\setcounter{footnote}{0}%
see, at en vidende Subjectivitet væsentlig maa være en \emph{vil\-%
lende}. Men, det forstaaer sig, den kan som Logik kun
fremhæve Selvbestemmelsens logiske, ikke dens ethiske Mo\-%
menter, og maa for Alting stræbe at give os disse Momenter
i Begrebets, ikke i Billedets Form. „Al Subjectivitetsbestem\-%
melse er væsentlig Selvbestemmelse, men en Selvbestemmelse,
der forudsætter Andet og betinges af Andet. At den fuld\-%
stændige Subjectivitetsbestemmelse antager en tredobbelt Form,
har netop sin Grund i, at Bestemmelsen af Andet skal gaae
op i Bestemmelsen af Selvet. Dette er nu paa Relativi\-%
tetens og Endelighedens Trin en Umulighed. Bestemmelsen
af Andet staaer i psychologisk Forstand som en Skranke, en
Hindring for Selvbestemmelsen. Ligeoverfor Andet, den mægtige
Skranke, er Selvbestemmelsen afmægtig, abstract og formel.
Afmagten sees da ogsaa baade paa Tanken og Ordet, thi
den rene Selvtænken er en tom Reflecteren, og Ordet et
Navn, der Intet formaaer over Gjenstanden. ... Men naar
Talen er, ikke om den psychologiske, men om den ontologiske,
ikke om den relative, men om den absolute Subjectivitet,
maa Nominalismen forsvinde. Det er det almægtige Ords
absolute Idealitet, hvorpaa Eftertrykket falder.“ (Grundid.
L. II. S. 336--37).

Opgaven er en logisk Formudvikling; den logiske Op\-%
gave er: nøiagtigt og strengt at udvikle den Form, den
eneste mulige Form, hvori en absolut Subjectivitet lader sig
tænke. Da Subjectiviteten skal i ubetinget Eenhed svare
baade til det Almene, det Særskilte og det Enkelte, maa dens
logiske Væsenhed sættes i en Slutning, en Midte med to
Extremer, der i det Uendelige modsættes og dog gjennem\-%
trænge hinanden, altsaa i en trefoldig Form. Opgaven er at
drive Formudviklingen til størst mulig Fuldendelse, ikke holde
den i det Graalige men angive Momenterne saa bestemt, at
den kan underkastes en Kritik.

Hvad vil nu skee, naar en philosophisk Tænker, en vir\-%
% --- p. 129 ---
\opage{129}\setcounter{footnote}{0}%
kelig Logiker, engang faaer isinde at revidere en saadan
Fremstilling? Her vil ganske vist ikke blive klaget over, at
det kun er Former, Logiken byder, at den ikke tillige kan
give os „Livet“ og „Fylden“ og „Kjærligheden“ og „de tre
Flammer“; nei, her vil blive spurgt, om Formen er correct,
om Formudviklingen er alsidig, conseqvent, stringent. Befindes
det saa, at Formen er mangelfuld, vil der blive anbragt
Correctiver, vel at mærke, logiske, ikke theologiske Correctiver.

Men dersom det nu ikke er en Logiker, men tvertimod
en Theolog, der, i Anledning af hvad han forstaaer ved Theisme
og Subjectivitet, faaer isinde at taxere Formudviklingens theo\-%
logiske Udbytte: hvad saa? hvad Dom ville vi saa faae?
En streng Dom! Det Logiske, det blot Logiske har jo ingen
theologisk Fylde; hvad er saa dette Logiske? Det Logiske,
der ikke strutter af theologisk Fylde, er magert som „Pro\-%
pædeutikens magre Kost“; det er „det Kenologiske“. „Det
Hele“, siger Leilighedsskriftet (S. 75--76), „er den tommeste
Formalisme, hule Nødder og Skyggespil op og ned ad Væg\-%
gene, ved hvilket man istedetfor det Antilogiske kommer til at
tænke paa det \emph{Kenologiske}, fordi den ethiske Virkelighed
ligger udenfor Systemet.“ Præmisserne for denne Dom
bygges ingenlunde paa nogen logisk Vurdering af Formud\-%
viklingen; thi i saa Fald kunde her maaskee dog opnaaes
Noget ved en videnskabelig Forhandling. Nei den Kritik,
hvormed vi her have at gjøre, udspringer psychologisk af „den
gode Villie“, ontologisk af en Speculation, der sætter Bil\-%
leder i Begrebers Sted, og er absolut ueensartet med viden\-%
skabelig Kritik.

„At søge det Absolute“, hedder det (S. 74), „i en
Sphære, som er høiere end den logisk-physiske, ligger i ethvert
Tilfælde langt udenfor de Intensioner, der vise sig i den
her omhandlede „Grundideernes“ (vel at mærke: Grund\-%
ideernes) Logik. Det Universum, hvori man her skuer ind,
er uden ethiske Beboere, som en af hine ubeboede Kloder i
% --- p. 130 ---
\opage{130}\setcounter{footnote}{0}%
Verdensrummet, hvor der kun findes Naturexistenser og et
tilsvarende logisk og mathematisk System.“

Paa en saadan Miskjendelse af det Logisk-Concrete, en
saadan Confusion af Billeder og Begreber, er Logiken
imidlertid saa vel forberedt, at Forfatteren af Leiligheds\-%
skriftet, blot ved at vende et Blad, kunde have fundet det
Svar, hvortil hans ethisk pathetiske Indsigelser saa høiligen
trænge. „Naar Særskilthedens Princip“ --- Talen er just
om et Blik i det Absolute --- (Grundid. Lgk. II S. 345)
„bringer det Almene i Zittren, naar Tanken lyser i Ordet,
og Ordet med Tankens Lys i Relativitetens Farveskjær ad\-%
spreder Mørket og aabner Erkjendelsen Indgang: hvilke
Syner fremstille sig da for Anskuelsen, hvilke Skikkelser
vil Tankens Blik da møde i Ideernes Rige? Det skulde
dog vel aldrig være den speculative Dogmatiks „legende Mulig\-%
heder“, „plastiske Forbilleder“, „magiske Virkeligheder“, en
„himmelsk Hærskare“ af de „Syner“, der ikke kunne nøies med
„en Aabenbaring \textit{ad intra}“, men ere i paatrængende Spænding
for at slippe ud og „løslades til Aabenbaring \textit{ad extra}“?
Ingenlunde!... „Magiske Syner“, speculative Hallucinationer
egne sig ikke for en guddommelig Subjectivitet, og maadelig nok
for en menneskelig. Med en Hærskare af slige Syner kan
Individualiteten maaskee blive en gjærende Mystiker, men
ikke en Digter, endnu mindre en Tænker. De Syner, der
oprinde for Digteren, de Forbilleder, der foresvæve Tænkeren,
have som Tilstrømninger fra Ideernes Kildevæld altid, om
Strømmen end gaaer nok saa stærk, det Særkjende, at de komme
med Ordet, ordne sig i Tanker og forklares i Domme; selv
i den mythiske Bevidsthed maae Synerne trods al deres
Magie underordne sig Ordets Magt. Men Syner, der
underordne sig Ordets Magt, have ikke den „creaturlige“
Nærgaaenhed, den apokryphiske Paatrængenhed, at de umid\-%
delbart skulde „begjære at løslades til en Aabenbaring \textit{ad
extra}“; nei, tvertimod! de klare sig just til Aabenbaring
% --- p. 131 ---
\opage{131}\setcounter{footnote}{0}%
\textit{ad intra}, thi Videns Ord i Ideernes Dyb er Aabenbaring
\textit{ad intra}.“

2) \emph{Theismen i Grundideernes Logik er kun en
Dualisme.}% numeral and parenthesis not letterspaced; the words are

Om denne Indsigelse kan jeg fatte mig kort. Det
Spørgsmaal, med hvad Ret Grundideernes Logik har sat
sin Modsætningseenhed af Viden og Magt istedenfor den
hegelske Identitet af Tænken og Væren, kunde vel fortjene
en Drøftelse; og det skal være mig en sand Ære at afhandle
dette Punkt, hvis Nogen med videnskabeligt Blik for Logik
og logiske Undersøgelser har Noget at fremføre. Men det
Martensenske Blik er ikke blot min Logik, men al Logik i den
Grad ugunstigt, at jeg nu næsten tør sige forud: Alt, hvad
Biskop Martensen --- fra sit hidtil fastholdte Standpunkt ---
for Fremtiden maatte have at meddele Publicum enten om
min eller Andres Logik, vil for mit Vedkommende blive
lagt \textit{ad acta}.

Kun som en Prøve paa det Eiendommelige i den Biskop\-%
martensenske Kritik skulle vi til Læserens Fornøielse anføre
et Exempel. „Vi komme“, siger Biskop Martensen (Leilig\-%
hedssk. S. 56), „til en ny Dualisme i Systemet og faae
Modsætningen mellem det Theoretiske og Practiske som en
Modsætning indenfor selve den theoretiske Philosophie, to
guddommelige Principer, den Vidende og den Mægtige, idet
der dog gives os den Forklaring, at „det er een og den samme
Magt, der i mægtig Selvmodsætning er ideal i Tænken og
Dommen, reel i Virken og Frembringen“. Hvad den% PRINTED: one closing “ only; the quotation opened at „til en ny Dualisme is never closed (quote balance +1)
biskoppelige Kritik her affærdiger som en forklarende Be\-%
mærkning er i Logiken en Grundtanke, der udvikles, bear\-%
beides og formes i lutter logiske Begrebsbestemmelser. Skulde
her virkelig være Noget at corrigere, da maatte en viden\-%
skabelig Kritik jo nødvendigviis gaae ind paa logiske Under\-%
søgelser; men Biskop Martensens Kritik og den videnskabelige
Kritik forholde sig ganske rigtigt som de tvende absolut ueens\-%
% --- p. 132 ---
\opage{132}\setcounter{footnote}{0}%
artede Træer, af hvilke det enes Frugter vare Mennesket
forbudne. Istedenfor at agte paa Tankebestemmelser, Re\-%
flexioner, Begrebsvendinger og slige Formaliteter, gaaer Bi\-%
skoppen som sædvanlig Totalanskuelsens brede speculative Vei,
og underholder saa sine Læsere ved følgende Forklaring.
„Vi kunne“, siger han, „maaskee tydeliggjøre os Forholdet
ved at forestille os den practiske Gud eller Naturprincipet
som den \emph{stærke Blinde}, der udfører alle Tankens Gjer\-%
ninger, og den theoretiske, den vidende Gud som den \emph{seende
Lamme}, der begriber, dømmer og slutter, anticiperer og
subjectivt bestemmer Alt, hvad den Blinde punktlig udfører,
men selv slet Intet kan gjøre, hverken røre Haand eller Fod.
Og Identiteten, Modsætningseenheden, eller „den absolute
Sandhed“, kunne vi maaskee forestille os saaledes, at det er
Naturprincipet eller den stærke Blinde, der paa sin Ryg
bærer den seende Lamme, som i sin „oversvævende Reflexion“
uden Tvivl maa holde sig til det Aristoteliske: at den rene
Theorie er det Saligste. Paa denne Maade faae vi dog en
vis Identitet. Thi i den Lamme kunne vi erkjende den
mægtige Viden, forsaavidt som den Blinde maa udføre Alt,
hvad denne paralytiske Krøbling tænker og indicerer; og i
den Blinde kunne vi erkjende den vidende Magt, forsaavidt
som Begrebet ikke blot „oversvæver ham“, men ogsaa er
ham instinctivt „iboende“ eller immanent.“

Dette er nu et Billede, og det Følgende er ogsaa et
Billede; ved „Kjærligheden“ og „de tre Flammer“, har ikke
Biskoppen sin hele Logik i Billeder! Biskop Martensen er
ved denne Leilighed endog saa nedladende, at han ligefrem
nedlader sig til, hvad Grundideernes Logik II (S. 194--201)
har charakteriseret som Opinionspolemik\footnote{At det imidlertid ikke gaaer slet saa let med den theologiske Opi\-%
nionspolemik i det nærværende Øieblik, som det gik for sytten Aar
% [printed p. 133 begins here, inside this note]
siden, vil Biskop Martensen selv kunne skjønne deraf, at allerede
R. Lindberg i sin Anmeldelse af „Leilighedsskriftet“ (Dansk Kirke\-%
tidende. 5te og 12te Mai) med logisk kategorisk Sikkerhed har
paapeget adskillige af de Confusioner, hvortil Forvrængningen
støtter sig.}.
% --- p. 133 ---
\opage{133}\setcounter{footnote}{0}%
Heraf sees nu, hvorledes ikke blot min Philosophies, men
al Philosophies Stilling maa blive, naar Biskop Martensen
en af Dagene faaer sit homerisk-aristotelisk-dogmatiske Uni\-%
versalmonarchie oprettet.

3) \emph{Theismen i Grundideernes Logik er en
slet og ret Naturalisme.}% numeral not spaced; "en" and "slet og" are set tighter than the rest (justification), taken as part of the spaced proposition

Det er en Indsigelse, der i Leilighedsskriftet atter og
atter vender tilbage, at Anlæget i Grundideernes Logik
umulig kan være theistisk, efterdi Theismens Princip er et
overnaturligt Princip, og Logiken for stedse har udelukket det
Overnaturlige fra Videnskaben. Hertil maa bemærkes, at
det gaaer med Begrebet: det Overnaturlige, som med saa\-%
mange andre: det kommer an paa at bestemme dets op\-%
rindelige Charakteer og den Region, hvortil det hører.

At Erkjendelsen af en høiere Natur, en høiere Tingenes
Orden end denne sandselige nærværende Verden, ikke skal
udelukkes af Videnskaben, derom kan Enhver, der blot vil
kaste et Blik paa Grundideernes Logik I (f. Ex. S. 452--456)
og II (f. Ex. S. 450--465) let overbevise sig. Der\-%
imod vil det til Forstaaelse, af hvad Logikens „udprægede
Theisme“ væsentlig betyder, og hvortil den sigter hen, være
nødvendigt at bestemme Forholdet imellem Logik og Reli\-%
gionsphilosophie noget nærmere.

Et af den nyere Philosophies store Grundspørgsmaal er,
i mine Tanker, Spørgsmaalet om Videns Grændse. At den
menneskelige Viden har sine Grændser, er der naturligviis
Ingen, der negter: det er kun Forholdet imellem Begrænds\-%
ningens absolute og relative Side, hvorom her tales. Fra
% --- p. 134 ---
\opage{134}\setcounter{footnote}{0}%
Fornuftkritikens Standpunkt satte Kant en absolut Videns\-%
grændse, idet han opfattede de aprioriske Erkjendelsesformer
saa udelukkende subjectivt, at de slet ikke lode sig anvende
paa det absolut Objective (Tingen i sig). Den Modsigelse,
hvori den kantiske Philosophie ved sin Dualisme hildede sig,
er fra logisk Side belyst i Grundideernes Logik (I. S. 76--85).
Da Schelling siden gik ud over Fichte ved at anskueliggjøre
de subjective Kategoriers objective Gyldighed, og Hegel i sin
„Logik som Videnskab“ udviklede Eenheden af det Subjec\-%
tive og Objective til Ideens immanente Begreb, vandt den
philosophiske Viden sit Absolute tilbage. Men i speculativ
Glæde over det Absolute og absolut Absolute glemte man at
gjøre en Adskillelse imellem den menneskelige og den gud\-%
dommelige Viden af det Absolute: Menneskets Viden af
Gud blev paa ukritisk-mystisk Maade til Guds Viden af sig
selv; den væsentlige Grændse blev udslettet. I Grundideernes
Logik er Grændsen sat paany: ikke i kantisk Forstand, saaledes
at Kategoriernes objective Gyldighed negtes, men i ontologisk
Forstand, saaledes at Forskjellen imellem den menneskelige og
den guddommelige Viden kan sees i vor Videns Ideal.

Hvor afgjørende det logiske Grændsespørgsmaal er for
Gudserkjendelsen og derved for Religionsphilosophien, er nu
let at indsee. Tages den absolute Grændse bort, saa for\-%
svinder Mysteriet, og Guds Væsen bliver gjennemsigtigt i Ideen.
At tage sit Standpunkt i den absolute Idee, er da at tage
sit videnskabelige Standpunkt i det guddommelige Centrum.
Fra et theocentrisk Standpunkt bliver saa Troens Gud ganske
rigtig til en Videns Gud; men Troen forsvinder, Religionen
bliver Videnskab. Sættes Grændsen derimod som absolut,
saa sættes med det Samme et absolut Mysterium; Guds
Væsen er da ikke længere et for Menneskets Blik gjennem\-%
sigtigt Væsen: et theocentrisk Standpunkt er umuligt.

Saasandt der er en absolut Vidensgrændse, et absolut
Mysterium, saasandt maa der være et med Vidensprincipet
% --- p. 135 ---
\opage{135}\setcounter{footnote}{0}%
absolut ueensartet Troesprincip; saasandt Viden og Tro
udgjøre to absolut ueensartede Regioner i vort høiere Be\-%
vidsthedsliv, saasandt vil ogsaa den ene og samme Gud vise
sig for os med en anden Skikkelse i Videns end i Troens
Lys. Det er fra dette Synspunkt, at jeg har gjort opmærk\-%
som paa en paulinsk Adskillelse imellem Videns Gud (Rom.
1, 19--20) og Troens Gud. At en saadan Adskillelse har
Gyldighed, kan da ogsaa sees, naar vi betragte den til\-%
svarende Adskillelse mellem det philosophiske og det evangeliske
Viisdomsideal. Naar Paulus (1 Cor. 1, 22) siger om
Grækerne, at de søge Viisdom, da overseer han ganske vist
den relative Forskjel i selve den græske Viisdom, Forskjellen
f. Ex. imellem en aristippisk og en platonisk, imellem en
epikuræisk og en stoisk Viisdom. Og naar han udtrykkelig
tilføier, at den Viisdom, han forkynder, idet han prædiker
den Korsfæstede, er Grækerne, d.\,v.\,s. de græske Viisdoms\-%
elskere (οἱ φιλόσοφοι), en Daarlighed, da indlader han sig
aldeles ikke paa at bestemme den philosophiske Forskjel imellem
disse Viisdomselskere. Den Viisdom, Paulus forkynder, er
og bliver en Daarlighed, ligesaavel for en Sokrates og en
Plato, som for en Aristipp og en Epikur; med andre Ord:
det evangeliske Viisdomsideal er absolut ueensartet med det
philosophiske, Troesidealet absolut ueensartet med Videns\-%
idealet.

Kan nu Gud som Gud erkjendes baade i Viden og i
Troen, men saaledes, at Gudserkjendelsen i Viden maa bygge
paa rationel Nødvendighed, Gudserkjendelsen i Troen der\-%
imod paa ubetinget Frihed, saa viser Forholdet imellem Logiken
og Religionsphilosophien sig for os i et nyt Lys. Aaben\-%
barer Guds Væsen sig i to Skikkelser, Nødvendighedens og
Frihedens, Videns og Troens: da maa der jo kunne gives
to Slags Theisme, en videnskabelig og en religiøs, en Videns
og en Troens Theisme. At de to Slags Theisme maae
% --- p. 136 ---
\opage{136}\setcounter{footnote}{0}%
forholde sig indbyrdes som de Erkjendelsessphærer, hvortil de
høre, og Videnstheismen i Principet være absolut ueens\-%
artet med Troestheismen, er en Selvfølge. Men den absolute
Ueensartethed udelukker ingenlunde Muligheden af, at Theismen
i den ene Skikkelse kan have Charakteermærker, der nøiagtig
svare til Theismen i den anden. Tvertimod! det ligger i
Sagens Natur, at Charakteermærkerne i Videnstheismen maae
speile Charakteermærkerne i Troestheismen som Nødven\-%
dighedsbestemmelser Frihedsbestemmelser, som Vidensbegreber% sic, elliptical as printed (no "og"/verb supplied)
Troesbegreber.

Under denne Forudsætning vil man forstaae: hvorledes
en christelig Verdensbetragtning vil, ikke umiddelbart, men
middelbart, kunne faae Indflydelse paa Philosophien som
autonomisk Videnskab; hvorledes Henblik til Troestheismen
kan være veiledende ved en streng videnskabelig Udvikling af
Videnstheismen; hvorledes den i Grundideernes Logik „ud\-%
prægede“ Theisme kan --- uden i noget Punkt at overskride
Grændsen for det Logisk-Physiske --- være bestemt ved og
bestemmende for en tilsvarende Religionsphilosophie.

Men hvad er da Religionsphilosophie? til hvilken Sphære
maa den henregnes? Forsaavidt Religionsphilosophien er en
Videnskab, maa den naturligviis høre under Videns Kreds;
men forsaavidt den har Religionen til Indhold, maa den
orientere os i Troens Sphære. Ved denne dobbeltsidige
Opgave bestemmes dens dobbeltsidige Charakteer som \emph{Grændse\-%
videnskab}. Den philosopherende Bevidsthed er, hvad der
da forstaaer sig af sig selv, det subjective Medium, hvori
de ueensartede Begreber mødes; heraf følger, at Bevidstheds\-%
eenheden maa forudsætte en formel, til Grændsebestemmelsen sva\-%
rende Begrebseenhed, at, med andre Ord, de absolut ueensartede
Begreber maae i Grændsen selv blive formelt eensartede. Hvad
en saadan formel Eensartethed har at betyde, kan let paavises.
Vi tale om absolut ueensartede Principer, men Princip er
% --- p. 137 ---
\opage{137}\setcounter{footnote}{0}%
jo \textit{in abstracto} altid Princip. Vi tale om ueensartede Er\-%
kjendelser, ueensartede Begreber o.\,s.\,v.; men idet vi abstra\-%
here fra det Eiendommelige, bliver Erkjendelse = Erkjendelse,
Begreb = Begreb.

Som Grændsevidenskab faaer Religionsphilosophien sine
Troesbegreber fra Troeslivet og sine Vidensbegreber fra
Videnskaben. I sidste Henseende viser det sig, hvad en logisk
Theisme har at betyde; i første Henseende kommer det an
paa religiøs Psychologie. Det nytter ikke, at Religions\-%
philosophen vil hjælpe paa Videnskaben ved Forsikkringer om
sin egen Tro, sine egne individuelle Erfaringer. Her spørges
ikke, hvorfra Philosophen har sine Erfaringer; her spørges om
Erfaringernes psychologiske Værdi i deres Forhold til Troes\-%
principet. Den tilsyneladende Modsigelse, at Religions\-%
philosophien paa de opstillede Vilkaar maa blive en Videnskab,
der handler baade om det, der falder \emph{udenfor}, og det, der
falder \emph{indenfor} Videnskaben, er let at hæve. Vil Nogen
paastaae, at der umulig kan gives en Viden om en Troes\-%
erkjendelse, naar Troeserkjendelsen efter sin Natur skal falde
udenfor Viden, lad ham da begynde med Begyndelsen og paa\-%
staae Umuligheden af, at der kan gives en Bevidsthed om
Jorden, naar den virkelige Jord efter sin Natur skal falde
udenfor Bevidstheden. En Indvending som den, at Troes\-%
begreberne ved at reflectere sig i Viden absolut maae blive
Vidensbegreber, har Intet at betyde; Videnskabens Opgave
er netop at vide det Ueensartede som ueensartet.

Som Grændsevidenskab maa Religionsphilosophien vise
os, hvorledes eet og samme Begreb kan blive ueensartet med
sig selv, eftersom det udvikles til Videns- eller til Troes\-%
begreb. Til Troesbegrebet Skabelse svarer f. Ex. et Videns\-%
begreb Skabelse, hiint udviklet af Nødvendigheden, dette
af Friheden; til Troesbegrebet Gudmenneske svarer Videns\-%
begrebet Gudmenneske o.\,s.\,v. Hvor forskjelligartede disse
% --- p. 138 ---
\opage{138}\setcounter{footnote}{0}%
Begreber ere, kan sees hos Strauß, Feuerbach og Kierkegaard.
Der skal nu paavises en alsidig articuleret Vexelbestemmelse
imellem de ueensartede Begreber; det skeer saaledes, at
Vidensbegreberne tjene som intellectuelt Underlag, Troes\-%
begreberne derimod træde i Forgrunden og udvikles systematisk.
Af Vexelforholdet mellem Begreberne fra de to Sphærer
fremlyser saa Vexelforholdet imellem den videnskabelige og
den religiøse Theisme, den Theisme, hvis Gudsbegreb kan
udvikles med logisk Nødvendighed i en Grundideernes Logik,
og den Theisme, hvis aabenbarede Gud er Gjenstand for Tro.

Men hvorledes er det muligt, at der kan tilveiebringes
Overeensstemmelse imellem de absolut ueensartede Begreber?
Fra hvilken Side vi betragte dem, synes de jo at være i
uforsonlig Strid. Stille vi os med Strauß og Feuerbach
paa Videns Standpunkt, løse Troesbegreberne sig op i
lutter Modsigelser; fastholde vi med Kierkegaard Troens
Standpunkt, medens vi tillige indrømme de af Videnskaben
paaviste Modsigelser, saa komme vi til at hævde Troen ved at
trodse Modsigelserne og tage til Takke med Paradoxet. Det
seer noget mørkt ud. Hvad om vi forsøgte det med en
mindelig Overeenskomst imellem de Stridende, et Slags
Concordat? Biskop Martensen, der har saamange Concordater
i Hovedet, at han knap kan rumme dem, har jo af sin
„gode Villie“ tilbudt at tjene os med et til midlertidig Af\-%
benyttelse. Ja, vilde vi være Theologer, saa maatte vi takke
for Tilbudet; men vi ville ikke være Theologer, nei ikke for
tusinde Concordater! Theologie og Concordat ere uadskillelige:
Biskop Martensen kan som Theolog umulig blive Con\-%
cordaterne qvit; han har dem ikke blot udenfor sig, han har
dem „tillige indeni sig selv“: han har dem i Dogmatiken.
Indvender man os, at der ifølge vor egen Indrømmelse ikke er
noget Concordat i Biskop Martensens „christelige Dogmatik“,
og at Biskop Martensens Dogmatik dog visselig maa være
Theologie, siden vi i den have villet angribe al Theologie,
% --- p. 139 ---
\opage{139}\setcounter{footnote}{0}%
saa svare vi: der er ikke blot et Concordat, men mangeslags
Concordater i „den christelige Dogmatik“; dog ikke φανερῶς,
men κατὰ κρύψιν, thi de ere alle medierede, medierede til
Vælling og skjulte i Vællingen.

Det absolute Mysterium gjør ethvert Concordat mellem
Tro og Viden ikke blot overflødigt, men umuligt. Den
Kjendsgjerning, at Troesbegreberne i deres absolut afgjørende
Modsætning til Vidensbegreberne ere beheftede med Mod\-%
sigelser, faaer nemlig i Conseqvens af det absolute Mysterium
en ganske anden Betydning i Religionsphilosophien end i
Theologien. Jeg har i mit lille Skrift: „Om Holbergs
Kirkehistorie og Theologie“ omhyggeligt søgt at paavise de
Modsigelser, hvori den orthodoxe Theologie hildede sig, naar
den ved sine „ydmyge Inconseqvenser“, tog, som det hedder,
Fornuften fangen under Troens Lydighed. Modsigelsen er,
at Theologien ikke blot anseer sig for en Videnskab, men
endog, som Biskop Martensen udtrykker det, for „den cen\-%
trale Videnskab“; thi at den centrale Videnskab, Videnskaben
med den centrale, den theocentrisk-centrale Fornuft, kommer i
Splid med Fornuften og vil tage Fornuften fangen, er unegtelig
et Tegn paa videnskabelig Fortvivlelse. Den orthodoxe Theo\-%
logie har et skarpt Blik for Troesbegrebernes Charakteristik;
men da den ikke er opmærksom paa Forskjellen imellem en
videnskabelig og en religiøs Motivering, saa kommer Moti\-%
veringen idelig i Strid med Charakteristiken.

Som Grændsevidenskab har Religionsphilosophien en
ganske anden Forklaring af de for Troesbegreberne charak\-%
teristiske Modsigelser end Theologien. At Troesbegreberne,
naar de udvikles af deres eget Princip, indeholde Modsigelser,
kan ikke, som fra et scholastisk-orthodox Standpunkt, forklares
deraf, at den syndige Fornufts oprindelige og nødvendige
Kategorier skulde staae i Strid med den aabenbarede Sand\-%
hed, heller ikke, som fra et kritisk-videnskabeligt Standpunkt,
% --- p. 140 ---
\opage{140}\setcounter{footnote}{0}%
deraf, at Troesbegreberne selv skulde være usande; nei, Mod\-%
sigelserne forklares, i Conseqvens af Principet, udtrykkeligt
deraf, at de manglende Mellembestemmelser ligge skjulte i det
absolute Mysterium. Da her altsaa ligesaa lidt kan være
Tale om at presse Troesbegreberne ind under en videnskabelig
som Vidensbegreberne ind under en religiøs Maalestok: saa
vil det heraf skjønnes, at der i Religionsphilosophien end ikke
vil kunne være Tanke om Concordat mellem Tro og Viden.
Den, der nu desuagtet har saa klare Øine, at han i noget af mine
Skrifter virkelig kan see, hvad en Biskop Martensen kan see,
enten „et Concordat“, eller Noget, der ligner „et Concordat“,
han kan see mere end jeg kan see; han kan see mere end
Schellings nye Klæder: han kan see „Keiserens nye Klæder“.

Her er, som sagt, et afgjørende Dilemma: \emph{enten maa
Theologien bort, eller den maa indtage sin Plads
som høieste Videnskab.} Biskop Martensen holder paa
det Sidste, jeg paa det Første.

Har nu Biskop Martensen Ret, da er han nok den
Mand, der vil vide at gjøre sin Ret gjældende, sætte sit
homerisk-aristoteliske: Een være Hersker! i Kraft og op\-%
rette det theologiske Universalmonarchie. Hvis dette skeer, da
ville mine Indsigelser ganske rigtig, hvad Leilighedsskriftet
(S. 142) jo ogsaa propheterer, blive at betragte som „en
literær Skandale (\textit{parturiunt montes})“. Men da ville
ogsaa Festlighederne i Monarchiet blive glimrende: paa den
palmestrøede Vei under Klokkernes Ringning vil „Dogma\-%
tiken med Christusbilledet“ bevæge sig i Procession gjennem
rigt smykkede Gader høitidelig opad mod den objective al\-%
mindelige Kirke med det objective almindelige Taarn.

Men skulde det vise sig, at Biskop Martensen ikke har
Ret; skulde det vise sig, at han med sit Begreb om Real\-%
videnskaberne ikke formaaer at indføre „en grundig Subordi\-%
nation i Ideeriget,“ altsaa --- trods sin „\emph{gode Villie}“ ---
% --- p. 141 ---
\opage{141}\setcounter{footnote}{0}%
heller ikke formaaer at oprette Monarchiet: da var det jo
muligt, at Biskoppen i Betragtning af de mange og store
Hindringer, der optaarne sig mod hans: \emph{Een være Her\-%
sker!} kunde komme paa de Tanker, at Touren nu var til
ham at skifte Standpunkt. Skulde dette blive Stridens Ud\-%
fald, og skulde der kunne lægges nogen Vægt paa mine Be\-%
stræbelser: da tilbyder jeg med velvillig Hu mine \textit{bona
officia}, og vilde i Sandhed ansee det for en stor Ære, om
jeg engang endnu kunde komme til at arbeide sammen med
en saa begavet Mand som Biskop Martensen.

% The book ends here: a short centred rule with a lozenge swelling at its
% middle. No signature, no date, no "Ende".
\begin{center}
\rule[0.5ex]{3em}{0.4pt}\,$\diamond$\,\rule[0.5ex]{3em}{0.4pt}
\end{center}

% --- Rettelser (unnumbered page, PDF 149) ---
% Heading in Fraktur with a short rule below it; a short closing rule at the
% foot of the list. "Sanctus"/"Sancti" are in antiqua. The ditto dashes are
% printed. The first entry's closing mark is unpaired by design (it is
% the missing mark being supplied): quote balance -1 on this page.
\section*{Rettelser.}
\noindent
Side 17 Linie 4 fra oven legitim; læs: legitim“;\\
--- 43 Linie 17 fra neden saadan læs: saadant\\
--- 46 Anm. Linie 7 fra neden \textit{Sanctus} læs: \textit{Sancti}\\
--- 61 Linie 10 fra oven hedder det om læs: hedder det, som oven\-%
for anført, om\\
--- 61 Anm. Lin. 1 fra neden S. 3 læs: S. 7.
% (this last fragment ends with the unnumbered "Rettelser." page,
%  PDF 149, set as \section*{Rettelser.} with NO page marker)

\end{document}
